% Compile with XeLaTeX (or LuaLaTeX), twice.
\documentclass[11pt,a4paper]{article}
\usepackage[margin=26mm,headheight=14pt]{geometry}
\usepackage{fontspec}
\setmainfont{Latin Modern Roman}
\setsansfont{Latin Modern Sans}
\setmonofont[Scale=0.90]{Latin Modern Mono}
\usepackage{amsmath,amssymb,amsthm,mathtools}
\usepackage{unicode-math}
\setmathfont{Latin Modern Math}
\usepackage{microtype,booktabs,longtable,array,enumitem,xcolor,fancyhdr}
\usepackage{xurl}
\usepackage[breaklinks,colorlinks,linkcolor=blue!45!black,urlcolor=blue!45!black,citecolor=blue!45!black]{hyperref}
\hypersetup{pdftitle={Audit of the claimed proof of RH and Dirichlet GRH},pdfauthor={Independent mathematical audit}}
\setlength{\parindent}{0pt}
\setlength{\parskip}{5.5pt}
\setlength{\emergencystretch}{2em}
\setlist{topsep=3pt,itemsep=3pt,parsep=0pt}
\pagestyle{fancy}
\fancyhf{}
\fancyhead[L]{\small Audit of RH and Dirichlet GRH}
\fancyhead[R]{\small 21 November 2026}
\fancyfoot[C]{\thepage}
\newtheorem{proposition}{Proposition}[section]
\newtheorem{lemma}[proposition]{Lemma}
\newtheorem{example}[proposition]{Example}
\theoremstyle{remark}
\newtheorem{remark}[proposition]{Remark}
\newcommand{\Z}{\mathbb Z}
\newcommand{\Q}{\mathbb Q}
\newcommand{\R}{\mathbb R}
\newcommand{\C}{\mathbb C}
\newcommand{\OO}{\mathcal O}
\DeclareMathOperator{\Sha}{Sha}
\DeclareMathOperator{\rank}{rank}
\DeclareMathOperator{\ord}{ord}
\DeclareMathOperator{\im}{im}
\DeclareMathOperator{\length}{length}
\DeclareMathOperator{\Rea}{Re}
\newcommand{\src}[1]{\cite{#1}}
\newcommand{\code}[1]{\texttt{#1}}
\begin{document}
\begin{center}
{\LARGE\bfseries Audit of the claimed proof of the Riemann hypothesis and Dirichlet GRH\par}
\vspace{6pt}
{\large Findings, counterexample searches, and limits of verification\par}
\vspace{7pt}
21 November 2026
\end{center}
\begin{abstract}
The announcement at \url{https://alpo.ge/remark.pdf} and all fourteen
linked archive parts were downloaded, identified by cryptographic
hashes, and decompressed. The archive contains 175\,187 Lean source
files. This continued audit examines the intended conclusions,
consequential mathematical interfaces, counterexamples, and the
integrity of the complete exported proof.

\textbf{No fatal mathematical flaw in the submitted proof has been
established.} The entire 61.57-gigabyte proof export has been
inventoried locally and contains only the standard axioms
\texttt{propext}, \texttt{Classical.choice}, and \texttt{Quot.sound}.
An independent structural comparison of all five target statements
and their transitive definitions with the canonical challenge has
passed. Separately, Andrew Sutherland's public verification records
report complete compilation, replay of every module, and acceptance
of the full export by the independent Rust kernel \emph{nanoda}.
We did not execute a complete kernel replay in this workspace.

The mathematical sections give independent proofs of the analytic
criterion, numerical comparison, integral-lattice facts, restricted
cohomological vanishing, and other tested ingredients. Explicit
counterexamples show why several tempting stronger assertions fail;
they do not satisfy the actual hypotheses of the submitted results.
The nine additional quasi-Riemann and Landau--Siegel references are
reviewed, and their useful consequences and limitations are recorded.
This version adds full ordinary proofs of the conditional finite-family
quadratic selection, the cyclotomic interval mechanism, the quotient
divisor formula, integral coefficient duality, the arithmetic complex
and homology bound, and the full local theta majorant induction,
as well as the unit/class-group and cyclotomic Selmer comparisons,
the actual coefficient-map bounds, and a stronger global ideal sum.
An independent algebraic $K$-theory argument now proves the coarse
cohomological estimate from explicitly cited classical theorems,
and the original finite-field homology construction has also been
completed in ordinary mathematics. Further additions prove the actual
adelic class coordinates and extra cosets, the compact and integral
branching laws with their exact scalars, and the integral lowering maps
for every auxiliary character. The archimedean coefficient identities,
separated interpolation sets, terminal radial positivity, and its
finite orthogonal decomposition are reconstructed. Explicit
counterexamples test invalid simplifications; the remaining sharp
arithmetic divisibility and noncompact integral inputs are identified.
The hook product proves the required local coefficient-lattice equality,
and the primitive unitary densities and exact mass ratio are checked,
using the explicitly stated classical global mass theorem.
The formal degree is now identified with its exact Haar normalization
and smallest-compact-type condition, and the column dimensions are
proved. The ordinary CM theta laws, both matrix supports, the
finite coefficient field, and the weighted Hecke commutation law
give the actual selected-prime column bound. The deep-cusp estimate
completes the ambient multiplicity comparison. Three complementary
Selmer cases have direct conventional proofs. All four polynomial
differential operators and the regular spectral inverse are
constructed and their radial actions verified; the actual recursive
frame and its compactly supported positive form are proved.
The actual Jacquet transforms now have proved rapid Fourier decay and
pointwise smooth compact-input marginal densities. Degenerate off-diagonal
vanishing, stronger Gram-volume estimates and sharper dyadic summation
are reconstructed. The whole-field character comparison retains its
integral torsion error, and the adjoint-vanishing implication is matched
to Thorne's classical theorem with all its hypotheses checked.
The complete source development has not been reconstructed in ordinary
mathematical prose. The companion manuscript isolates the remaining
uniform arithmetic construction and does not claim a finished proof
of GRH.
\end{abstract}
\setcounter{tocdepth}{1}
\hypersetup{bookmarksdepth=2}
\begingroup\small\setlength{\parskip}{0pt}
\tableofcontents
\endgroup
\newpage

\section{Scope and evidentiary standard}
There are three different questions:
\begin{enumerate}
\item Does the announcement give a complete argument in ordinary mathematical prose?
\item Do the final formal statements express the intended Riemann hypotheses?
\item Are all the proofs of those statements valid, with only the intended foundational axioms?
\end{enumerate}
The announcement describes a method; it is not itself a complete conventional proof. The source inspection and the complete structural comparison support an affirmative answer to the second question. For the third, Section~\ref{sec:complete-export-verification} distinguishes the independent published kernel acceptance from the complete local export inventory and statement comparison. These formal checks are substantial evidence, but do not constitute an exhaustive human reconstruction of the mathematics.

A serious mathematical objection would require a false assertion with all its hypotheses satisfied, an inference that does not follow, or an essential hypothesis that is not supplied. A theorem whose proof has not been checked is not thereby a false theorem. Likewise, an error in a simplified paraphrase is not a counterexample to a more carefully qualified formal statement.

The counterexamples in Sections~\ref{sec:landau}, \ref{sec:integral}, and \ref{sec:tempered} are counterexamples to \emph{weakened or overgeneralized alternatives}. The inspected source excludes them. They are included to document the counterexample searches and to prevent their being misreported as flaws in the submitted proof.

\subsection{The exact source examined}
\begin{center}
\begin{tabular}{@{}lp{9cm}@{}}
\toprule
Item & Recorded value \\
\midrule
Announcement & \url{https://alpo.ge/remark.pdf}; 4 pages; 245,114 bytes. \\
Archive parts & \url{https://alpo.ge/lean.tar.gz.part01} through \code{part14}, in numerical order. \\
Combined compressed size & 239,275,554 bytes. \\
Regular archive members & 175,190: 175,187 Lean files, a toolchain file, and two Lake configuration files. \\
Uncompressed regular-file bytes & 1,030,536,719. \\
Pinned Lean version & \path{leanprover/lean4:v4.33.1}. \\
Pinned Mathlib revision & \path{0df444a360eaa60ab8c11dca51a86af692955474}. \\
\bottomrule
\end{tabular}
\end{center}
The SHA-256 hashes identify the examined version; subsequent changes at the public URLs would require a new audit.
\begin{description}
\item[Announcement:] \mbox{}\\[-5pt]
{\small\ttfamily 3fc97784b71ddd44ff8aa0a177e99622a1c\\61b13cfa843a0989ee4a7a0cfeb8d}
\item[Combined archive:] \mbox{}\\[-5pt]
{\small\ttfamily 9631b4f020199e580ba86178cb56c6e9b09\\f690c2369a80dc4334aa93f49c7fc}
\end{description}
The displayed line breaks are not part of the hashes. The accompanying reproducibility files contain the unbroken hashes, per-part hashes, source-inventory checks, exact finite tests, and build records.

All local source references below are relative to the archive's \path{lean/src/} directory. Their full paths appear in the source bibliography. Reading every file programmatically to form an import graph is explicitly distinguished from reading and checking every mathematical argument.

\section{The actual conclusions}
The file \code{Statements.lean} imports Mathlib and defines propositions asserting
\begin{align}
&\zeta(s)=0,\quad s\notin\{-2,-4,-6,\ldots\},\quad s\ne1
       &&\Longrightarrow\quad \Rea s=\tfrac12, \label{eq:rh}\\
&\Rea s>\tfrac12,\quad (\chi\ne1\ \text{or}\ s\ne1)
       &&\Longrightarrow\quad L(s,\chi)\ne0, \label{eq:grh}
\end{align}
where $\chi$ ranges over all complex Dirichlet characters modulo every positive integer. It also gives the critical-strip conclusion, both for all characters and for primitive characters. The endpoint \code{Main.lean} claims these propositions and Mathlib's \code{RiemannHypothesis}. The pinned Mathlib definition of the latter has the usual exclusions in~\eqref{eq:rh}. \src{S-endpoint,mathlib}

Thus the inspected endpoint is not merely a statement about a newly defined substitute for $\zeta$ or $L(s,\chi)$. No endpoint redefinition by notation was identified in the syntax checks. This paragraph concerns the meaning of the declarations. The later complete export comparison and the published independent proof check are recorded separately in Section~\ref{sec:complete-export-verification}. The exposed conclusions concern RH and Dirichlet GRH. They are not statements of every Hecke or Artin version sometimes included under the name GRH.

\subsection{Imprimitive characters do not supply a counterexample}
\begin{proposition}
Let $\chi$ modulo $N$ be induced by a primitive character $\chi^{*}$ of conductor $f$. The extra Euler factors in passing from $L(s,\chi^{*})$ to $L(s,\chi)$ have no zero in $\Rea s>0$.
\end{proposition}
\begin{proof}
The standard identity is
\[
 L(s,\chi)=L(s,\chi^{*})
 \prod_{\substack{p\mid N\\p\nmid f}}
       \bigl(1-\chi^{*}(p)p^{-s}\bigr).
\]
For $\sigma=\Rea s>0$, each $p\nmid f$ has $|\chi^{*}(p)|=1$, and consequently
$|\chi^{*}(p)p^{-s}|=p^{-\sigma}<1$.
No factor can vanish. A zero of an added Euler factor lies on $\Rea s=0$, outside the open critical strip.
\end{proof}
This is the reduction used in the inspected character argument. \src{S-character}

\subsection{Reflection for complex characters}
For a primitive nonprincipal character, put $a\in\{0,1\}$ with $\chi(-1)=(-1)^a$, and
\[
 \Lambda(s,\chi)=\left(\frac f\pi\right)^{(s+a)/2}
                  \Gamma\!\left(\frac{s+a}{2}\right)L(s,\chi).
\]
In $0<\Rea s<1$, the gamma factor is finite and nonzero. The functional equation relates $\Lambda(s,\chi)$ to a nonzero constant times $\Lambda(1-s,\overline\chi)$. Therefore a zero with $\Rea s<1/2$ produces a zero for $\overline\chi$ with real part greater than $1/2$. Nonvanishing for \emph{all} characters on the latter half-plane gives the critical line conclusion. It is unnecessary, and in general incorrect, to assume that the zeros of one complex character are invariant under conjugation without changing the character. The inspected source changes the character correctly. \src{S-character}

\section{Assessment of the proposed method}
The announcement describes a multiplicity argument suggested by Stepanov's method over finite fields. In its proposed number-field analogue, a chain of theta lifts of length comparable with $\sqrt X$ is meant to provide the required multiplicity. Norm estimates, divisibility, comparisons of integral structures, arithmetic slopes, and bounds on negative-twist Galois cohomology then lead to an upper bound for a smoothed von Mangoldt sum. An analytic criterion is used to exclude zeros to the right of $1/2$. \src{announcement}

This outline identifies a possible architecture for a proof; its analogies do not themselves prove the required assertions. In particular, the informal finite-field identity involving a term $M^k$ and a remainder does not specify all the varieties, sheaves, smoothness conditions, and normalizations needed to formulate one universal theorem. Refuting an arbitrary trace function or an arbitrary singular quadric would not, by itself, refute the construction intended in the archive.

The most consequential obligations are as follows.
\begin{center}
\begin{tabular}{@{}p{3.1cm}p{10cm}@{}}
\toprule
Part of the method & What has to be proved \\
\midrule
Iterated theta lifts & Nonvanishing, the local correspondences and infinity types, and the precise effect of every lift on the relevant periods. \\
Growing ranks & Bounds with explicit dependence on ranks, weights, conductors, fields, and the number of lifts. A theorem with an unspecified constant for a fixed group is insufficient by itself. \\
Integral structures & Denominator control, saturation, lattice indices, and compatibility under change of coefficient field. A complex inner-product identity alone does not give these. \\
Cohomological estimates & The exact global and local conditions, Tate twists, parity, and sufficiently uniform bounds on torsion lengths. \\
Analytic conclusion & The sign of the estimate, the length and quantifiers of the intervals, and nonvanishing of the Euler and gamma factors used in reductions. \\
\bottomrule
\end{tabular}
\end{center}
The source does contain explicit statements addressing these obligations. The report cannot claim that an obligation is absent merely because it is not explained in the four-page announcement.

\subsection{Where these arguments meet}
The final character theorem is assembled in
\path{CharacterBoxes/MainConclusion.lean}. Its main inputs are a Galois representation and local--global compatibility statement, a quantitative cohomology assertion, and the theta-integral and covolume estimates assembled in
\code{EveryCharacterOfLiteratureOfVWin.lean}. The latter file has a large list of imports and short theorems supplying the various input propositions. Those short theorems are references to further proofs, not independent explanations of them. Their names, including references to published literature, are not a substitute for checking the further proofs. \src{S-assembly}

The inspected analytic chain then passes through the character families, bounds on the smoothed prime sum, and a Landau-type criterion. The strongest independent positive checks below concern several implications within this chain and the finite-dimensional algebra used in the covolume argument. They do not establish the difficult existence theorems supplying their hypotheses.

\section{An independent proof of the final numerical inequality}\label{sec:numeric}
The following implication is a translation of the scalar estimate used near the analytic end of the argument. Its explicit lower bounds on the ranks and on the number of stages prevent the most immediate degenerate examples. \src{S-size}

\begin{proposition}\label{prop:size}
Let $x,\varepsilon,C,n,n_T,n_2,m,b,k,A,E,G$ be real numbers satisfying
\[
 x\ge3,\qquad \varepsilon>0,\qquad C\ge0,\qquad n>0,\qquad b>0,
 \qquad 0\le A\le2xb,
\]
\[
 n_T\ge\frac n2x^{-\varepsilon/8},\qquad
 x^{-\varepsilon}\le\frac{n_2}{n}\le1,\qquad
 x^{1/2-\varepsilon}\le m\le x^{1/2+\varepsilon/2},
\]
\[
 xm\le k\le n^2x^{1/2-\varepsilon},\qquad
 x^{3\varepsilon/8}\ge26.
\]
Suppose also that
\begin{align*}
E\le x^{C\varepsilon}\biggl[
&n_TA\left(1+\frac{m^2}{k}\right)\\
&+m\left(2nmbx^\varepsilon+
 \frac{n_2kbx^\varepsilon}{n^2x^{-\varepsilon}}\right)\biggr]
\end{align*}
and
\[
 n_Tmb\frac{n_2}{2n}\bigl(G-4x^{1/2+\varepsilon}\bigr)
       \le E+\frac{2n_Tmb}{x}.
\]
Then $G\le x^{1/2+(C+3)\varepsilon}$.
\end{proposition}
\begin{proof}
Set $Q=n_Tmb>0$, $d=n_2/n>0$, and $Y=x^{1/2+13\varepsilon/8}$. Dividing the expression in square brackets by $Q$ gives exactly
\[
 \frac{A}{mb}+\frac{Am}{kb}
 +\frac{2nmx^\varepsilon}{n_T}
 +\frac{n_2kx^{2\varepsilon}}{n^2n_T}.
\]
The respective bounds are
\begin{align*}
 \frac{A}{mb}&\le\frac{2x}{m}\le2x^{1/2+\varepsilon},\\
 \frac{Am}{kb}&\le\frac{2xm}{k}\le2,\\
 \frac{2nmx^\varepsilon}{n_T}&\le4mx^{9\varepsilon/8}
                  \le4x^{1/2+13\varepsilon/8},\\
 \frac{n_2kx^{2\varepsilon}}{n^2n_T}
       &\le\frac{2k}{n^2}x^{17\varepsilon/8}
                  \le2x^{1/2+9\varepsilon/8}.
\end{align*}
Each displayed power is at most $Y$, and $Y\ge1$. Hence $E/Q\le10x^{C\varepsilon}Y$. The second assumed inequality implies
\[
 d\bigl(G-4x^{1/2+\varepsilon}\bigr)
       \le 2E/Q+4/x
       \le22x^{C\varepsilon}Y,
\]
because $4/x\le4/3\le2x^{C\varepsilon}Y$. Dividing by $d$ and using $d^{-1}\le x^\varepsilon$ gives
\begin{align*}
 G&\le4x^{1/2+\varepsilon}
           +22x^{1/2+(C+21/8)\varepsilon}\\
  &\le26x^{1/2+(C+21/8)\varepsilon}
   \le x^{1/2+(C+3)\varepsilon}.
\end{align*}
The last step is precisely $26\le x^{3\varepsilon/8}$.
\end{proof}

This calculation checks the implication for arbitrary data satisfying its hypotheses, not just a numerical sample. It does not establish the construction of such data. In particular, it does not prove the required lower bounds for $n_T$, $n_2$, or $m$, or the estimate for $E$.

The archive itself contains a counterexample to a weaker error estimate. That weaker estimate is explicitly negated in its source. The inspected final chain uses a stronger assertion, so the internal counterexample is not a contradiction in the claimed proof. \src{S-weak}

\section{The analytic criterion: signs and intervals}\label{sec:landau}
The abstract analytic input uses a real measurable function $g(u)$, initially with $|g(u)|\le C e^{2u}$ for $u\ge0$. Its Laplace transform is assumed to extend meromorphically to $\Rea s>0$, with isolated simple poles; the parameter $\theta$ satisfies $0<\theta<1$. A real pole has a \emph{positive real residue}; a nonreal pole has a nonzero residue. The upper estimate concerns a nonnegative majorant $R\ge g$ on intervals
\[
 [U,(1+\delta)U],\qquad \delta>0\ \text{fixed},
\]
available arbitrarily far out, with a bound of the form
\[
 \int_U^{(1+\delta)U}R(u)e^{-\theta u}\,du\le e^{\eta U}
\]
for every fixed $\eta>0$. In the original variable $x=e^u$, these intervals are $[Y,Y^{1+\delta}]$. \src{S-landau}

An independent proof of the abstract analytic implication is supplied below in Proposition~\ref{prop:landau-independent}. Its application still requires the correct Laplace-transform identities and the arithmetic construction of the interval estimates. First, the following sign issue is useful. If $w\ge0$ and the integrals exist, then
\[
 \left|\int g(u)w(u)e^{-i\gamma u}\,du\right|
       +\int g(u)w(u)\,du
 \le\int\bigl(|g(u)|+g(u)\bigr)w(u)\,du.
\]
This follows from the triangle inequality. On the chosen interval,
$|g|+g=2\max(g,0)\le2R$. Thus a one-sided bound can control the expression used in the argument. It need not be replaced by a bound on $|g|$.

\begin{example}[Why the sign of a real residue matters]
Fix $0<\theta<\beta<1$ and let $g(u)=-e^{\beta u}$. Its Laplace transform is
\[
 \int_0^\infty g(u)e^{-su}\,du=-\frac1{s-\beta}
 \qquad(\Rea s>\beta).
\]
The majorant $R=0$ satisfies every displayed one-sided interval estimate, although there is a pole at $\beta>\theta$. Its residue is $-1$. This refutes the criterion \emph{without} its positive-real-residue hypothesis, and is excluded by the actual hypothesis.
\end{example}

\begin{example}[Why a fixed interval ratio matters]
Let $g(u)=e^{\beta u}\cos u$, with $0<\theta<\beta<1$. Its transform is
$(s-\beta)/((s-\beta)^2+1)$, with simple poles at $\beta\pm i$ and residues $1/2$. There are arbitrarily large intervals of fixed bounded length on which $g\le0$, so $R=0$ works there.

This does not give the required estimates on proportional intervals. Once $\delta U>3\pi$, the interval $[U,(1+\delta)U]$ contains an interval of length $\pi/3$ with $\cos u\ge1/2$, beginning between $U$ and $U+2\pi$. Any nonnegative majorant therefore satisfies
\[
 \int_U^{(1+\delta)U}R(u)e^{-\theta u}\,du
 \ge\frac{\pi}{6}e^{(\beta-\theta)U}.
\]
For $0<\eta<\beta-\theta$ this exceeds $e^{\eta U}$ for all sufficiently large $U$. Hence the actual proportional-interval hypothesis fails.
\end{example}
The source fixes $\delta$ before passing to arbitrarily large intervals. It also allows constants in the auxiliary analytic estimates to depend on that fixed $\delta$. Neither of the preceding examples establishes a flaw in the stated criterion.

% Optional self-contained insertion for the main report.
\subsection{An independent proof of the one-sided criterion}

\begin{proposition}\label{prop:landau-independent}
Let $g:[0,\infty)\to\R$ be measurable and satisfy
$|g(u)|\le C_0e^{2u}$. Suppose that its Laplace transform
\[
 \Phi(s)=\int_0^\infty g(u)e^{-su}\,du\qquad(\Rea s>2)
\]
extends meromorphically to $\Rea s>0$, with only isolated simple poles,
and that every real pole has positive real residue.
Fix $0<\theta<1$ and $\delta>0$. Suppose that for every $\eta>0$
there are arbitrarily large $U$ and a nonnegative function $R_U\ge g$
on $[U,(1+\delta)U]$ such that $R_U(u)e^{-\theta u}$ is integrable there and
\[
 \int_U^{(1+\delta)U}R_U(u)e^{-\theta u}\,du\le e^{\eta U}.
\]
Then $\Phi$ has no pole with real part greater than $\theta$.
\end{proposition}

\begin{proof}
We give the tail estimate and the pole argument explicitly.

Put $a=1+\delta/2$ and $b=1+\delta$. Both numbers
\[
 q_- = \frac{\exp(1-1/a)}{a},\qquad
 q_+ = \frac ba\exp(1-b/a)
\]
are strictly less than $1$. For the first, use
$\log a>1-1/a$ for $a>1$; for the second, use
$\log r<r-1$ for $r=b/a>1$. Choose $\lambda>0$ so large that
\[
 \lambda a\ge1,\qquad
 q:=\max(q_-,q_+)<\frac{\lambda}{\lambda+2},
 \qquad \sigma:=\lambda+2>2.
\]
For $U\ge1$, set $k=\lceil\lambda aU\rceil$, so that $k\ge U$.
Elementary exponential tilting and
$\int_0^\infty u^ke^{-tu}\,du=k!t^{-k-1}$ give
\begin{align*}
 \int_0^U\frac{u^ke^{-\lambda u}}{k!}\,du
 &\le \frac{e^{\lambda(a-1)U}}{(\lambda a)^{k+1}}
 \le \frac1{\lambda a}\left(\frac{q_-}{\lambda}\right)^k,\\
 \int_{bU}^\infty\frac{u^ke^{-\lambda u}}{k!}\,du
 &\le \frac{e^{-\lambda(b-a)U}}{(\lambda a/b)^{k+1}}
 \le \frac{b\exp((b-a)/a)}{\lambda a}
          \left(\frac{q_+}{\lambda}\right)^k.
\end{align*}
For the first inequality, write the integrand as
$e^{\lambda(a-1)u}u^ke^{-\lambda au}/k!$ and bound the first factor
by its value at $U$. For the second, write it as
$e^{-\lambda(1-a/b)u}u^ke^{-(\lambda a/b)u}/k!$ and bound the first
factor by its value at $bU$. The final bounds use
$U\le k/(\lambda a)$ and $U\ge(k-1)/(\lambda a)$, respectively.
Thus, for a constant $C_1$ depending only on the chosen $\delta,\lambda$,
\begin{equation}\label{eq:landau-tail-independent}
 \int_{[0,\infty)\setminus[U,bU]}
      \frac{u^ke^{-(\sigma-2)u}}{k!}\,du
 \le C_1\sigma^{-k}.
\end{equation}

For $\Rea s>2$, differentiation under the integral is justified by
the exponential bound on $g$. Write
\[
 A_k(s)=\frac{(-1)^k}{k!}\Phi^{(k)}(s)
       =\int_0^\infty\frac{g(u)u^ke^{-su}}{k!}\,du.
\]
The triangle inequality gives, for every real $\gamma$,
\[
 |A_k(\sigma+i\gamma)|+\Rea A_k(\sigma)
 \le\int_0^\infty (|g(u)|+g(u))
               \frac{u^ke^{-\sigma u}}{k!}\,du.
\]
On $[U,bU]$, use $|g|+g\le2R_U$ and
\[
 \frac{u^ke^{-(\sigma-\theta)u}}{k!}
 \le(\sigma-\theta)^{-k},
\]
which follows from the power series for the exponential.
Off this interval, use $|g|+g\le2C_0e^{2u}$ and
\eqref{eq:landau-tail-independent}. Since $\theta>0$, there is a constant
$C_2>0$, independent of $U,k,\eta,\gamma$, such that, for every
$\eta>0$ and arbitrarily large suitable $U$,
\begin{equation}\label{eq:landau-coeff-independent}
 |A_k(\sigma+i\gamma)|+\Rea A_k(\sigma)
 \le C_2\frac{e^{\eta U}}{(\sigma-\theta)^k}
 \quad\text{for all }\gamma\in\R,
 \qquad k\ge U.
\end{equation}

Suppose, for a contradiction, that a pole $\rho_0$ has
$\Rea\rho_0>\theta$. All poles have real part at most $2$, since the
original Laplace integral is holomorphic on $\Rea s>2$.
Let $Z$ be the pole set. Every compact subset of $\Rea s>0$ meets $Z$
in a finite set. If the closed disk with centre $\sigma$ and radius
$\sigma-\theta$ contains a pole, let $D$ be the distance to its nearest
pole; otherwise put $D=\sigma-\theta$. In either case
\[
 0<D\le\sigma-\theta,
 \qquad B(\sigma,D)\cap Z=\varnothing.
\]

First suppose there is a pole $\rho_1$ with
$\sigma-\Rea\rho_1<D$. There is then a real $\gamma$ for which
$z=\sigma+i\gamma$ has a unique nearest pole $\rho$ at distance $d<D$.
Here is the elementary justification. Start with
$\gamma=\Im\rho_1$ and vary it in a sufficiently small interval so
that $\rho_1$ remains at distance less than $D$. Only finitely many
poles can enter the disks of radius $D$ under this variation. For two
distinct such poles, equality of their distances from
$\sigma+i\gamma$ occurs at at most one value of $\gamma$: squared
distances have a difference affine in $\gamma$, and if their imaginary
parts coincide their different real parts, both less than $\sigma$,
give unequal distances for every $\gamma$. Avoid these finitely many
exceptional values. Choose $r$ with $d<r<D$ such that $B(z,r)$ contains
only $\rho$, and choose $r_1$ with $d<r_1<r$.

If $c\ne0$ is the residue at $\rho$, subtracting $c/(s-\rho)$ makes
$\Phi$ holomorphic on $B(z,r)$ after filling the removable singularity.
Cauchy's estimates consequently give constants $M,M'>0$ such that
\[
 \left|A_k(z)-\frac{c}{(z-\rho)^{k+1}}\right|\le Mr_1^{-k},
 \qquad |A_k(\sigma)|\le M'r_1^{-k}
 \quad(k\ge0).
\]
The second estimate uses the pole-free disk $B(\sigma,D)$.
Combining these inequalities with~\eqref{eq:landau-coeff-independent}
gives
\begin{equation}\label{eq:landau-final-independent}
 \frac{|c|}{d}\,d^{-k}
 \le (M+M')r_1^{-k}
       +C_2e^{\eta U}(\sigma-\theta)^{-k}.
\end{equation}

It remains to consider the case in which
$D\le\sigma-\Rea\rho$ for every pole $\rho$.
The pole-free closed-disk alternative is impossible because of
$\rho_0$. A nearest pole $\rho$ must satisfy
\[
 D\le\sigma-\Rea\rho\le|\sigma-\rho|=D.
\]
Hence it is the real pole $\beta=\sigma-D$, and it is the unique nearest
pole. Moreover $\beta>\theta$, since
$D\le\sigma-\Rea\rho_0<\sigma-\theta$.
Choose $D<r_1<r\le\sigma-\theta$ so that $B(\sigma,r)$ contains no
other pole. Its positive residue $c$ gives
\[
 \Rea A_k(\sigma)\ge cD^{-k-1}-Mr_1^{-k}
\]
by the same Cauchy estimate. Equation~\eqref{eq:landau-coeff-independent},
with $\gamma=0$, gives~\eqref{eq:landau-final-independent} with
$d=D$, $|c|=c$, and $M'=0$.

In either case $0<d<r_1\le\sigma-\theta$. Fix
$0<\eta<\log(r_1/d)$. Since $k\ge U$, multiplying
\eqref{eq:landau-final-independent} by $d^k$ shows
\[
 0<\frac{|c|}{d}
 \le (M+M'+C_2)
      \exp\bigl((\eta-\log(r_1/d))U\bigr).
\]
The right-hand side tends to zero along the arbitrarily large admissible
$U$, which is a contradiction.
\end{proof}

This proves the abstract analytic implication under its stated
hypotheses. It does not prove the arithmetic construction of the
majorants or the interval estimates for the character families.

\section{The cohomological assertions and their quantifiers}\label{sec:cohomology}
The arithmetic bound is not an assertion that all class groups, higher $K$-groups, or negative-twist cohomology groups vanish. Its exact formulation matters.

For a number field $E$ and a rational prime $\ell$, let $G_{E,S_\ell}$ denote the Galois group of the maximal extension unramified outside the places above $\ell$ and the infinite places. The group used in the source is
\[
 \Sha^1_\ell(E,-e)=
 \ker\left(
 H^1(G_{E,S_\ell},\Q_\ell/\Z_\ell(-e))
 \longrightarrow
 \prod_{v\mid\ell}H^1(E_v,\Q_\ell/\Z_\ell(-e))
 \right).
\]
The coefficients are discrete, with the cyclotomic action of exponent $-e$. Write $v_\ell(E,e)$ for its length as a $\Z_\ell$-module. The condition denoted by \code{VVle E t B} requires both
\begin{enumerate}
\item $v_\ell(E,e)<\infty$ for every prime $\ell$ and every integer $1\le e<t$;
\item for every finite set $S$ of rational primes,
\[
 \sum_{\ell\in S}(\log\ell)
       \sum_{1\le e<t}v_\ell(E,e)\le B.
\]
\end{enumerate}
An infinite extended length is not merely converted to the natural number zero and then used in the estimate: the finiteness requirement is a separate conjunct. \src{S-cohomology,S-vwin}

\subsection{The field family and the interval}
Fix a primitive Dirichlet character $\psi$ of conductor $f$. For an integer $h\ge0$ and auxiliary primes $q,d$, the field $E$ is the cyclotomic subfield cut out by the simultaneous conditions
\[
 a\equiv\pm1\pmod{2^{h+2}},\qquad
 \psi(a)=1,\qquad
 \left(\frac aq\right)=\left(\frac ad\right)=1.
\]
Equivalently, it is the compositum of the maximal real subfield of
$\Q(\mu_{2^{h+2}})$, the field fixed by $\ker\psi$, and the two quadratic fields specified by the Jacobi characters. The prescribed congruences
$q\equiv3\pmod8$ and $d\equiv7\pmod{16}$ make the latter two fields
$\Q(\sqrt{-q})$ and $\Q(\sqrt{-d})$. This field is therefore totally imaginary.

For a real parameter $V$, the condition on the auxiliary primes also requires
\[
 q\mid d+1,\qquad
 p\mid d+1\quad
 \text{for every odd prime }p\le V^{\varepsilon/5},
\]
\[
 \frac12V^{\varepsilon/5}<\log q\le V^{\varepsilon/5},
\qquad
 \frac12V^{\varepsilon/5}(\log V)^2<\log d
       \le V^{\varepsilon/5}(\log V)^2.
\]
These requirements are part of the statement being audited; an arbitrary pair of primes cannot replace $q,d$.

The principal cohomology assertion, with exponent $401$, has the following order of quantifiers. Fix $\psi$ and $0<\varepsilon<1/52$. For each nonnegative integer-valued choice $h:\R\to\Z_{\ge0}$ satisfying
\[
 x^{2+8\varepsilon}\le 2^{h(x)}
             \le2x^{2+8\varepsilon}\qquad(x\ge1),
\]
there exist a real number $J\ge1$ and $X_0$ such that, for every $X\ge X_0$ and every pair of primes satisfying the preceding conditions with $V=X^{1/J}$, there exists $Y$ with
\[
 X^{1/J}\le Y,\qquad Y^2\le X,
\]
for which, for every $x\in[Y,Y^2]$ and its prescribed field $E$,
\[
 \operatorname{VVle}
   \left(E,\lfloor x^{401}\rfloor,\exp(x^\varepsilon/2)\right)
\]
holds. The constants may depend on the fixed choices made before them. They are not claimed to be uniform as $\psi$ varies or $\varepsilon\to0$. \src{S-vwin}

Neither an isolated nonzero cohomology group nor a polynomially large contribution, by itself, refutes this eventual exponential bound. A counterexample would have to satisfy the field and auxiliary-prime conditions and violate the quantified estimate. On the other hand, uniformity throughout $[Y,Y^2]$, across the allowed pairs $q,d$, and across all finite sets of primes is genuinely required. This audit has not proved that uniform estimate.

\subsection{Parity prevents the first irregular-prime objection}
In the characterwise part of the proof the coefficient action is
\[
 \psi^{-1}\chi_{\mathrm{cyc}}^{-e},
\]
and the difficult case explicitly assumes $\psi(-1)=(-1)^e$. If $c$ is complex conjugation, then
\[
 \bigl(\psi^{-1}\chi_{\mathrm{cyc}}^{-e}\bigr)(c)
   =\bigl((-1)^e\bigr)^{-1}(-1)^{-e}=1.
\]
Thus the relevant character is even. Odd eigenspaces occurring in the usual irregular-prime examples cannot be substituted for it.

For example, the character exponent associated with the Bernoulli index $32$ at $\ell=37$ is $1-32\equiv5\pmod{36}$, which is odd. With $\psi=1$, the corresponding negative twist has $e=31$, and the hypothesis $1=(-1)^{31}$ is false. This elementary parity check excludes that proposed example; it does not prove the vanishing of any even eigenspace. \src{S-counts}

The separate assertion \code{PropZ} is restricted still further:
\[
 \ell\ge5,\qquad e\ge2,\qquad e\text{ even},
 \qquad \ell-1\mid2e.
\]
It is split into the cases $\ell-1\mid e$ and the resulting quadratic-character alternative. It is not a statement of unrestricted class-group vanishing. An independent proof of this restricted assertion is supplied in Section~\ref{sec:further-arithmetic}, including the passage from divisible coefficients to residual coefficients with the local conditions preserved.

\subsection{Zero twists, real places, and rational coefficients}
The inspected comparison with Poitou--Tate duality has the explicit guard
\[
 \ell\ne2\quad\text{or}\quad r_1(E)=0.
\]
The invariant-module finiteness statements require $e\ne0$. The infinite untwisted module $\Q_\ell/\Z_\ell$ and unaccounted-for real-place terms at $\ell=2$ therefore do not refute the versions inspected here. Adding an extra zero local restriction is correctly used to make the relevant Selmer subgroup smaller. \src{S-cohomology,S-counts}

The relevant modular form has weight $k=2e+1$, with $e\ge1$, and $W(-e)$ is attached to it. Thus $W$ corresponds to the integer twist $e$, up to finite-order characters. This is a critical noncentral twist: $1\le e\le k-1$ and $e\ne k/2=e+1/2$. The dual critical index $e+1$ is noncentral as well. Kato's introductory statement distinguishes finiteness at noncentral critical twists from the central case, where a nonvanishing hypothesis is required. The obvious objection based on a positive central analytic rank does not address this application. Its conclusion is the vanishing of a \emph{rational} Bloch--Kato Selmer space; finite integral torsion does not contradict it. \src{S-kato,kato}

Similarly, the inspected Ferrero--Washington assertion is made for abelian number fields, with a bound allowed to depend on the field and prime. It is not the unproved assertion of a vanishing Iwasawa $\mu$-invariant for arbitrary nonabelian fields. \src{S-fw}

\subsection{Substantial arithmetic claims that remain unverified}
One auxiliary assertion supplies an absolute exponent $C_K$ such that, for a totally complex number field $L$ containing $\mu_\ell$ (and $\mu_4$ when $\ell=2$), and every $e\ge1$,
\[
 H^2(\OO_L[1/\ell],\Z_\ell(e+1))
 \quad\text{is finite}
\]
and
\[
 \log\#H^2(\OO_L[1/\ell],\Z_\ell(e+1))
 \le
 \exp\left((2e[L:\Q])^{C_K}
                 \log(2|\operatorname{disc}L|)\right).
\]
Its size makes mere torsion nonvanishing irrelevant as a counterexample. Its absolute uniformity nevertheless needs a proof, and that complete proof has not been independently reconstructed here. \src{S-lemmaK}

The refined bounds also use rank-three and rank-four deformation and automorphy constructions. The inspected files provide reductions to ordinary deformation rings, modular ideals, local vanishing, residual cases, and tower enlargement. These are supplied arguments with further dependencies, rather than explicit axioms. Their full validity has not been established by this audit. In particular, the construction of the modular ideals and the quantitative passage from those constructions to the cohomological interval estimate remain major verification tasks. \src{S-automorphy}

One possible objection to the rank-three argument can be tested algebraically. A general induced three-dimensional representation can become reducible on an open subgroup. The inspected argument, however, uses a residual character value $b$ with $b^2\ne1$ and reduces the alleged failure of irreducibility to
\[
 a b^n+1+(a b^n)^{-1}=0\qquad(n\ge0),
 \qquad a,b\ne0.
\]
This system is impossible when $b^2\ne1$. Indeed, all the elements $ab^n$ would be roots of the nonzero quadratic polynomial $T^2+T+1$. The three elements $a,ab,ab^2$ are distinct if $b^2\ne1$, contradicting the degree bound for roots of a polynomial over a field. Thus a generic example of automorphic induction does not refute this particular reduction; one must also meet the residual separation hypotheses. \src{S-irreducibility}

\section{Integral lattices, projectors, and covolumes}\label{sec:integral}
This part of the audit tests whether a rational direct-sum decomposition is silently treated as an integral one. Such an error would be serious. In the inspected step, the index is present and controlled.

\subsection{The denominator of a central idempotent}
\begin{proposition}\label{prop:idempotent}
Let $G$ be a finite group and let $K$ be a number field. If $e$ is a central idempotent of $K[G]$, then
\[
 |G|e\in\OO_K[G].
\]
\end{proposition}
\begin{proof}
After extension to a characteristic-zero splitting field, a central idempotent is a sum of primitive central idempotents. The primitive idempotent associated with an irreducible character $\chi$ is
\[
 e_\chi=\frac{\chi(1)}{|G|}
          \sum_{g\in G}\chi(g^{-1})g.
\]
Every value $\chi(g)$ is a sum of roots of unity, because a finite-order matrix is diagonalizable with roots of unity as eigenvalues in characteristic zero. It is therefore an algebraic integer. Every coefficient of $|G|e$ is an algebraic integer. Those coefficients already belong to $K$, so they belong to $\OO_K$, as required.
\end{proof}
The factor $|G|$ cannot generally be removed: for the trivial character of a group of order greater than one, the averaging idempotent is $|G|^{-1}\sum_{g\in G}g$. The inspected central-idempotent argument retains the denominator. \src{S-idempotent}

\subsection{Removing a small orthogonal summand}
Let $K$ have degree $d$ and $r_2(K)$ pairs of complex embeddings. For an $\OO_K$-lattice $A$ of rank $s$ with specified Hermitian norms, use the normalized covolume
\[
 h(A)=2^{r_2(K)s}|\operatorname{disc}K|^{-s/2}
                   \operatorname{covol}(A).
\]
This includes the factor used in the source's totally complex setting. Define, for $s>0$,
\[
 \mu_2(A)=-\frac{2}{ds}\log h(A).
\]
Let $\mu_{2,\max}(M)$ be the supremum of $\mu_2(A)$ over nonzero sublattices $A\subseteq M$, endowed with the induced norms.

\begin{proposition}\label{prop:slope}
Let $M$ be an $\OO_K$-lattice in a Hermitian $K$-vector space, let $0\ne L\subseteq M$ be a sublattice, and put $V=KL$. Let $P$ be a $K$-linear idempotent of $V$ which is orthogonal at every infinite place. Suppose that $g\ge1$ is an integer and $gP(L)\subseteq L$. Put
\[
 L_+=L\cap\im P,\qquad L_-=L\cap\ker P,
 \qquad r_\pm=\rank_{\OO_K}L_\pm,\qquad r=r_++r_-.
\]
Suppose that $0\le\delta\le1/12$ and
\[
 r_-\le3\delta r,\qquad
 \mu_{2,\max}(M)\le Q,\qquad
 \mu_2(L)\ge-Q,\qquad
 2\log g\le Q.
\]
Then $r_+>0$ and
\[
 \mu_2(L_+)\ge\mu_2(L)-12\delta Q.
\]
\end{proposition}
\begin{proof}
First, $Q\ge0$ because $g\ge1$. The denominator condition implies
$g(1-P)L\subseteq L_-$ and $gPL\subseteq L_+$. Thus $L_+\oplus L_-$ has finite index $I$ in $L$. The map
\[
 L/(L_+\oplus L_-)\longrightarrow g^{-1}L_-/L_-,
 \qquad x\longmapsto(1-P)x\bmod L_-,
\]
is well defined and injective. If $(1-P)x\in L_-$, then
$Px=x-(1-P)x$ belongs to $L\cap\im P=L_+$. Hence
\[
 I\le \#(g^{-1}L_-/L_-)=g^{dr_-}.
\]
The equality follows by regarding the lattice $L_-$ as a free abelian group of rank $dr_-$.

Orthogonality and the covolume-index identity give
\[
 h(L_+)h(L_-)=h(L_+\oplus L_-)=Ih(L).
\]
When $r_->0$, taking logarithms yields
\[
 r\mu_2(L)=r_+\mu_2(L_+)+r_-\mu_2(L_-)
                      +\frac2d\log I.
\]
Since $L_-\subseteq M$, one has $\mu_2(L_-)\le Q$. Also
\[
 \frac2d\log I\le2r_-\log g\le r_-Q.
\]
As $r_-\le r/4$, we have $r_+\ge3r/4>0$ and
\[
 \frac{r_-}{r_+}\le
       \frac{3\delta}{1-3\delta}\le4\delta.
\]
Solving the slope identity for $\mu_2(L_+)$ now gives
\begin{align*}
 \mu_2(L_+)
 &\ge \mu_2(L)
      +\frac{r_-}{r_+}\mu_2(L)
      -\frac{r_-}{r_+}Q-\frac{r_-}{r_+}Q\\
 &\ge \mu_2(L)-3\frac{r_-}{r_+}Q
 \ge \mu_2(L)-12\delta Q.
\end{align*}
If $r_-=0$, then $\ker P=0$, so $P=1$ and $L_+=L$; the conclusion is immediate.
\end{proof}

The source's application bounds both $2\log |G|$ and the relevant ambient slopes. Thus the integrality loss in Proposition~\ref{prop:idempotent} has a place in the inequality; it is not silently set to zero. \src{S-cut}

For example, if $\delta=e^{-x^\alpha}$ and $Q$ is bounded by a fixed power of $x$, the loss $12\delta Q$ is smaller than any prescribed negative power of $x$ eventually. Indeed, for fixed $A,B,\alpha>0$,
\[
 x^{A+B}e^{-x^\alpha}\longrightarrow0,
\]
since $x^\alpha-(A+B)\log x\to+\infty$. This checks the final use of an exponentially small exceptional-rank fraction. It does not prove the estimates that produce that fraction.

\subsection{Intersection and change of coefficient field}
\begin{proposition}\label{prop:flat}
Let $L/K$ be an extension of number fields, $R=\OO_K$, $S=\OO_L$, $V$ a finite-dimensional $K$-vector space, $M\subset V$ an $R$-lattice, and $W\subset V$ a $K$-linear subspace. Inside $L\otimes_K V$, one has
\[
 SM\cap (L\otimes_K W)=S(M\cap W).
\]
\end{proposition}
\begin{proof}
The $R$-module $S$ is torsion-free. Since $R$ is a Dedekind domain, $S$ is flat over $R$. This can be checked after localization at each nonzero prime, where $S$ is a finite torsion-free, hence free, module over a discrete valuation ring.

The kernel of $M\to V/W$ is $M\cap W$. Tensoring this exact sequence with $S$ therefore identifies
\[
 S\otimes_R(M\cap W)
 =
 \ker\bigl(S\otimes_R M\longrightarrow S\otimes_R(V/W)\bigr).
\]
Flatness also preserves the injection $M\hookrightarrow V$. Moreover
$S\otimes_R K=L$, and consequently the last kernel, viewed inside
$L\otimes_K V$, is exactly
$SM\cap(L\otimes_K W)$. This proves the identity.
\end{proof}
Ramification of $L/K$ does not invalidate this proof. In fact, flatness also gives $SM\cap SN=S(M\cap N)$ for arbitrary $R$-submodules $M,N$ of $V$. The danger in the application is different: replacing an integral submodule by its $K$-span can lose an index. The inspected source supplies a rational subspace and separately checks the field of definition of the relevant character. \src{S-rational}

\begin{example}[Why an arbitrary integral submodule is different]
Let $R$ be a discrete valuation ring with uniformizer $\pi$, fraction field $K$, and lattice $M=R\subset K$. If the proposed subspace is instead the integral submodule $N=\pi R$, then $KN=K$ and
\[
 M\cap KN=R\ne\pi R=N.
\]
Thus replacing the integral submodule $N$ by its $K$-span loses an index before any change of coefficient field. The archive contains a counterexample to an integral-submodule formulation of its lattice restriction. Its actual rational-subspace argument does not make this replacement.
\end{example}

\subsection{The small-section input is stronger than these algebraic checks}
The principal geometric input asserts that certain small theta sections span the relevant space over its coefficient field. Write $[E:\Q]=2f$ and let $r$ be the unitary rank, with signature $(p,q)$. In the inspected formulation the hypotheses include $p,q\ge1$, $f>2$, an even rank $n\ge8fr+4$, the dominance conditions, and inequalities $\lambda_{\tau,k}>n/2$ for $k<p$ at the infinite places. It also bounds the level and field size and assumes a specified good cover. A section norm is bounded by an expression of the form
\[
 \exp\bigl(C(frL)^2\kappa\bigr).
\]
The statement asks for spanning over the coefficient field, not generation of the integral module by those sections. An objection based solely on failure of integral generation therefore does not refute that statement. \src{S-gmin}

The source gives further theta-spanning, archimedean norm, and denominator-clearing arguments. In the inspected theta-spanning reduction, a reproducing-kernel calculation and a strict bound for a sum of off-diagonal terms lead to a Neumann-series inversion. This audit has not independently verified all the analytic and geometric hypotheses of that reduction. Neither Propositions~\ref{prop:slope} and~\ref{prop:flat} nor the observation about field-linear spanning proves this input.

\subsection{The strict estimate in the theta-spanning argument}
The occurrence of a contraction constant $0.65$ in the theta-spanning argument initially suggests a possible missing dependence on the level or the number of lattices. The underlying count has a direct explanation.

\begin{proposition}\label{prop:latticecount}
Let $V$ be a rational vector space of dimension $\rho$, with a fixed nondegenerate alternating form. For a full lattice $N$, put
\[
 N^\#=\{v\in V:\langle v,N\rangle\subseteq\Z\}.
\]
Fix a self-dual lattice $L=L^\#$. For each $m\ge1$, there are at most $m^{2\rho}$ self-dual lattices $L'$ satisfying $[L:L\cap L']=m$. Consequently, if $n\ge4\rho+4$,
\[
 \sum_{L'\ne L}[L:L\cap L']^{-n/2}<\frac{13}{20}.
\]
\end{proposition}
\begin{proof}
A free abelian group of rank $\rho$ has at most $m^\rho$ subgroups of index $m$. For prime $m=p$, each such subgroup is the kernel of a nonzero map to $\Z/p\Z$, giving the bound $p^\rho$. For composite $m$, choose a prime $p\mid m$. Every subgroup $H$ of index $m$ is contained in a subgroup $K$ with $[K:H]=p$ and $[L:K]=m/p$, because the finite abelian group $L/H$ has a subgroup of order $p$. Induction bounds the number of pairs $(K,H)$ by $(m/p)^\rho p^\rho=m^\rho$.

Put $N=L\cap L'$. There are at most $m^\rho$ choices for $N$. Duality reverses finite-index inclusions and preserves indices, so $[N^\#:L]=m$. Self-dual lattices have the same covolume, whence $[L':N]=m$ and $[N^\#:L']=m$. For each $N$ there are at most $m^\rho$ possible subgroups $L'$ of index $m$ in $N^\#$. This proves the count. If $m=1$, inclusion and duality force $L'=L$.

Grouping by $m$ now gives
\[
 \sum_{L'\ne L}[L:L\cap L']^{-n/2}
 \le\sum_{m=2}^{\infty}m^{2\rho-n/2}
 \le\sum_{m=2}^{\infty}\frac1{m^2}.
\]
An exact integral-test estimate is
\[
 \sum_{m=2}^{\infty}\frac1{m^2}
 <\sum_{m=2}^{10}\frac1{m^2}+\frac1{10}
 =\frac{825257}{1270080}<\frac{13}{20};
\]
the difference in the last inequality is $59/254016>0$.
\end{proof}

The application has $\rho=[E:\Q]r$. The overlap statement used in the source is
\[
 \|\Phi_{L,\Lambda}\|^2=|J|,
 \qquad
 |\langle\Phi_{L,\Lambda},\Phi_{L',\Lambda}\rangle|
 \le |J|\,[L:L\cap L']^{-n/2}.
\]
The same finite level factor $|J|$ occurs on the diagonal and off it. The normalized overlap of the underlying invariant vectors is calculated by averaging an orthonormal orbit indexed by a lattice quotient. If that quotient has order $d$, the average has squared norm $1/d$, giving an overlap of absolute value $d^{-1/2}$. Tensoring with a projective module of rank $n$ raises the relative index to the $n$th power, as can be checked locally over discrete valuation rings. The source uses the same semicharacter on both lattices; choosing incompatible characters on their intersection would change this calculation and does not satisfy the hypothesis. \src{S-theta-count}

A related field hypothesis is also essential. In $\Q(i)$ every integral element is $a+bi$ with $a,b\in\Z$, so its trace is $2a$; there is no integral element $w$ with $w+\overline w=1$. This would refute a construction asserting such an element for every CM field. The source's level structure, however, explicitly requires this element, and its actual field family supplies
\[
 w=\frac{1+\sqrt{-d}}2,\qquad d\equiv7\pmod{16}.
\]
The polynomial $T^2-T+(d+1)/4$ is monic with integer coefficients and vanishes at $w$. Thus $w$ is integral, and complex conjugation gives $w+\overline w=1$. The application uses this element in its level data. The example $\Q(i)$ therefore fails an explicit hypothesis and does not refute that construction. \src{S-theta-count}

The weighted finite genus sum contributes a positive mass $\mathfrak m$ to both sides. Thus, for its coefficient matrix $C$, the diagonal is
$a_0=\mathfrak m|J|>0$ and
\[
 \sum_{L'\ne L}|C(L,L')|\le(13/20)a_0.
\]
This is sufficient for the subsequent inversion. On bounded families with norm
$\|F\|=\sup_{L,x}h(x)|F_L(x)|$, where $h$ is positive and continuous, set
\[
 (KF)_L(x)=-a_0^{-1}\sum_{L'\ne L}C(L,L')F_{L'}(x).
\]
The sum converges absolutely and $\|K\|\le13/20<1$. Hence
\[
 F'=a_0^{-1}\sum_{j=0}^{\infty}K^jF,\qquad
 CF'=F,\qquad
 \|F'\|\le\frac{\|F\|}{a_0(1-13/20)}.
\]
Since $h^{-1}$ is bounded on compact subsets, the convergence is locally uniform and preserves holomorphicity. A covariance law preserved by $C$ is likewise preserved by this inverse. No uniform positive lower bound on $\mathfrak m$ is needed for this existence assertion; the spanning conclusion allows unrestricted coefficients.

These checks explain the strict estimate and the inversion. They do not prove the theta kernel identity, the definite Hermitian Siegel--Weil formula, or the explicit reproducing formula which turn the inverse into a theta-spanning statement. Those are precisely the deeper inputs still unverified here.

\section{Cyclotomic orders and the auxiliary combinatorics}\label{sec:behrend}
The auxiliary selection of stages uses orders in real cyclotomic Galois groups and restrictions on three-term linear relations. These are useful places to look for off-by-one errors and degenerate counterexamples. The following checks concern the statements in \path{Formalism/BehrendSets/}; they do not independently prove every asymptotic existence estimate in that development. \src{S-behrend}

\subsection{Orders modulo a power of two}
\begin{proposition}
For an integer $s\ge0$, put
\[
 G_s=(\Z/2^{s+2}\Z)^\times/\{\pm1\}.
\]
If $a$ is an odd integer, let $d_s(a)$ be its order in $G_s$. For every integer $j\ge0$,
\[
 d_s(a)\le2^j
 \quad\Longleftrightarrow\quad d_s(a)\mid2^j
 \quad\Longleftrightarrow\quad
 a\equiv\pm1\pmod{2^{\max(s+2-j,0)}}.
\]
If $a\ge2$ and $d_s(a)\le J$, then $2^{s+2}\le J(a+1)$.
If $a\equiv7\pmod{16}$, then $d_s(a)=2^{\max(s-1,0)}$.
\end{proposition}
\begin{proof}
Every class modulo sign has a unique representative congruent to $1$ modulo $4$. For an odd integer $b\equiv1\pmod4$ with $b\ne1$, repeated squaring gives
\[
 v_2(b^{2^t}-1)=v_2(b-1)+t.
\]
For completeness, the induction step follows from
$b^{2^{t+1}}-1=(b^{2^t}-1)(b^{2^t}+1)$ and
$v_2(b^{2^t}+1)=1$. It follows also that the class of $5$ has order $2^s$ modulo $2^{s+2}$, so the group of representatives congruent to $1$ modulo $4$ is cyclic of order $2^s$.

Choose $b=\pm a$ with $b\equiv1\pmod4$. Unless $b=1$, the order of its class is
\[
 2^{\max(0,s+2-v_2(b-1))}.
\]
For $j<s$ this formula gives the claimed equivalence. For $j\ge s$, the order divides $2^j$ and the congruence modulus is at most $4$, so both conditions hold automatically. The case $b=1$ is immediate.

Write $d_s(a)=2^t$ with $0\le t\le s$. The equivalence just proved shows that $2^{s+2-t}$ divides one of $a-1,a+1$. If $a\ge2$, both are nonzero, so $2^{s+2-t}\le a+1$. Multiplication by $2^t\le J$ proves the second assertion. Finally, for $a\equiv7\pmod{16}$ take $b=-a$; then $v_2(b-1)=3$, which gives the last formula, including $s=0,1$.
\end{proof}

\subsection{Spacing in a short interval}
\begin{proposition}
Let $n,L$ be positive integers, and let $T$ be a set of nonnegative integers with
$t<n/(8L)$ for all $t\in T$. Suppose $T$ contains no nonconstant solution to
\[
 \alpha x+\beta y+\gamma z=0,
 \qquad \alpha+\beta+\gamma=0,
\]
where $\alpha,\beta,\gamma$ are nonzero integers of absolute value at most $2L$.
If $h\in\Z/n\Z$ has order greater than $2L$, a translate of $T$ contains at most two of the terms $u+ih$ when $i$ ranges over any $2L$ consecutive integers.
\end{proposition}
\begin{proof}
Suppose three distinct times $i_1<i_2<i_3$ give representatives $t_1,t_2,t_3\in T$ after the fixed translation. Put $a=i_2-i_1$ and $b=i_3-i_1$. Then $0<a<b<2L$ and
\[
 b(t_2-t_1)-a(t_3-t_1)\equiv0\pmod n.
\]
The absolute value of this integer is less than
$(a+b)n/(8L)<n$, so it is zero. Its coefficients $a-b,b,-a$ are nonzero, have sum zero, and have absolute value at most $2L$. The hypothesis on $T$ gives $t_1=t_2=t_3$. Then $ah=0$ with $0<a<2L<\ord(h)$, a contradiction.
\end{proof}

If instead $\ord(h)=d\le8L$, every coset of $\langle h\rangle$ meets $T$ in at most one point. Indeed, $\langle h\rangle$ is the subgroup of multiples of $n/d$, so distinct representatives in one coset differ by at least $n/d\ge n/(8L)$. Distinct elements of $T$ have smaller difference. For $L=0$ the first window is empty; the hypothesis $d\le8L$ of the second statement is impossible. Thus the boundary cases do not yield exceptions.

\subsection{Random subset estimates}
Let a random subset $T$ of a finite ground set $X$ be formed by retaining each point independently with probability $\delta\in[0,1]$. Let $\mathcal W$ be a finite family of sets of cardinality at most $L$, and fix a point $t\in X$ and a nonnegative integer $K$. For a single set $W$ containing $t$, the event
\[
 t\in T,\qquad |T\cap W|\ge K+1
\]
requires the retention of $t$ and some $K$ distinct other points in $W$. A union bound gives, for $L\ge1$,
\[
 \Pr(t\in T,\ |T\cap W|\ge K+1)
 \le \binom{L-1}{K}\delta^{K+1}
 \le\delta(L\delta)^K.
\]
If $t\notin W$, independence and the same argument give the bound
\[
 \binom{L}{K+1}\delta^{K+2}
 \le\delta(L\delta)^{K+1}.
\]
Points of $W$ outside $X$ cannot belong to $T$ and do not spoil either inequality. Summing over $\mathcal W$ gives the more precise bound according to whether $t\in W$. If $L\delta\le1$, the common upper bound is
\[
 \delta\,|\mathcal W|\,(L\delta)^K.
\]
Empty sets, $\delta=0$, and $K=0$ can be checked directly, with the usual convention $0^0=1$ for this polynomial formula.

For the additional selection assertion, let $Z=|T|$ and let $Y(T)\ge0$ be any real-valued function. Then
\[
 \mathbf EZ=\delta|X|,\qquad
 \operatorname{Var}(Z)=|X|\delta(1-\delta).
\]
If $\lambda>0$ and $2\operatorname{Var}(Z)<\lambda^2$, Chebyshev's inequality gives
\[
 \Pr(|Z-\delta|X||>\lambda)<\tfrac12.
\]
Markov's inequality gives
$\Pr(Y>2\mathbf EY)\le1/2$; when $\mathbf EY=0$ this event has probability zero. The union of the two events has probability strictly below one. Consequently there is a subset $T$ satisfying
\[
 \bigl||T|-\delta|X|\bigr|\le\lambda,
 \qquad Y(T)\le2\mathbf EY.
\]
These proofs verify the relevant finite probabilistic implications without relying on a numerical approximation.

\subsection{Exact bounded counterexample searches}
The accompanying Python program uses integers and \code{fractions.Fraction}. It found no counterexample in the following finite searches.
\begin{center}
\begin{tabular}{@{}lr@{}}
\toprule
Assertion tested & Cases checked\\
\midrule
Cyclotomic order equivalences & 95,000\\
Order lower bound & 9,990\\
Order for $a\equiv7\pmod{16}$ & 1,250\\
At most two points in a window & 7,402,538\\
At most one point in a short orbit & 39,864\\
Refined random-subset bound & 29,580\\
Common random-subset bound & 15,948\\
Simultaneous size and expectation bound & 1,704\\
\midrule
Total & 7,595,874\\
\bottomrule
\end{tabular}
\end{center}
The order tests use $0\le s\le9$, odd $1\le a\le1999$, and $0\le j\le s+4$. The spacing tests use $1\le n\le64$, $1\le L\le6$, all eligible short subsets, all group elements, and all window translates. The probability tests use ground sets of size at most four and $\delta\in\{0,1/4,1/2,3/4,1\}$; they exhaust all families through size three and families of at most two subsets at size four. The simultaneous selection tests use ground sets of size at most two, every function $Y$ with values in $\{0,1,2\}$, and $\lambda\in\{1/4,1/2,1,2,3\}$ when the variance hypothesis holds.

The program and its output specify the ranges exactly. These finite searches are additional evidence. The preceding proofs, rather than the number of tested cases, establish the universally quantified elementary statements.

\section{Theta lifting and a temperedness counterexample}\label{sec:tempered}
Theta lifting does not give nonzero lifts for arbitrary source representations and arbitrary target spaces. Local occurrence and special values in inner-product formulas matter, as do the range and cuspidality hypotheses. The primary references inspected state these qualifications explicitly. Likewise, increasing the dimension of positive definite spaces does not by itself give a Witt tower; a Witt tower is obtained by adding hyperbolic planes. \src{gan-qiu-takeda,gan}

The source's inspected nonvanishing input \code{TLiftK} assumes a CM extension, admissible splitting data, a positive definite source Hermitian space, sufficiently large target dimension, prescribed Hecke characters, and inequalities
\[
 d+m_{W,v}\le2\lambda_{v,k}
\]
at every relevant infinite place and weight component. Its assertion is therefore not unrestricted nonvanishing of every theta lift. The rank and weight restrictions must be checked in its application; their presence rules out a counterexample which simply chooses a representation outside that range. \src{S-tlift}

\begin{example}[Definite unitary cuspidality does not imply temperedness]
Let $E=\Q(i)$ and let $G$ be the unitary group of
\[
 h(x,y)=x\overline{x}+y\overline{y}.
\]
This Hermitian form is positive definite at the real place. It has no nonzero rational isotropic vector, so $G$ has no proper rational parabolic subgroup. The constant automorphic function is consequently cuspidal according to the vanishing-of-constant-terms definition.

The prime $5$ splits in $E$, and $G(\Q_5)\simeq\mathrm{GL}_2(\Q_5)$. The local representation generated by the constant function is trivial. Its matrix coefficient is identically $1$, so it is not tempered: modulo the centre, this coefficient is not in $L^{2+\epsilon}$ for any $\epsilon>0$. Indeed, $\mathrm{PGL}_2(\Q_5)$ has infinite Haar measure, since the distinct diagonal Cartan classes give infinitely many disjoint translates of a compact open subgroup of positive measure.
\end{example}

This is a counterexample to the assertion that every cuspidal automorphic representation on a definite unitary group is tempered at finite places. It is also the obstruction discussed in Gan's notes. \src{gan}

The archive's inspected theorem \code{LocalComponentTempered} concerns a regular algebraic, conjugate-self-dual \emph{cuspidal} automorphic representation of $\mathrm{GL}_n$ over a CM field, with $n\ge2$, together with the specified Galois representation and local--global compatibility input. The example above does not satisfy those hypotheses. In particular, transfer from a definite unitary group need not be cuspidal on $\mathrm{GL}_n$. The source's base-change statements inspected here allow the relevant discrete isobaric factors. Thus the example is not a demonstrated flaw in that chain. \src{S-tempered}

\section{Abstract models and their actual use}
Some source files admit models with zero ranks, zero-valued functions, or empty auxiliary types. Such models merit scrutiny when a complicated geometric statement is expressed through an abstract interface. They constitute a defect only if the actual final application can obtain a required nonzero geometric quantity from an unrelated degenerate model.

Two relevant uses were traced. First, a default local object is used at indices which are not primes; the prime-index branch supplies the actual local construction. Second, an existence construction for prescribed points replaces the period lattice and its form by the genuine specified lattice and form. The later divisibility expression inspected in that branch is built from this prescribed lattice. It does not use the auxiliary zero-valued function of the abstract model. \src{S-models}

The source also contains a proposition named \code{Sw} defined to be \code{False}. A theorem with hypothesis $s:\code{Sw}$ and proof by elimination of $s$ is simply a vacuous conditional statement. It would yield an inconsistency only if a proof of \code{Sw} were obtained and used. No such derivation was established in the audited chain. Likewise, auxiliary propositions explicitly proving that a weaker estimate is false do not make a strengthened estimate false.

These observations concern the particular uses traced, not every abstract interface in the archive. No semantic substitution of the zero lattice for the required genuine period lattice has been demonstrated.


\section{Further arithmetic vanishing and nonvanishing checks}\label{sec:further-arithmetic}
The next two assertions have been proved independently during the continued audit. The subsequent nonzero Selmer class is a test of stronger vanishing claims; it does not contradict the submitted polynomial bounds.

\subsection{Independent proof of the Bernoulli-number bound used in \texttt{LemmaX}}

The statement checked here is
\texttt{UpperTheory.V.X.lemmaX} in
\path{LevelCohomology/LemmaX.lean}, with the hypotheses
\texttt{XData} from
\path{LevelCohomology/UnramifiedCocycleIsCoboundaryOfBernoulliUnit.lean}.
This is a verification of a substantive input, not a flaw in that input.

\paragraph{Proposition.}
Let $\ell$ be an odd prime, let $d\geq1$, and suppose that
$4\mid\ell^d-1$.  Put
\[
 q=\ell^d,\qquad a=v_2(q-1)\geq2.
\]
Let $f$ be an odd positive integer and let $\varphi$ be a Dirichlet
character modulo $f$ whose order divides $q-1$.  Let $\xi$ be a
primitive character of conductor $2^n$, with
\[
 \varphi(-1)\xi(-1)=-1,
 \qquad f\,4^a\leq2^n.
\]
Regard their values as elements of $\overline{\mathbf Q}_\ell$.
Then the generalized Bernoulli number $B_{1,\varphi\xi}$ is an
$\ell$-adic unit.  In particular, this holds when $\varphi$ is
primitive, as required by the submitted statement.

\paragraph{Proof.}
Set
\[
 N=2^n,\qquad m=n-a,\qquad s=f\,2^a.
\]
The size hypothesis implies $m\geq a\geq2$ and $s\leq2^m$.
Let $\mathcal O$ be the valuation ring of
$\overline{\mathbf Q}_\ell$, with residue field
$\overline{\mathbf F}_\ell$.  All character values are integral.
The reductions of the values of $\varphi$ belong to $\mathbf F_q$.
Reduction preserves the orders of the values of $\xi$, since these
orders are powers of two.

We first remove the possible denominator $f$ from the Bernoulli sum;
this also covers the case $\ell\mid f$.  For $u\in\mathbf Z/N\mathbf Z$,
write
\[
 A(u)=\sum_{\substack{0\leq b<fN\\ b\equiv u\pmod N}}b\varphi(b),
 \qquad
 T=\sum_{b=0}^{f-1}\varphi(b),
 \qquad
 c=\frac{T}{N},
 \qquad
 g(u)=\frac{A(u)-A(0)}{fN}.
\]
The Chinese remainder theorem shows that
$\sum_{0\leq b<fN,\ b\equiv u\bmod N}\varphi(b)=T$ for every $u$.
Shifting the summation interval by $ft$ and using the periodicity of
$\varphi$ gives, for every integer $t\geq0$,
\[
 g(u)-g(u-ft)
 =tc-\sum_{\substack{0\leq b<ft\\b\equiv u\pmod N}}\varphi(b).
 \tag{1}
\]
For completeness, the corresponding identity before division is
\[
 A(u)-A(u-ft)
 =ftT-fN\sum_{\substack{0\leq b<ft\\b\equiv u\pmod N}}\varphi(b).
\]
It follows by comparing the sum of $b\varphi(b)$ over
$0\leq b<fN$ with its sum over $ft\leq b<fN+ft$; the terms outside
the common part differ by $fN\varphi(b)$, while translating the
linear factor contributes $ftT$.

Since $f$ is invertible modulo $N$, for any $u$ there exists $t\geq0$
with $ft\equiv u\pmod N$.  Equation~(1), together with $g(0)=0$,
expresses $g(u)$ as an integral linear combination of $c$ and the
values of $\varphi$.  The denominator of $c$ is a power of two.
Consequently $g(u)\in\mathcal O$ and its reduction belongs to
$\mathbf F_q$ for every $u$.

If $r(u)$ denotes the representative of $u$ in $[0,N)$, applying
(1) with $t=2^a$ gives
\[
 g(u)-g(u-s)
 =2^ac-
 \begin{cases}
   \varphi(r(u)),&r(u)<s,\\
   0,&r(u)\geq s.
 \end{cases}
 \tag{2}
\]
The product character $\varphi\xi$ is odd and hence nontrivial.
The usual Bernoulli sum formula and $\sum_{u\bmod N}\xi(u)=0$ give
\[
 B_{1,\varphi\xi}
 =\frac1{fN}\sum_{b=0}^{fN-1}b\varphi(b)\xi(b)
 =\sum_{u\bmod N}\xi(u)g(u).
 \tag{3}
\]
In particular, this Bernoulli number is integral.

We prove that its reduction is nonzero.  In the rest of the proof,
write $\xi,\varphi,g,c$ for their reductions.  Suppose, to the contrary,
that the right side of (3) is zero.  Define
\[
 U=\{x\in(\mathbf Z/N\mathbf Z)^\times:
                    \xi(x)^{2^a}=1\}.
\]
For every unit $y$, changing variables shows that
\[
 \sum_{x\bmod N}\xi(x)g(xy)=0.
 \tag{4}
\]
The subgroup generated by $q$ in $(\mathbf Z/N\mathbf Z)^\times$
is $1+2^a\mathbf Z/N\mathbf Z$.  Indeed, $q=1+2^a v$ with $v$ odd,
and the elementary formula
$v_2(q^{2^j}-1)=a+j$ shows that its order modulo $2^n$ is
$2^{n-a}$, the order of that subgroup.
The formula follows by induction from
$q^{2^{j+1}}-1=(q^{2^j}-1)(q^{2^j}+1)$, since the second factor
has $2$-adic valuation one.
Raise (4) to powers $q^r$.  Since all values of $g$ are in
$\mathbf F_q$, it follows, for every integer $j\geq0$, that
\[
 \sum_{x\bmod N}\xi(x)^{1+2^aj}g(xy)=0.
\]
Sum this identity over $0\leq j<2^n$.  The geometric sum is $2^n$
when $\xi(x)^{2^a}=1$ and is zero otherwise.  As $2^n$ is invertible
in characteristic $\ell$, we obtain
\[
 \sum_{x\in U}\xi(x)g(xy)=0
 \quad\hbox{for every unit }y.
 \tag{5}
\]

Primitivity of $\xi$ means that $\xi(5)$ has order $2^{n-2}$.
Writing an odd residue as $(-1)^\epsilon5^j$ therefore shows that
\[
 x\in U\quad\Longrightarrow\quad x\equiv1\ \hbox{or}\ -1\pmod{2^m}.
 \tag{6}
\]
The two possibilities are distinct, because $m\geq2$.  Write
$\eta(x)=1$ or $-1$ according to these two possibilities.
Since $2^a\mid s$, relation (6) implies
$xs=s$ or $-s$ modulo $N$, respectively.  Thus (2) gives
\[
 \xi(x)\bigl(g(x)-g(x(1-s))\bigr)
 =\eta(x)\xi(x)2^ac
       -\mathbf1_{x=1}-\mathbf1_{x=-1}
 \qquad(x\in U).
 \tag{7}
\]
Here is the endpoint check in this identity.  On the residue class
$1\pmod{2^m}$, the interval $[0,s)$ contains only $1$ among the
possible representatives, since $s\leq2^m$, and $\varphi(1)=1$.
On the residue class $-1\pmod{2^m}$ one applies (2) at $x+s$;
the interval contains only the representative $s-1$, which occurs
precisely when $x=-1$.  Since $f\mid s$,
\[
 \varphi(s-1)\xi(-1)=\varphi(-1)\xi(-1)=-1.
\]
These observations prove (7), including the boundary case $s=2^m$.

Finally, put $u_0=1+N/2$.  Primitivity gives $\xi(u_0)=-1$.
Multiplication by $u_0$ permutes $U$ and preserves $\eta(x)$,
because $u_0\equiv1\pmod{2^m}$.  Consequently
\[
 \sum_{x\in U}\eta(x)\xi(x)=0.
\]
Sum (7) over $U$.  Its left side is zero by (5), applied to
$y=1$ and $y=1-s$; the latter is a unit since $s$ is even.
Its right side is $-2$.  This contradicts $\ell\ne2$.
Thus the reduction of (3) is nonzero, proving that
$B_{1,\varphi\xi}$ is an $\ell$-adic unit.\hfill$\square$

\paragraph{Exact computation used as a separate check.}
The accompanying program \path{check_lemma_x.py} uses only exact
integer sums and polynomial arithmetic over finite fields.  Its
coefficient primes are the odd primes $3\leq\ell\leq101$, its odd
conductors are $f=1$ and the primes $3\leq f\leq31$, and
$4\leq n\leq11$.  There are $2200$ quadratic-character checks,
of which $617$ satisfy the proposition's size hypothesis, and
$1305$ additional checks of nontrivial $\mathbf F_\ell$-valued
characters with that hypothesis.  Some of the latter repeat
quadratic cases; these numbers count checks, not distinct examples.
Every check simultaneously includes all primitive $\xi$ of the
specified parity, by taking the greatest common divisor of the
Bernoulli-sum polynomial and $X^{2^{n-3}}+1$ over $\mathbf F_\ell$.
No counterexample satisfies the hypothesis.

There are $18$ detected nonunits below the threshold in the quadratic
family.  For example, for $f=1$, $n=6$, and $\ell=17$, the polynomial is
\[
 P(X)=-1-X-X^2+X^3+X^4+X^5-X^6+X^7,
\]
and
\[
 \gcd(P(X),X^8+1)\equiv X+6\pmod{17}.
\]
Here $a=v_2(17-1)=4$, so the hypothesis would require $n\geq8$.
The detected nonunit is therefore a check on the computation,
not a counterexample to the submitted bound.

\src{S-bernoulli-continued}

\subsection{The restricted cyclotomic vanishing assertion \texttt{PropZ}}

The assertion \texttt{PropZ} in
\path{HeightAndVanishing/Counts.lean} assumes
$\ell\geq5$, $e\geq2$, $e$ even, and $\ell-1\mid2e$.
The theorem \path{LevelCohomology/PropZ.lean} splits it into
\texttt{zOne} in \path{LevelCohomology/TorsClass.lean} and
\texttt{zQuad} in \path{LevelCohomology/ZqField.lean}.
The following proof independently verifies this exact vanishing
statement.  It does not assert Vandiver's conjecture for the other
even characters of the cyclotomic Galois group.

\paragraph{Proposition.}
Let $\ell\geq5$ be prime and let $e\geq2$ be even, with
$\ell-1\mid2e$.  Let $\mathcal O$ be the ring of integers of a
finite extension $L/\mathbf Q_\ell$, and put
\[
 M=(L/\mathcal O)(\varepsilon_\ell^{-e}).
\]
Fix a decomposition group $D_\ell$ at $\ell$, and inertia groups
$I_p$ at all finite primes.  Then
\[
 S(M):=\ker\left(H^1(\mathbf Q,M)\longrightarrow
       H^1(D_\ell,M)\times\prod_{p\ne\ell}H^1(I_p,M)\right)=0.
\]
Cohomology here means continuous cohomology with discrete coefficients.
This is the group whose $\mathcal O$-length is
\texttt{Chr.V l e 1}; the source uses all conjugate local groups,
which gives the same condition.

\paragraph{Proof.}
Every cohomology class in $H^1(\mathbf Q,M)$ has finite
$\ell$-power order.  Indeed, a continuous cocycle from the compact
group $G_{\mathbf Q}$ to the discrete torsion group $M$ has finite
image, and a single power of $\ell$ kills that image.  Thus it is
enough to show that $S(M)[\ell]=0$.

Set $W=M[\ell]$.  The action on $W$ is the scalar action of
\[
 \chi=\overline\varepsilon_\ell^{-e}:
 G_{\mathbf Q}\longrightarrow\mathbf F_\ell^\times.
\]
As an $\mathbf F_\ell$-vector space $W$ is finite dimensional; it
is a direct sum of copies of this one-dimensional character.  This
description remains valid when $L/\mathbf Q_\ell$ is ramified:
$W$ identifies with $\mathcal O/\ell\mathcal O$, and the action
comes from the scalar subfield $\mathbf F_\ell$.

We first justify that an $\ell$-torsion class satisfying the local
conditions has a lift satisfying the corresponding local conditions
with coefficients in $W$.  The cyclotomic image of $D_\ell$ is the
same as that of $G_{\mathbf Q}$, namely $\mathbf Z_\ell^\times$.
For example, each extension
$\mathbf Q(\mu_{\ell^r})/\mathbf Q$ is totally ramified at $\ell$
of degree $\ell^{r-1}(\ell-1)$, as follows from the Eisenstein
polynomial $\Phi_{\ell^r}(X+1)$ over $\mathbf Q_\ell$; passing to
the inverse limit proves the image assertion.  Since the action
on $M$ factors through this cyclotomic image,
\[
 H^0(\mathbf Q,M)=H^0(D_\ell,M)
 \quad\hbox{as subgroups of }M.
 \tag{1}
\]

Let $x\in S(M)[\ell]$.  The long exact sequence associated with
\[
 0\longrightarrow W\longrightarrow M
   \xrightarrow{\ \ell\ }M\longrightarrow0
\]
gives $y\in H^1(\mathbf Q,W)$ mapping to $x$.
As $x|_{D_\ell}=0$, the class $y|_{D_\ell}$ is a connecting
homomorphism image of an element of $H^0(D_\ell,M)$.
By (1), subtracting the global connecting homomorphism image of
the same element arranges that
\[
 y|_{D_\ell}=0.
 \tag{2}
\]
For $p\ne\ell$, inertia acts trivially on the full module $M$.
Multiplication by $\ell$ on $M$ is surjective, so
\[
 H^0(I_p,M)/\ell H^0(I_p,M)=M/\ell M=0.
\]
The same long exact sequence therefore makes
$H^1(I_p,W)\to H^1(I_p,M)$ injective.  Since the latter image of
$y$ is zero, $y|_{I_p}=0$.  Thus $y$ is unramified at all finite
primes, and it is zero on $D_\ell$.

The congruence $\ell-1\mid2e$ makes $\chi$ either trivial or
quadratic.  If it is trivial, the class $y$ is a continuous
homomorphism $G_{\mathbf Q}\to W$, unramified at all finite primes.
It is also unramified at the real place, because a homomorphism
from a group of order two to an $\ell$-group is zero.  By the
unramified part of global class field theory it factors through
the ideal class group of $\mathbf Q$, which is trivial.  Hence
$y=0$ in this case.

Suppose now that $\chi$ is nontrivial.  Then
\[
 e\equiv\frac{\ell-1}{2}\pmod{\ell-1}.
\]
Since $e$ is even, $\ell\equiv1\pmod4$.  The fixed field of
$\ker\chi$ is the quadratic subfield
\[
 K=\mathbf Q(\sqrt\ell)\subset\mathbf Q(\mu_\ell),
 \qquad \operatorname{disc}(K)=\ell.
\]
Restriction
$H^1(\mathbf Q,W)\to H^1(K,W)$ is injective: the kernel is
$H^1(\operatorname{Gal}(K/\mathbf Q),W)$, which is zero since
the group order two is invertible on $W$.
Over $K$ the module $W$ is trivial.  The restriction of $y$ is
therefore an unramified continuous homomorphism $G_K\to W$.
Again its restrictions at real places are zero because $\ell$
is odd.  It consequently factors through the ordinary ideal
class group $\operatorname{Cl}(K)$.

Here that class group has order strictly less than $\ell$.
One elementary bound suffices.  Minkowski's theorem supplies an
integral ideal of norm at most $\sqrt\ell/2$ in each ideal class.
Put $B=\lfloor\sqrt\ell/2\rfloor$.  In a quadratic field the number
of integral ideals of norm $r$ is at most the divisor function
$\tau(r)$: at a split prime power $p^j$ there are $j+1$ ideals,
at a ramified prime power there is one, and at an inert prime
power there are at most one.  By multiplicativity the bound
follows for every $r$.  Since $\tau(r)\leq r$,
\[
 h_K\leq\sum_{r=1}^{B}\tau(r)
       \leq\frac{B(B+1)}2<\ell.
\]
For the last inequality, $\ell>4$ gives
$B\leq\sqrt\ell/2$ and $B+1<\sqrt\ell$.
It follows that $\operatorname{Cl}(K)$ has no $\ell$-torsion and
admits no nonzero homomorphism to $W$.  Hence $y|_{G_K}=0$, and
injectivity of restriction gives $y=0$.

In both cases $x=0$.  Therefore $S(M)[\ell]=0$; since every
element of $S(M)$ has $\ell$-power order, $S(M)=0$.
\hfill$\square$

\paragraph{Consequence for the counterexample search.}
An even quadratic field with nonzero $\ell$-class group can give
nonzero integral Selmer classes for other finite characters, as
the discriminant-$2021$ example shows.  It cannot contradict this
proposition: the congruence forces the only possible nontrivial
residual character to have fixed field $\mathbf Q(\sqrt\ell)$,
whose discriminant is exactly $\ell$ and whose class number is
too small to be divisible by $\ell$.  The proof also addresses
the possible loss of the residual local condition on passing
from divisible coefficients to their $\ell$-torsion subgroup.

\src{S-propz-continued,milne-ant,milne-cft}

\subsection{A nonzero class with an even quadratic character and zero local restrictions}

The following example is a test of vanishing assertions, not a counterexample
to the polynomial upper bounds in the submitted source.

Let
\[
 f(X)=X^3-8X+1,\qquad D=2021=43\cdot47,
 \qquad E=\mathbf Q(\sqrt D),
\]
and let $H$ be the splitting field of $f$.  Let $\psi$ be the quadratic
character of $E/\mathbf Q$.  For every even integer $e\geq2$, put
\[
 M_e=(\mathbf Q_3/\mathbf Z_3)(\psi^{-1}\varepsilon_3^{-e}),
\]
where $\varepsilon_3$ is the $3$-adic cyclotomic character.
There is a nonzero class in $H^1(\mathbf Q,M_e)$ which is unramified
at every finite prime and whose restriction to the full decomposition
group is zero at each prime in
\[
 \{2,3,11,43,47,263,2039,2099\}.
\]
Here unramified means that the restriction to inertia is zero in
cohomology; at $43$ and $47$ the stronger decomposition-group assertion
holds as well.

\paragraph{Proof.}
The rational-root test proves that $f$ is irreducible.  Its discriminant is
$-4(-8)^3-27=2021$, a positive squarefree nonsquare.  Therefore
$\operatorname{Gal}(H/\mathbf Q)\simeq S_3$, and its quadratic resolvent is
$E$.  In particular, $\psi$ is even and has conductor $2021$.
Squarefreeness of the polynomial discriminant implies that
$\mathbf Z[X]/(f)$ is the full ring of integers of the cubic field.
Only $43$ and $47$ ramify.  At either prime the inertia group in the
Galois closure is generated by a transposition: these primes are tame,
and discriminant exponent one gives precisely a single two-cycle.
Since inertia is normal in a decomposition group and the normalizer of
a subgroup generated by a transposition in $S_3$ is that same subgroup,
the full decomposition groups at $43$ and $47$ have order two.

Modulo $2$ and $3$ one has respectively
\[
 f=(X+1)(X^2+X+1),\qquad
 f=(X-1)(X^2+X-1),
\]
with the quadratic factors irreducible in the corresponding residue
fields.  Both primes are unramified; consequently their Frobenius
elements are transpositions and their decomposition groups in $H$
also have order two.  The primes $11$, $263$, $2039$, and $2099$ do not divide $D$ and
\[
 \left(\frac{D}{11}\right)=\left(\frac{D}{263}\right)=\left(\frac{D}{2039}\right)=
 \left(\frac{D}{2099}\right)=-1.
\]
Thus their Frobenius elements are odd permutations, hence
transpositions, and their decomposition groups also have order two.

Write $S_3=\mathbf F_3\rtimes\{\pm1\}$, with multiplication
$(a,s)(b,t)=(a+sb,st)$.  The map $z(a,s)=a$ is a one-cocycle for the
sign module $\mathbf F_3(\psi)$.  It is not a coboundary, since its
restriction to the normal subgroup $\mathbf F_3$ is nonzero whereas
every coboundary restricts to zero there.  Inflation gives a nonzero
class in $H^1(\mathbf Q,\mathbf F_3(\psi))$; nontriviality can also be
checked directly on elements mapping to the three-cycles.  On every
subgroup of order two this class is zero, because multiplication by two
is invertible on $\mathbf F_3$ and finite-group cohomology in positive
degree is killed by the group order.  Outside $43$ and $47$, inertia
maps trivially into $S_3$.  These observations give all of the asserted
local restrictions for the residual class.

Since $e$ is even and the mod-$3$ cyclotomic character has order two,
there is an isomorphism $M_e[3]\simeq\mathbf F_3(\psi)$.  Moreover
$H^0(\mathbf Q,M_e)=0$: an inertia element at $43$ acts by $-1$, because
$\psi$ is ramified there and $\varepsilon_3$ is unramified there.  Its
fixed elements are killed by two, and a $3$-primary group has no
nonzero such elements.  The long exact sequence for
\[
 0\longrightarrow M_e[3]\longrightarrow M_e
 \xrightarrow{\,3\,}M_e\longrightarrow0
\]
therefore shows that
$H^1(\mathbf Q,M_e[3])\longrightarrow H^1(\mathbf Q,M_e)$ is injective.
The nonzero residual class remains nonzero.  Every residual local
coboundary remains a coboundary after this inclusion, proving all
local assertions for $M_e$.\hfill$\square$

\paragraph{Relation to the submitted definitions.}
The definition \texttt{Chr.V'} in
\path{CharacterLevelBounds/ShaFamMono.lean}, with coefficient ring
$\mathbf Z_3$, is the length of precisely this divisible torsion
cohomology module with zero restrictions at the primes dividing $3N$
and zero inertia restrictions elsewhere; the family of local groups is
specified in \path{GaloisLibrary/SelmerTransport/ConjStableShaFam.lean}.
Taking
\[
 N=16\cdot2021\cdot11\cdot263
\]
makes the class admissible for that group.  The auxiliary prime
congruences are $11\equiv3\pmod8$ and $263\equiv7\pmod{16}$; moreover $11\mid263+1$. These primes still do not meet all the eventual size requirements of the asymptotic construction.
The induced quadratic character has the required parity for even $e$;
at $3$ its prime-to-$3$ component still has order two, and $3$ is
explicitly included in the exceptional finite set defining Case~II.
Thus the corresponding integral Selmer length is at least one.
This does not contradict a bound by a positive power of $3Ne$, and
the definition \texttt{NoEnd} concerns irreducible characteristic-zero
blocks with specified distinct weights.  The example does not by
itself meet those different hypotheses.  More explicitly,
\texttt{PropZ} in \path{HeightAndVanishing/Counts.lean} assumes the
trivial finite character and a prime $\ell\geq5$, while here $\psi$ is
nontrivial and $\ell=3$.  The assertions \texttt{CountII} and
\texttt{CountIII} are polynomial upper bounds, not vanishing
statements.  The vanishing condition \texttt{Raw.Good.xvan} in
\path{SymplecticGSp4/DecompSigmaOrd.lean} applies to the specified
odd auxiliary characters and their inverses, whereas this example has
an even residual character.  No assertion in the submitted source has
been identified that this example contradicts.

\src{S-even-selmer}

\subsection{A central extension and the determinant trace condition}

A possible objection to the deformation argument is that a nonabelian
central extension can make a cup product a coboundary while leaving its
two factors nonzero. The following calculation tests that objection
against the additional trace identity actually required in the source.
It is not a counterexample to the submitted theorem.

Let $p$ be odd, $k=\mathbf F_p$, and let $G$ consist of the tuples
$(u,v,z,\psi,\chi)\in k^3\times\{\pm1\}^2$, with product
\[
 (u,v,z,\psi,\chi)(u',v',z',\psi',\chi')
 =\bigl(u+\psi u',v+\chi\psi v',
 z+\chi z'+\chi\psi uv',\psi\psi',\chi\chi'\bigr).
\]
This is the semidirect product of the Heisenberg group
$(u,v,z)(u',v',z')=(u+u',v+v',z+z'+uv')$ with the two sign
characters. The signs act by automorphisms because
$(\psi u)(\chi\psi v')=\chi uv'$.
Put $G'=\ker\chi$ and
\[
 r(g)=\begin{pmatrix}1&0\\0&\psi(g)\end{pmatrix},\quad
 A=\begin{pmatrix}0&1\\-1&0\end{pmatrix},\quad
 B=\begin{pmatrix}0&1\\1&0\end{pmatrix},\quad
 D=\begin{pmatrix}1&0\\0&-1\end{pmatrix}.
\]
The span of $r(G')$ is the full diagonal algebra: it contains $I,D$
and $2$ is invertible. One has $A^t=-A$, $B^t=B$, $D^t=D$,
$AB=D$, $BA=-D$, and $[A,B]=2D$.

Define $\kappa(g)=u(g)A$ and $\lambda(g)=v(g)B$.
Conjugation by $r(g)$ multiplies both off-diagonal matrices by
$\psi(g)$, so the group law gives
\begin{align*}
 \kappa(gh)&=\kappa(g)+r(g)\kappa(h)r(g)^{-1},\\
 \lambda(gh)&=\lambda(g)+\chi(g)r(g)\lambda(h)r(g)^{-1}.
\end{align*}
These cocycles have the required skew-symmetric and symmetric
polarizations respectively. Neither is a coboundary. Indeed, at
$x=(1,0,0,1,1)$ the first is $A$ and every ordinary coboundary
vanishes, while at $y=(0,1,0,1,1)$ the second is $B$ and every
twisted coboundary vanishes. The same observation proves that
no nonzero scalar multiple of $\kappa$ is a coboundary.

The diagonal symmetric function $\phi_0(g)=2z(g)D$ satisfies
\begin{equation}\label{eq:heisenberg-inhom}
 \phi_0(gh)=\phi_0(g)+\chi(g)
 \left(r(g)\phi_0(h)r(g)^{-1}
       +[\kappa(g),r(g)\lambda(h)r(g)^{-1}]\right).
\end{equation}
This follows by substituting $[A,B]=2D$ and the formula for the
$z$ coordinate. Thus the central extension has removed the first
obstruction. The additional required identity, however, is
\begin{equation}\label{eq:heisenberg-trace}
 \operatorname{tr}(\phi(g)r(g))+
 \operatorname{tr}(\lambda(g)\kappa(g)r(g))=0
                         \qquad(g\in G').
\end{equation}
For $\phi_0$ its left side is
$(2z(g)-u(g)v(g))(1-\psi(g))$, which equals $4$ at
$(0,0,1,-1,1)$. It is nonzero for every odd $p$.

No other solution of \eqref{eq:heisenberg-inhom} can satisfy
\eqref{eq:heisenberg-trace}. Set $c=(0,0,1,1,1)$ and
$t=(0,0,0,-1,1)$. Then $xy=c(yx)$. The equation at $(x,y)$,
$(y,x)$, and $(c,yx)$ gives $\phi(c)=2D$: the inhomogeneous term
is $2D$ for the first pair and zero for the other two.
At $t$ and $ct$, both $\kappa$ and $\lambda$ vanish, while
$r(t)=r(ct)=D$ and $\phi(ct)=\phi(c)+\phi(t)$.
Subtracting \eqref{eq:heisenberg-trace} at these two elements
would give $0=\operatorname{tr}(2D^2)=4$, a contradiction.

The trace condition therefore detects precisely the central term
introduced by this candidate. The source imposes it on the full
subgroup $G'$, not only on the kernel of $r$. Removing that
hypothesis produces a false stronger assertion; it does not
refute the assertion actually used.

\src{S-detector-continued}

\section{Regular semisimple elements in four-dimensional representations}
\label{sec:regular-cosets}

The four-dimensional deformation argument uses the assertion
\code{AlgebraTools.regCosetFour}, through
\path{WeakShimuraTheory/CompactIntegralFields/H4.lean}.
The following statement isolates its substantive algebraic content.
It is important here to retain the hypothesis concerning
noncommuting images of finite-index normal subgroups.

\begin{proposition}\label{prop:regular-cosets}
Let $k$ be an algebraically closed field of characteristic different
from two. Let $G$ be a topological group and
$r:G\to\operatorname{GL}(V)$ a representation with $\dim_kV=4$.
Suppose that:
\begin{enumerate}
\item every polynomial zero locus in the matrix entries of $r(g)$
has closed inverse image in $G$;
\item the linear span of $r(G)$ is $\operatorname{End}_k(V)$;
\item some $u\in G$ has eigenvalues $a_1,\ldots,a_4\in k^\times$
such that $a_i/a_j$ is not a root of unity when $i\ne j$;
\item the image of every open normal subgroup of finite index in
$G$ is noncommutative.
\end{enumerate}
For every finite-index subgroup $U\subseteq G$ and every $\sigma\in G$,
the coset $\sigma U$ contains an element whose characteristic
polynomial has four distinct roots.
\end{proposition}

\begin{proof}
Let $H$ be the Zariski closure of $r(G)$ in $\operatorname{GL}(V)$,
and let $H^\circ$ be its identity component. The group $H^\circ$ is
irreducible and has finite index in $H$. Put
$N=r^{-1}(H^\circ)$. This is a normal subgroup of finite index.
It is closed by the first hypothesis: equations involving an inverse
determinant can be multiplied by a power of the determinant without
changing their zero sets on $\operatorname{GL}(V)$. A closed subgroup
of finite index is open, since its complement is a finite union of
closed cosets. Thus $r(N)$ is noncommutative.

The image of $U\cap N$ is Zariski dense in $H^\circ$.
Indeed, $r(N)$ is dense there, while $N$ is a finite union of cosets
of $U\cap N$. Their closures cover $H^\circ$; irreducibility forces
one of these translates to be all of $H^\circ$, and hence forces
the closure of $r(U\cap N)$ to be all of $H^\circ$.
Consequently finitely many polynomial functions which are each
nonzero somewhere on $r(\sigma)H^\circ$ are simultaneously nonzero
at some $r(\sigma g)$ with $g\in U\cap N$.

Choose a positive integer $b$ with $u^b\in U\cap N$ and diagonalize
$r(u)$, writing $r(u)=\operatorname{diag}(a_1,\ldots,a_4)$.
All powers of its eigenvalues remain pairwise distinct.
Let $A$ be the linear span of $r(N)$. This is a unital subalgebra of
$\operatorname{End}(V)$, normalized by $r(G)$. Polynomial interpolation
in $r(u^b)$ shows that $A$ contains the four diagonal matrix units.
Thus, for every $i,j$, either $E_{ij}\in A$, or every element of $A$
has zero $(i,j)$-entry: multiply any element of $A$ on the left and
right by $E_{ii}$ and $E_{jj}$ to see this.

The Jacobson radical $J$ of $A$ is nilpotent and is normalized by
$r(G)$. Its common annihilator in $V$ is nonzero: if $J\ne0$, take
the last nonzero member of the sequence $V,JV,J^2V,\ldots$.
That annihilator is $G$-invariant. The second hypothesis makes $V$
irreducible as a $G$-module, so $J$ annihilates all of $V$, and $J=0$.
It follows that $A$ is a direct sum of full matrix algebras on a
partition of the four coordinate lines. One can see this directly
from its matrix units: a one-way edge between distinct strongly
connected components would lie in the radical, whereas within a
component multiplication gives every matrix unit.
Conjugation by $r(G)$ permutes these blocks transitively; otherwise
the sum of a proper orbit of blocks would be $G$-invariant.
The blocks therefore have equal size. Noncommutativity excludes
four blocks of size one. The remaining possibilities are one block
of size four, or two blocks of size two.

If $r(\sigma)$ preserves the blocks, then every diagonal entry of
$r(\sigma)r(n)$, and every diagonal entry of its adjugate, is
nonzero for some $n\in N$. For the entries, use the fact that their
linear span is the full block diagonal algebra; for the adjugate,
use $\operatorname{adj}(x)=\det(x)x^{-1}$ and normality of $N$.
The density observation therefore supplies $g\in U\cap N$ for
which all eight entries are nonzero. If $r(\sigma)$ exchanges two
blocks, it instead supplies $g$ such that
\[
 r(\sigma g)=\begin{pmatrix}0&B\\ C&0\end{pmatrix},
\]
where $B,C$ are invertible $2\times2$ matrices and every entry of
both matrices is nonzero.

We next construct distinct integral weights compatible with the
multiplicative relations among $a_i^b$. Let $L\subseteq\Z^4$ consist
of the vectors $(m_i)$ for which $\prod_i(a_i^b)^{m_i}$ is a root
of unity. The quotient $\Z^4/L$ is a finitely generated torsion-free
abelian group. The images of its four coordinate vectors are
distinct, by the hypothesis on ratios. An integral linear
functional on this quotient can be chosen nonzero on each of their
six differences: in a free basis, avoid the finitely many proper
rational hyperplanes and clear denominators. Its values $w_i$
are distinct integers and satisfy
\[
 \prod_i(a_i^b)^{m_i}=1\quad\Longrightarrow\quad
 \sum_i m_iw_i=0.                                      \tag{*}
\]

Put $x=r(\sigma g)$ and $D(T)=\operatorname{diag}(T^{w_i})$.
The characteristic polynomial of $xD(T)$ is separable over $k(T)$.
Here are the needed elementary details. In the block-preserving
case, reorder the coordinates so that $w_1>w_2>w_3>w_4$, and write
the characteristic polynomial as
$Y^4+\beta Y^3+\gamma Y^2+\delta Y+\epsilon$.
Its coefficient degrees as Laurent polynomials in $T$ satisfy
\[
 \deg\beta=w_1,\quad \deg\delta=w_1+w_2+w_3,
 \quad\deg\epsilon=w_1+w_2+w_3+w_4.
\]
The first two equalities use the chosen nonzero diagonal and
adjugate entries. If the principal minor on coordinates $1,2$
is nonzero, then $\deg\gamma=w_1+w_2$; the Newton polygon at
$T=\infty$ has four segments of length one with distinct slopes,
so all four roots are distinct. If that minor is zero, then
$\deg\gamma\leq w_1+w_3$. The two outer segments have length one,
and the middle segment has length two; its intermediate point lies
strictly below the corresponding upper Newton polygon.
After passing to $k((T^{-1/2}))$ and rescaling, its residual polynomial
is $Y^2+c$ with $c\ne0$. Its two roots are distinct because
$\operatorname{char}(k)\ne2$, and Hensel's lemma separates the two
corresponding roots. The two outer roots have different valuations
from these and from each other. This proves separability.

In the block-exchanging case, order the weights in each block as
$v_1>v_2$ and $z_1>z_2$. Then the characteristic polynomial is
\[
 Y^4-t(T)Y^2+d(T),
\]
where $t(T)$ is the trace of
$C\operatorname{diag}(T^{v_1},T^{v_2})
 B\operatorname{diag}(T^{z_1},T^{z_2})$ and $d(T)$ its determinant.
All entries being nonzero gives $\deg t=v_1+z_1$; invertibility gives
$d\ne0$ and $\deg d=v_1+v_2+z_1+z_2<2\deg t$.
Thus $t^2-4d\ne0$, and $d\ne0$ and characteristic different from two
make this biquadratic polynomial separable.

Finally separability at $D(T)$ implies separability at
$D_m=\operatorname{diag}((a_i^b)^m)$ for some integer $m\geq0$.
For completeness, let $F$ be the discriminant polynomial in the four
diagonal variables. If $F(D_m)=0$ for every $m$, group its monomials
according to the values $c_\nu=\prod_i(a_i^b)^{\nu_i}$.
The equalities $\sum_\nu F_\nu c_\nu^m=0$ for all $m\geq0$ and
the Vandermonde determinant imply that the coefficient sum in
each group is zero. Relation~(*) gives the same exponent
$\sum_i\nu_iw_i$ for all monomials in that group. Hence $F(D(T))=0$,
a contradiction. Thus $r(\sigma g u^{bm})=xD_m$ is separable,
and $g u^{bm}\in U$, as required.
\end{proof}

\begin{example}[An obstruction excluded by the noncommutativity hypothesis]
Over $\C$, let
\[
 T=\{\operatorname{diag}(a,a^{-1},b,b^{-1}):a,b\in\C^\times\}.
\]
Let $W$ be the group of permutations of the four coordinates generated
by $(12)$, $(34)$, and $(13)(24)$, and let $G=T\rtimes W$ act on
$\C^4$ by the resulting monomial matrices. The action is irreducible:
a $T$-invariant subspace is a sum of coordinate lines, and $W$ is
transitive on those lines. In fact its matrix span is all of $M_4(\C)$,
by diagonal interpolation and transitivity. The element
$\operatorname{diag}(2,1/2,3,1/3)$ has all eigenvalue ratios of
infinite order. All the polynomial closedness hypotheses hold in
the usual topology.

Nevertheless, put $U=T$ and $\sigma=(12)(34)$. For every $t\in T$,
$(\sigma t)^2=1$ and each of its two $2\times2$ blocks has eigenvalues
$1,-1$. Therefore
\[
 \det(Y-\sigma t)=(Y^2-1)^2
\]
throughout this finite-index coset. No element of the coset is
regular semisimple. The missing hypothesis is precisely that the
open finite-index normal subgroup $T$ has noncommutative image.
The submitted theorem includes this hypothesis.
\end{example}

In the arithmetic application, noncommutativity is supplied by
\path{StabilizerIndexBounds/NonComm.lean}. Its assumption is that a
ratio of two residual characters has order different from $1,2,3,4$.
If a finite-index normal subgroup had commuting image, simultaneous
eigenlines and irreducibility would make the representation monomial.
Reduction of the characteristic polynomials forces all residual
characters to agree on the kernel of the permutation action.
Their ratios would factor through a subgroup of $S_4$. A field-valued
character of such a subgroup has order at most four: every
permutation has order dividing three or four; in an abelian image,
the coexistence of an element not killed by three and an element
not killed by four would give a product killed by neither.
This contradicts the stated residual ratio condition.
Thus the example does not invalidate the actual application.
The separate deformation-theoretic construction of that residual
condition is not replaced by the algebraic proof above.

\src{S-regfour,milne-ag}

\section{Formal-source integrity and actual compilation results}\label{sec:formal}
\subsection{What the static checks establish}
All 175,187 local Lean files were read programmatically to extract their import declarations. Every local file was reachable from \code{Main} in the resulting import graph. No missing local module or import cycle was found. This is \emph{module import reachability}, not proof-term dependency: an imported auxiliary theorem need not occur in the proof term of the final theorem.

The lexical and targeted metaprogramming checks found no explicit use of
\code{axiom}, \code{sorry}, \code{admit}, \code{native\_decide},
\code{sorryAx}, or the inspected operations for bypassing kernel checking.
The five \code{run\_cmd} registrations in \path{BuildAid/D1.lean} through
\code{D5.lean} add names to a library-suggestion denylist; they do not assert mathematical theorems.

Custom tactics deserve a more specific distinction. The source includes tactics that assign terms directly to goals and code that copies theorem and definition declarations with substituted constant names. Direct goal assignment can omit an elaborator-level compatibility check, but the resulting declaration still undergoes ordinary kernel checking. The inspected declaration-copying code uses the checked declaration interface, and no override disabling the kernel check was identified. These observations do not certify every generated term. \src{S-meta}

The pinned Lean version is 4.33.1. Its own release notes include the kernel fixes also discussed in the later 4.34.0 notes. It would therefore be incorrect to infer that the pinned version lacks those fixes merely from the later release number. No version-specific kernel exploit was established by this audit. \src{lean-release}

\subsection{A negative control for the custom assignment tactic}
The file implementing the custom \code{kcopy} tactic was compiled with the pinned compiler. A separate audit file, outside the submitted source, then deliberately attempted
\begin{verbatim}
import Formalism.LocalToGlobal.NumberFieldFamilies.OwnStruct
theorem audit_invalid_assignment : False := by
  kcopy True.intro
\end{verbatim}
\begin{minipage}{\linewidth}
The compiler rejected it with exit status $1$ and the kernel diagnostic
\begin{verbatim}
declaration type mismatch, 'audit_invalid_assignment' has type
  True
but it is expected to have type
  False
\end{verbatim}
\end{minipage}

This is the expected rejection of an invalid audit probe. It is not an error in an archive theorem. It confirms that this tactic does not, merely by assigning a term, turn every proposition into a theorem.

\section{Further verification of the complete exported proof}
\label{sec:complete-export-verification}

\subsection{What is now established, and by whom}

The formal evidence is substantially stronger than a successful compilation of
selected files.  Andrew Sutherland's public verification page reports that all
175\,187 submitted modules compiled, that all were replayed through the Lean
kernel without rejection, and that the independent Rust implementation
\emph{nanoda} accepted the complete exported proof.  The corresponding public
log records exit status zero.  These are results reported by an independent
verifier; a complete kernel check was not run in this workspace.  The page was
retrieved on 20 November 2026 and still bore its update time of 19:45 UTC.
The primary records are
\url{https://math.mit.edu/~drew/grh/} and
\url{https://math.mit.edu/~drew/grh/nanoda.log}.
The timing log names a configuration file and records exit status zero,
but does not print that configuration's contents or the input file's hash.
Its association with the hash-identified proof export therefore relies on
the verifier's accompanying primary report. We did not authenticate the
remote checker executable or independently reproduce that execution.

Our additional work examined the entire published export, rather than a
selection of source files.  Its SHA-256 digest agrees with the published digest,
and every axiom and declaration record in it has been inspected by the
reproducible inventory program supplied with this audit.  This program is not
a type checker: it inventories the serialized data and verifies its byte
identity.  The independent kernel result and our inventory answer different
questions and should not be conflated.

\subsection{Identity and complete axiom inventory}

The downloaded compressed proof contains 7\,699\,933\,588 bytes.  Its digest is
\begin{quote}\small\ttfamily
\path{61a8938ad5f5e57397c1c0e3b954916af546d9004ade01c3555744a681b05381}.
\end{quote}
The decompressed stream contains 61\,570\,595\,788 bytes, with digest
\begin{quote}\small\ttfamily
\path{5ea74b912147c5c93cc2b6d70cc727cf908df2f62ac6a20480a3fa02f6be244e}.
\end{quote}
Both digests were verified locally.  The decompressed stream was processed
without retaining a second 61-gigabyte copy.  The complete inventory gives
\begin{center}
\begin{tabular}{lr}
Records in the stream & 1\,080\,769\,508\\
Expression nodes & 1\,065\,214\,335\\
Declarations & 1\,484\,733\\
Theorem declarations & 1\,246\,110\\
Definition declarations & 214\,241\\
Opaque declarations & 91\\
Inductive types & 7\,970\\
Constructors & 8\,342\\
Recursors & 7\,972\\
Quotient declarations & 4\\
Axiom declarations & 3
\end{tabular}
\end{center}
The three axiom names are exactly
\[
 \texttt{propext},\qquad \texttt{Classical.choice},\qquad
 \texttt{Quot.sound}.
\]
No \texttt{sorryAx}, \texttt{Lean.ofReduceBool},
\texttt{Lean.ofReduceNat}, or \texttt{Lean.trustCompiler} occurs as an axiom
in this export.  No declaration is marked unsafe or partial.  There are no
repeated declaration names or unknown top-level record forms.  These are
findings about every record in this particular hash-identified export.
They do not assert that a lexical search, by itself, establishes soundness.

The axiom inventory of the entire imported interactive environment is a
different object.  Sutherland's \texttt{audit.log} lists fifteen core axioms,
including compiler-support declarations and \texttt{sorryAx}; their mere
presence in the environment is not a dependency of a target theorem.  The
same log lists only the above three axioms as dependencies of each of the five
final theorems.  The local export inventory independently confirms that the
exported proof contains only those three axiom declarations.  The audit source
and output are available at
\url{https://math.mit.edu/~drew/grh/Audit.lean} and
\url{https://math.mit.edu/~drew/grh/audit.log}.

The canonical challenge was also downloaded and scanned completely.  Its
uncompressed digest is
\begin{quote}\small\ttfamily
\path{bd484ac9886a96f26b1d9a7a59ebd45b7a70d102c015d4f31f9302bbdd6ab2b9}.
\end{quote}
Its four axioms are the three standard ones and \texttt{sorryAx}, as expected
for a file supplying the five statements with placeholder proofs.  Its exact
local byte count is 350\,862\,311; the public page's stated count
350\,862\,970 differs by 659 bytes.  The canonical digest agrees exactly.
This is a discrepancy in reported file size, not a counterexample to a
mathematical statement or evidence of a different challenge.

\subsection{The statements and their definitions}

The five final declaration types refer respectively to the intended Riemann
hypothesis, its explicit version for the zeta function, nonvanishing of all
Dirichlet $L$-functions to the right of $\Re s=1/2$, and the critical-line
statements for all and for primitive Dirichlet characters.  The source
statements quantify over every positive modulus.  For the principal character,
the nonvanishing assertion excludes its pole at $s=1$.  The critical-line
assertions concern $0<\Re s<1$; hence the additional zeros on $\Re s=0$ of
imprimitive Euler factors are irrelevant.  The zeta statement excludes the
negative even integers and the exceptional assigned value at its pole.
Here ``GRH'' means the hypothesis for Dirichlet $L$-functions; the statements
do not quantify over arbitrary automorphic $L$-functions.

A check of names alone would not establish this interpretation: a definition
could be changed while retaining its name.  We therefore carried out an
independent structural comparison against the canonical challenge, starting
from the types of all five target declarations and following every constant
reference.  The comparison succeeded for 53\,009 reached declarations and
6\,075\,060 pairs of expression nodes.  It compared the bodies of definitions
and opaque constants, inductive declarations and their constructors,
recursor rules, universe levels, and every reached non-target theorem body.
All 53\,721 challenge declaration names other than the placeholder axiom
occur in the solution.

This comparison used the actual mathematical expressions, not their printed
names or a hash of their names.  String and numerical components of names
were kept distinct.  Bound-variable indices, universe arguments, literal
values and projection indices were preserved.  Only bound-variable names,
expression annotations and reducibility hints were disregarded; these do
not alter the underlying expression.  The five target proof bodies were
excluded, since those are precisely the bodies the challenge leaves to be
supplied.  Thus the successful comparison establishes the identity of the
statements and their transitive definitions in these two exports, while the
validity of the supplied proofs remains the separate kernel question.
The parser and comparison program are independently written audit tools,
not formally verified software. Their reported equality is reproducible
computational evidence under the explicitly stated representation and
normalization, rather than an additional kernel certificate.

Seven comparison controls passed.  In particular, replacing a definition
body while preserving its name, changing a target type, or changing a
transitively used inductive declaration was detected.  Renaming a bound
variable was accepted.  A control that changed only the target proof was
also accepted, intentionally demonstrating that this program compares
statements and does not check their proofs.  Three additional controls
checked that the optimized and ordinary expression parsers preserve the
order of function and argument in an application.

\subsection{The remaining distinctions in the verification record}

The public verification page states that its build and export ran without a
sandbox.  It reports that a sample of 77 modules agrees byte for byte with
ordinary Lake output, and that the official Mathlib and dependency artifacts
were unchanged.  Its end-to-end canonical Comparator run and its Lean replay
of the full export were still described as running when retrieved.  The
separate Deicyde verification notes likewise still described their own
sandboxed build as unfinished.  None of these unfinished checks is counted
as a successful run here.  The relevant primary records are the public page
above and
\url{https://github.com/Deicyde/riemann-hypothesis/blob/main/verification/NOTES.md}.

The local build has produced 31 submitted modules.  All 31 source files were
compared directly with their original archive members and are identical.
Comparison of their compiled files with Sutherland's published manifest gives
25 identical hashes and six differing hashes.  A different compiled-file
hash need not mean a different proposition or an invalid proof, but it
precludes claiming complete local reproduction of those artifacts.  The
175\,187-entry manifest itself was downloaded and its published hash checked.
The earlier direct-body axiom traversal of 527 declarations reached 79\,855
constants, found only the standard three axioms, and found no missing
constants.  Its scope remains the selected compiled environment.

The two independent checker implementations matter because the pinned Lean
release predates several runtime repairs.  In particular, the official
4.34.1 release notes discuss deliberately constructed reference-counting
inputs that could compromise the Lean runtime.  An implementation in Rust
addresses a distinct implementation risk.  Conversely, the fact that a
checker cannot process this development is not itself a logical rejection:
the reported failures of \emph{con-leche} and \emph{con-ron} concern an
unsupported mutual reflexive inductive declaration, while the bundled
\emph{lean4lean} failure concerns an incompatible core-library version.
See \url{https://lean-lang.org/doc/reference/latest/releases/v4.34.1/}
and \url{https://lean-lang.org/doc/reference/latest/External-Kernels/}.

The more recent issue \#16489 concerns acceptance and rewriting of an
ill-scoped constructor type.  Its primary report expressly records no derived
inconsistency; it should not be described as an established proof of
\texttt{False}.  The inspected source metaprograms construct neither raw
inductive declarations nor forged compiled files.  The primary issue and fix
are \url{https://github.com/leanprover/lean4/issues/16489} and
\url{https://github.com/leanprover/lean4/pull/16512}.

\subsection{Reproducibility and limits}

The evidence bundle contains the inventory source, its dependency provenance,
positive and negative controls, all resulting JSON summaries, and the primary
verification logs.  The controls test forbidden axioms, unsafe declarations,
changed record order, missing final newline, hidden declarations, unknown
record types, and repeated names.  Ten inventory controls passed.  The
optimized application-record parser was also checked against the original
parser on the entire canonical challenge, producing the same inventory apart
from elapsed time.  The full solution inventory completed in approximately
346 seconds, with a largest recorded child-process resident set of approximately 1.1 gigabytes.

The correct conclusion from these records is that there is substantial
independent formal evidence for the five stated theorems.  The complete export
contains no additional axiom of the kinds specifically sought in this audit,
and its target statements and transitive definitions agree with the canonical
challenge under the explicit structural comparison above.  The records do not constitute a complete human mathematical reading of
175\,187 source files, nor an independently executed full kernel replay in this
workspace.  No fatal mathematical flaw has been established by the checks
reported here.


\section{Review of the quasi-Riemann and Landau--Siegel arguments}
\label{sec:related-literature}

The nine additional references were examined for two purposes: to test
their consequential assertions and to determine which parts of the
announced proof they could simplify. Their conclusions must be kept
separate. A fixed zero-free half-plane for all Dirichlet characters,
an effective lower bound on the distance of a real zero from $1$, and
the critical-line assertion are different statements.

\subsection{The quasi-Riemann papers and catalogue}

The September paper has 199 pages; the October paper has 48. The
versions examined are fixed by the repository commit
\path{4eaee87548250b2682affda7d05d947c05554bae}.
The September theorem asserts nonvanishing in $\Re s>7/8$ for
the specified finite-order Hecke characters over
$\Q(\sqrt{-3})$. The October argument yields $\Re s>11/12$;
its shorter proof does not improve $7/8$. Factorization of the
norm-lifted Hecke $L$-function transfers these conclusions to ordinary
Dirichlet characters, with the possible cancellation at $s=1$
treated separately. The companion manuscript proves this transfer
and the subsequent analytic deductions.\src{qrh-september,qrh-october}

For the October paper, the review covered the amplification, the
parameter ranges in the iteration, the Poisson transformation and
its arithmetic cancellation, and the stated cusp inputs. For the
September paper it covered the main theorem, the passage from
smoothed sums to nonvanishing, the marked-moment statements,
amplification, and the restrictions on the outer indices.
The full metaplectic transformation calculations and every step of
the long marked-moment induction have not been independently
reconstructed. These papers are therefore not promoted to fully
verified inputs of an unconditional proof in this report.

The catalogue was examined at commit
\path{5b7217781269248f39fbcec2432394d3d1610344}; its maintained
collection contains 4\,310 records. Selected entries concerning
Mellin transforms, inverse moments, amplification, and the alleged
improved exponent were read with their evidence qualifications.
An entry being marked as proved relative to another source does not
verify that source. In particular, the smaller advertised distance
from the critical line remains dependent on a stronger weighted
estimate not established in the examined reconstruction.
The companion manuscript independently proves the useful
finite-polynomial obstruction and the absence of common zeros of
the indicated family of oscillatory Mellin transforms.\src{qrh-catalogue}

One potentially useful improvement can be stated precisely. Suppose
the fixed-character smoothed M\"obius sums satisfy
$A_1(D)\ll_\epsilon D^{\beta+\epsilon}$, where $\beta>1/2$.
Exact deletion of one prime, retaining this cancellation, would
improve the exponent to
\[
 T_k(\beta)=\frac{k\beta}{k-1+2\beta}
\]
if the required complete mean square could be proved over a range
of length $D^{h_*(\beta)}$, where
\[
 h_*(\beta)=\frac{k(2\beta-1)}{k-1+2\beta}.
\]
For sixth-power characters the first step from $7/8$ would use
$h_*=2/3$ and yield $7/9$. The proof of this conditional implication
is given in full in the companion manuscript.

The needed shorter mean square is \emph{not} supplied by the cited
results. The October iteration needs the original outer range to
exceed the coefficient length by a fixed power. Below that length,
its transformed range violates the next proposition's hypothesis,
and the unit divisor retains the original scale. Later catalogue
estimates either retain a weaker exponent or sum only over
sixth-power-free indices, which exclude the sixth powers used in
this amplification. A direct use of marked prime factors also
fails its required range inequalities. These checks identify an
unproved new estimate; they are not an iteration to RH.

There is a concrete obstruction to replacing M\"obius coefficients
by arbitrary bounded coefficients. Put coefficient $1$ on prime
ideals and $0$ elsewhere, and use a fixed nonnegative smooth weight
supported on norms comparable to $D$. At an argument $p^k$ with
$Np\asymp H^{1/k}=D^{h/k}$, $h/k<1$, every nonzero column
character is $1$. The prime-ideal theorem in the fixed field gives
an inner sum $\gg D/\log D$ and
$\gg H^{1/k}/\log D$ distinct such arguments. The complete
mean square is consequently
\[
 \gg \frac{D^2H^{1/k}}{(\log D)^3}.
\]
This contradicts $D^{1+\epsilon}H$ when
$h<k/(k-1)$ and $\epsilon$ is sufficiently small. The argument
does not contradict an estimate for the special M\"obius
coefficients, or an outer set which excludes $k$th powers.

\subsection{The six Landau--Siegel references}

The mathematical content of all six references in this group was
read. The original OpenAI paper supplies the finite determinant
method; Guo supplies an explicit constant and a uniform paired-tail
estimate; Tao explains the mechanism; Pescador's version 3 gives
the weighted interpolation and a sharper inert-prime estimate;
van Dijk gives a particularly short rectangular determinant proof;
Kerr gives an alternative earlier interpolation argument.
The complete six-page text of van Dijk's paper was obtained through
the PDF reader when direct download returned an access error.
No access restriction was bypassed.\src{sz-openai,sz-guo,sz-tao,sz-pescador,sz-vandijk,sz-kerr}

The companion manuscript reconstructs the determinant argument
directly over a quadratic order. Its finite-field interpolation
proof is a root-count argument and does not require a polynomial
irreducibility assumption. It keeps ordinary derivatives, proves
the required divisibility by each inert prime, and bounds the
determinant using row-normalized Hadamard. The resulting finite
bound has $N\log N$ in place of $2\log(N!)$ in the corresponding
size term. It leads to
\[
 \sum_{\substack{p\le X\\(d/p)=-1}}\frac{\log p}{p}
 \le \frac12\log X+\frac14\log|d|+\log\log(3X)+O(1)
\]
and to the effective exclusion of a real zero too close to $1$.
These improvements concern the determinant estimate; they do not
change a real-zero exclusion into GRH.

Five exact quadratic-ring examples, of matrix dimensions
$6,10,15,21,14$, were checked with fraction-free determinants.
The calculations verify the row changes, determinant preservation,
nonvanishing modulo the auxiliary prime, and the stated norm
divisibility. There are useful negative controls: divided
derivatives can destroy divisibility, since the divided third
derivative of $z^3-w$ is $1$ whereas its ordinary third derivative
is $6$; and allowing the interpolation order to exceed the
characteristic admits the nonzero polynomial $w^p-w$ in the kernel.
Neither change is made in the actual rectangular determinant proof.

Guo's paired-tail argument was generalized to any bounded sequence
with bounded interval sums and proved for every $0<\beta<1$.
The kernel maximum is computed exactly. Its error stays uniform
as $\beta$ tends to $1$; an unaccounted factor
$(1-\beta)^{-1}$ was specifically sought and does not occur.
The explicit rational losses in Guo's and van Dijk's conclusions
were independently checked with exact arithmetic and certified
logarithm bounds. The class-number-one residue sieve in van Dijk's
appendix was also reproduced: the $5\,641\,415$ remaining large
candidates give precisely the listed final possibilities. This
sieve is not an input to the GRH reconstruction.

Tao's article is explicitly a sketch: its temporary irreducibility
assumption is not a full proof of injectivity in every case. The
finite-field proof above avoids that assumption. Kerr uses an
external zero estimate of Masser; the primary bibliographic record
was located, but its full text was not independently checked here.
The replacement determinant proof does not use it. These are
scope distinctions, not identified contradictions in the claimed
zero-exclusion theorem.

\subsection{A repairable exceptional conductor in Kerr's field setup}

One actual local defect was found in an auxiliary reference. At
the positive fundamental discriminant $q=8$, the field
$\Q(\sqrt q,\sqrt2)$ equals $\Q(\sqrt2)$ and has degree two,
not four. Thus a four-element basis and two independent sign
automorphisms cannot be used in that case. The exceptional case
is easily removed, as the following complete argument shows.
It is not a fatal defect in the zero-exclusion theorem, and it
is not a flaw identified in Alp\"oge's proof.\src{sz-kerr}

\begin{proposition}
For the primitive real character $\chi_8$ of conductor $8$,
$L(s,\chi_8)>0$ for every real $s>0$.
\end{proposition}
\begin{proof}
The character takes the values $1,-1,-1,1$ on residues
$1,3,5,7$ modulo $8$ and vanishes on even integers. Its
partial sums are bounded. Summation by parts therefore makes
its Dirichlet series converge for every real $s>0$.
Let $f(t)=t^{-s}$. Then $f''(t)=s(s+1)t^{-s-2}>0$.
For $x>0$,
\[
 f(x)-f(x+2)-f(x+4)+f(x+6)
 =\int_0^2\int_0^4 f''(x+u+v)\,dv\,du>0.
\]
Set $x=8n+1$ and sum over $n\ge0$. Grouping consecutive
blocks of eight terms preserves the value of the convergent
Dirichlet series. Every block is strictly positive, so their
sum is strictly positive.
\end{proof}

\subsection{What can be removed from the proposed GRH construction}

Under a uniform Dirichlet quasi-Riemann theorem, a buffered
logarithmic-derivative estimate and Mellin inversion give ordinary
prime supply in reduced arithmetic progressions at a polynomial
scale. The companion manuscript proves this with all conductor
dependence stated. It then supplies the simultaneous ordinary
prime congruences and the relative sizes required of the two
auxiliary quadratic-field parameters. A uniform real-zero gap
also gives $|L'/L(1,\chi)|\ll\log q$ for real primitive
characters, removing an exceptional-character choice when that
derivative bound is its only purpose.

These deductions have strict limits. The source also needs
prime supply for ray-class characters over varying quadratic
fields; a theorem only about Dirichlet characters does not
establish it. The source's scalar $A$ bounded by $2xb$ is a
height quantity and cannot be replaced by a M\"obius estimate.
The specially chosen quadratic smoothed sums must meet an
$x^{1/2+\epsilon}$ bound, for which $7/8$ is insufficient.
The integral theta-lift, covolume, and uniform cohomological
estimates consequently remain substantial mathematical work
for a full conventional reconstruction.

One part of that reconstruction is now simpler. For integral maps
between Hermitian lattices, their slope filtrations and a trace-square
estimate for the sum of the matrices $A_iA_i^*$ give a direct
comparison of the target mean slope with the weighted source mean.
The companion manuscript proves the filtration properties, the matrix
dimension bound, and the resulting integral estimate. Its elementary
form replaces a coefficient $4$ by $2$ in the logarithmic error and
removes a dependence on the number of maps. The hypotheses still
require the actual integral maps and the quantified near-orthogonality
estimates; a nonzero determinant cannot replace them. Normalized
columns of the unimodular matrix
$\left(\begin{smallmatrix}1&M\\0&1\end{smallmatrix}\right)$
give an explicit counterexample to that proposed replacement.


\section{Further review of the uniform arithmetic estimates}
\label{sec:audit-uniform-continuation}

This continuation concentrates on the estimates\cite{alpo-uniform-source} whose constants must
remain controlled as the auxiliary primes, weights, and number of theta
lifts grow. No counterexample satisfying the actual hypotheses of a
consequential submitted theorem was found in these checks. The companion
manuscript now contains ordinary proofs of several additional groups of
results: simultaneous selection of the quadratic twists under the
explicit QRH hypothesis; the cyclotomic common-interval argument
and its coarse cohomological input, using stated classical
$K$-theory results; the determinant and lattice
argument for the divisor inequality; and local harmonic-analysis
ingredients of the theta estimates. None of these statements promotes
the unfinished global construction to a proved theorem.

\subsection{The full family of character twists}

There is an important correction to the scope of the earlier shorthand.
The three characters $\chi_q,\chi_d,\chi_q\chi_d$ suffice for the RH
specialization. In the Dirichlet-GRH argument the chosen primes must
also work for a fixed finite family of base characters $\xi$, with the
three nontrivial quadratic twists of every $\xi$. The number of base
characters may depend on the fixed character being studied, but must
not grow with the scale. This is a correction to the description of
our earlier partial reconstruction; the submitted development contains
the larger family.

The new conventional proof establishes the following conditional
selection. Assume that every primitive Dirichlet $L$-function has no
zero in $\Re s>11/12$, with the pole of $\zeta$ allowed. Fix a modulus,
$0<\epsilon<1/52$, and $A\ge1$. For every sufficiently large $Y$
there are primes of the required sizes
\[
 \tfrac12Y^{\epsilon/5}<\log q<Y^{\epsilon/5},\qquad
 \tfrac12Y^{\epsilon/5}(\log Y)^2<\log d
                  <Y^{\epsilon/5}(\log Y)^2,
\]
with $q\equiv3\pmod8$, $d\equiv7\pmod{16}$,
$q\mid d+1$, and every odd prime at most $Y^{\epsilon/5}$
dividing $d+1$, for which all the requisite primitive twists obey
\[
       |G_\chi(y)|\le y^{1/2+\epsilon/2},\qquad
       |L'(1,\chi)/L(1,\chi)|\le4\log y
                     \quad(Y\le y\le Y^A).
\]
The character with no quadratic factor is omitted from this selection;
it is part of the fixed function to be bounded later. Exponent-zero
quadratic factors are omitted, rather than interpreted as principal
characters at a larger modulus.

The proof has three substantive checks. A smooth cutoff gives
\[
 G_\chi(y)=-\sum_n\Lambda(n)\chi(n)\min(1,y/n)w(n/Z)
             +O(yZ^{-1/24}\log(Q+3)),
\]
uniformly in the conductor $Q$; it is used only for nonprincipal
characters. Weighted prime averages in the required progression bound
all nonsquare quadratic correlations, with the factor $\varphi(m)$
from normalization explicitly retained. Finally a derivative estimate
for the finite sum gives a maximal second moment on the whole
interval, avoiding a growing discretization error. For the second
prime, $\log m=O(Y^{\epsilon/5})$, whereas its averaging scale has
logarithm $Y^{\epsilon/5}(\log Y)^2$, so the correlation errors are
smaller than every fixed negative power of $Y$, uniformly in the
first prime.

This is a useful integration of the supplied QRH work: a fixed
zero-free half-plane can support a \emph{selection by averaging}, even
though it gives no square-root bound for every individual character.
The proof and all conductor estimates appear in the companion
manuscript. Its use of the QRH conclusion remains conditional, since
the two original QRH papers have not been completely reconstructed
here. It supplies neither the varying quadratic ray-class estimate
nor the theta-map construction.

The compatibility with the later cohomological interval has also
been checked. Fixing the cohomological constant $J$ first, one
selects the primes at $U=X^{1/J}$ with character estimates on
$[U,X]$. Cohomology then chooses its subinterval inside this
already covered interval. The stabilized degree has the form
$[E_h:\Q]=4c2^h$, where $c$ depends only on the fixed character.
The substitution $x=c^{1/(2+8\epsilon)}x'$ changes the allowable
weight range from $(x')^{401}$ to $x^{400}$ and leaves an interval
of exponent $3/2$. The logarithms of $q,d$ retain their original
scale $U$ throughout. No circular selection of a conductor or
of the prime pair was found in this part of the assembly.

\begin{example}[Too many forced splitting conditions]
For $t\ge2$ let $m=16\prod_{3\le\ell\le t}\ell$, with the product
over primes, and require $p\equiv7\pmod{16}$ and
$p\equiv-1\pmod\ell$ for every odd $\ell\le t$.
For every prime $p>t$ in this reduced class, quadratic reciprocity
gives $(\ell/p)=(-p/\ell)=1$ at odd $\ell\le t$, and
$(2/p)=1$ by the supplementary law. Hence $(n/p)=1$ for every
$n\le t$, and
\[
      \sum_{n\le t}\Lambda(n)(n/p)=\psi(t).
\]
The elementary inequality
$\psi(2r)\ge\log\binom{2r}{r}
\ge2r\log2-\log(2r+1)$, followed by monotonicity, gives
$\psi(t)\gg t$ for large $t$. Thus the second moment over any
nonempty such prime family is $\gg t^2$, not $O(t\log t)$.
The binomial inequality follows because each prime valuation of
$\binom{2r}{r}$ is at most $\lfloor\log_p(2r)\rfloor$;
the lower bound uses that the largest of the $2r+1$ binomial
coefficients is the central one. Under the QRH hypothesis the
progression contains primes at sufficiently large averaging scales,
so the example is nonempty.

The actual construction forces splitting only at primes at most
$Y^{\epsilon/5}\le t^{1/5}$. Their total contribution to the
finite sum is $O(\sqrt t)$ and is retained in the second moment.
The example therefore refutes an unjustified extension, not the
required selection theorem.
\end{example}

\subsection{Cohomological lengths in the cyclotomic towers}

The group used in the source has coefficients
$\Q_p/\Z_p(-e)$, with $e\ge1$. It is unramified away from $p$
and has zero localization on the full decomposition groups above
$p$. Over $K_\infty=K(\mu_{p^\infty})$, let $X$ be the compact
dual of the group of unramified $p$-primary characters that vanish
on those decomposition groups. If $\gamma$ is a topological
generator of $\operatorname{Gal}(K_\infty/K)$, $T$ its action on
$X$, and $u=\varepsilon_p(\gamma)$, the exact descent formula is
\[
 \lambda_p(K_m,e)=
 \length_{\Z_p}X/\bigl((u^eT)^{p^m}-1\bigr)X.
\]
The companion manuscript proves this formula with the local
conditions. The divisible coefficient module makes the global and
local inflation kernels zero; the negative twist is nonzero, so
the scalar defining those kernels is nonzero.

For a fixed abelian field $M$, and the actual primes $q,d$, the
module for $M(i,\sqrt{-q},\sqrt{-d})$ has a generator bound
\[
                         R\le4(g+6[M:\Q]),
\]
where $g$ is a fixed Ferrero--Washington bound for $M(i)$.
The proof uses at most $2[M:\Q]$ primes above $q$ and at most
$4[M:\Q]$ above $d$ in the infinite cyclotomic tower, one tame
inertia generator at each, and two index-two character bounds.
The latter bounds hold in characteristic $2$ because the
involution satisfies $(T-1)^2=0$, with the restriction obstruction
accounted for. No dependence on the sizes of $q,d$ is discarded.

For a module generated by $R\ge1$ elements, put
$j=1+\min\{k:p^k\ge R\}$ and $c=v_p(u-1)$.
There are $s_1,\ldots,s_R\in\Z_p$ for which
\[
 \lambda_p(K_m,e)\le
      \sum_{i=1}^R(c+\max(m,j)+v_p(e-s_i)).
\]
This estimate is proved by lifting the endomorphism of $X$ to a
free cover, bounding its cokernel by a determinant, and choosing
nearest points in the compact set
$\{u^{-p^j s}:s\in\Z_p\}$ for its eigenvalues. Thus finite
torsion in $X$ is included. For $1\le e\le Z$, at most $R$
integers have $v_p(e-s_i)>\lfloor\log_p Z\rfloor$ for some $i$.
The exceptional sets for all the required coefficient primes are
united before selecting an interval.

The resulting elementary interval lemma improves the source's
constant: for a union of at most $R_*$ exceptional weights and
$C_t\ge\eta>0$, one can take
$J=2^{R_*+1}(C_t/\eta)^{R_*}$ and find
$X^{1/J}\le Y$, $Y^2\le X$, such that
$e\le x^{C_t}$ implies $e\le x^\eta$ for every exceptional
$e$ and every $x\in[Y,Y^2]$. The proof removes one obstruction
at a time on the logarithmic scale; it does not intersect unrelated
intervals. Full proofs are in the manuscript.

The remaining input at the exceptional weights is explicitly
identified there: an absolute $C_K$ must bound
\[
 \log\#H^2(G_{F,S},\Z_p(e+1))
       \le\exp\bigl((2e[F:\Q])^{C_K}\log(2|D_F|)\bigr).
\]
The source constructs such an estimate through integral homology
torsion bounds and a global class-counting argument; it is not
inferred from the generator bound. The comparison route used in
the final assembly has now been traced explicitly. A separately
defined interface for a comparison through algebraic $K$-theory
is conditional, and is not being counted as a proof of that
comparison.

The manuscript now proves the orbit-and-stabilizer torsion bound
\[
 \#H_m(\Gamma,\Z)_{\mathrm{tors}}
       \le B^A s^{A\sum_{j=1}^m s^j},
\]
for an augmented exact permutation complex with at most $A$
orbits, stabilizers of order at most $s$, and incoming boundary
columns of norm at most $B$. It proves the actual arithmetic
stabilizer bound $s\le3^{nN^2}$ by reduction modulo $3$.
The geometric complex and its orbit bound have now also been
reconstructed. The complex is the ordered nerve of the normal
cones of the Hermitian reduction polyhedron. A finite closed-convex
nerve argument proves exactness, including tuples with repetitions.
The determinant estimate comes from the elementary lattice box
bound and the exact discriminant covolume. Integral matrices of
bounded norm determinant are reduced by the residues of
$I(A)A^{-1}$ modulo $I(A)$, where
$I(A)=|N_{F/\Q}(\det A)|$. This gives the required orbit bound
without assuming that $\OO_F$ is a principal ideal domain.
The subsequent numerical exponent can therefore be $13$ rather
than $24$ in the unconditional arithmetic homology estimate.
The local cohomology calculation is also complete:
its error is at most
$n(2+\lfloor\log_p n\rfloor+v_p(e))\log p$.
Thus a strict Selmer bound of the stated exponential form gives
the required $H^2$ bound.

The intervening class count has now been reduced further. The
manuscript proves the class-number bound, the logarithmic unit
theorem, the unit-and-ideal-class comparison, cyclotomic enlargement,
and the passage from uniformly bounded equivariant class-group
characters to strict Selmer classes. The finite-field homology
functionals and the integral arithmetic cycles have now also been
constructed, so that the original coarse Selmer route is complete
with its stated classical inputs. A separate $K$-theory argument,
given below, supplies a sharper coarse $H^2$ estimate. The estimates for the other
prime and character components remain uncompleted.
The newly proved argument gives the required simultaneous summed
bound $\exp(x^{4\epsilon/5})$ for the $2$-adic contribution and
the fixed exceptional odd-prime layers, uniformly in $q,d$.

\begin{example}[Discarding finite torsion changes the answer]
Take $X=\Z/p^a\Z$, $T=1$, $u=1+p^c$, where $a,c\ge1$ and
$c\ge2$ if $p=2$. For $e\ge1$,
\[
 \length X/((u^eT)^{p^m}-1)X
                  =\min(a,c+m+v_p(e))>0.
\]
The lifting-the-exponent formula proves the equality. Replacing
$X$ by its torsion-free quotient gives zero and loses this entire
length. The source and the new proof both use generator bounds
for the full module, so this is not a counterexample to either.
\end{example}

\subsection{The divisor inequality and its conjugate primes}

The quotient-metric determinant inequality has now been proved
without assuming that the coefficient field is CM. Its proof
pairs an embedding $v$ with
$a\mapsto\overline{v(\bar a)}$ and applies determinant
Cauchy--Schwarz. The exact patched formula has an explicitly
nonnegative difference between the finite Gram-determinant sum
and the quotient covolume. Discarding that difference gives the
required inequality. The manuscript also proves the local dual
containment bound, the lattice-index formula, a weighted
operator-matrix norm bound, and the descended pairing estimate
for chains of rectangular matrices.

The initial polynomial coefficient duality has now been
reconstructed as well. Restriction to lower triangular matrices
is injective by Gaussian elimination. If that restriction is
integral, localization at the principal minors recovers the
whole polynomial with integral coefficients: none of those
minors vanishes identically modulo a prime. The normalized
Fock pairing with the extremal polynomial is evaluation at
the identity. Adjoint divided row operations therefore identify
the exact integral dual of the Weyl lattice with the dual-Weyl
lattice. This yields a full elementary proof, including small
residue characteristics, without importing a tableau theorem.
Tensor products, negative determinant characters, local
transport, and conjugation are treated explicitly. For
unimodular base lattices the coefficient dual containment has
$s=1$.

\begin{example}[Saturating the Weyl lattice destroys the pairing]
For binary quadratics, use $H=x^2$ and normalize the Fock form
by $[H,H]=2$. The monomials $x^2,xy,y^2$ are orthogonal
with squared bilinear values $1,1/2,1$ in that normalization.
Thus the dual of
$\Z x^2\oplus\Z xy\oplus\Z y^2$ is
$\Z x^2\oplus2\Z xy\oplus\Z y^2$.
Replacing the latter by its saturation would give
$B(xy,xy)=1/2$, so the claimed integral pairing would fail.
If the normalization were omitted instead, the dual would be
$\tfrac12\Z x^2\oplus\Z xy\oplus\tfrac12\Z y^2$.
Both changes are excluded by the actual definitions.
\end{example}

\begin{example}[A Hermitian basis at only one prime gives a wrong index]
In $K=\Q(i)$ take $L=\Z[i]$ with $h(x,y)=\bar x y$ and
$\mathfrak p=(2+i)$, $\bar{\mathfrak p}=(2-i)$.
The element $e=2-i$ is a basis of $L_{\mathfrak p}$, but
$h(e,e)=5$ has valuation $1$ at $\mathfrak p$.
The Hermitian dual is $L$ itself, so its signed index is $0$,
not $-1$. A formula equating that index to the negative Gram
valuation therefore fails for a basis chosen at only one prime.
The source uses a basis simultaneously at
$\mathfrak p,\bar{\mathfrak p}$, exactly excluding this example.
\end{example}

\begin{example}[An equality cannot replace the quotient inequality]
Let $\alpha=\sqrt[3]2$, let $\zeta$ be a primitive cube root of
unity, and set $K=\Q(\alpha,\zeta)$ with complex conjugation
fixing the real $\alpha$. Give $K^2$ the identity Hermitian form
and take $P=(1,\alpha)$. Its quotient lattice is $\OO_K$.
At all three complex places the quotient metric is
$(1+\alpha^2)^{-1}$, whereas the algebraic Gram scalar is
$\Gamma=(1+\alpha^2)^{-1}\in K$ with its distinct embeddings.
In the normalization $c=2[K:\Q]^{-1}\log$ of the covolume
relative to $\OO_K$, one obtains
\[
 c=-\log(1+2^{2/3}),\qquad
 \sum_{\mathfrak p}\frac{\log N\mathfrak p}{6}
                  v_{\mathfrak p}(\Gamma)=-\tfrac13\log5.
\]
Indeed $N_{K/\Q}(1+\alpha^2)=25$, because
$\prod_{j=0}^2(1+\alpha^2\zeta^{2j})=5$.
The inequality is strict since $(1+\alpha^2)^3>5$.
The field is totally imaginary but not CM: the fixed field
$\Q(\alpha)$ is not totally real. The actual source needs and
proves the inequality, so strictness is compatible with it.
\end{example}

\begin{example}[Multiplying minors does not bound an index above]
Over a discrete valuation ring, let
$A=(1,1)$ and $B=(1,-1+\pi^N)^t$. Each matrix has its first
determinantal ideal equal to the unit ideal, while $AB=(\pi^N)$
has valuation $N$, arbitrarily large. Cauchy--Binet gives an
inclusion of determinantal ideals in the direction producing a
lower valuation bound, not an upper one. The actual local-pairing
estimates must therefore be verified for the actual products.
This remains an identified task, not a claim that the source
reverses the inclusion.
\end{example}

\subsection{Local theta estimates and the remaining global sums}

The local phase
$f_D(x)=\sum_{i=1}^m x_i+D/(\prod_{i=1}^m x_i)$ has critical
points $x_i=z$, $z^{m+1}=D$, with Hessian
$z^{-1}(I_m+\mathbf1\mathbf1^t)$ and determinant
$(m+1)z^{-m}$. The new proof gives the exact stationary-phase
formula with self-dual Haar measure, including residue
characteristics dividing $2(m+1)$. The Hessian is still
invertible over the characteristic-zero local field. Its inverse
and the factor $|m+1|^{-1/2}$ are retained. Uniformity holds for
fixed local data and compact parameter sets; it is not asserted
uniformly over growing dimensions or residue characteristics.

At a distinguished split prime the local order is
\[
 R_\pi=\{(x,y)\in\OO_F^2:x-y\in\pi\OO_F\},\qquad
 \psi_R(x,y)=\psi_F((x-y)/\pi).
\]
It is self-dual. Writing $(x,y)=(y+\pi z,y)$ shows that a
pair $(a,b)$ annihilates it exactly when $a\in\OO_F$ and
$(a-b)/\pi\in\OO_F$. The self-dual Haar measure is
$q=\#(\OO_F/\pi)$ times product Haar measure normalized on
$\OO_F^2$, and gives $R_\pi$ volume one. This factor is
present in the source.

\begin{example}[The Gaussian admissibility and measure hypotheses]
For $m=3$ over $\Q_2$, consider
$Q(x)=\sum_i x_i^2+\sum_{i<j}x_ix_j$ and the character
$\exp(2\pi iQ(x)/2)$ on $\Z_2^3$.
Grouping the eight residue vectors by their number of odd
coordinates gives the sum $1-3-3+1=-4$.
The squared absolute value of the normalized integral is
$1/4$. The determinant formula without its admissibility
condition would give $1/2$. The dual quadratic form is not
trivial on the annihilator lattice, so this does not satisfy
the actual Gaussian hypothesis.

For the split order above with $F=\Q_p$, $p$ odd, the integral
of $\psi_R(x^2,y^2)$ over $\Z_p^2$ is $1/p$ with product
Haar measure and $1$ with self-dual Haar measure. To verify
this, reduce modulo $p$: the sum is
$|\sum_x\exp(2\pi i x^2/p)|^2=p$; the latter equality
follows by the invertible substitution $(x,y)\mapsto(x-y,x+y)$
and additive-character orthogonality. Division by $p^2$ and
then multiplication by the measure factor $p$ give the two
values. Thus omitting the measure factor changes the answer.
\end{example}

The manuscript also proves positivity of the finite-cokernel
kernel by an explicit Gram identity, determines the nonzero
reproducing scalar by an idempotent convolution operator, and
gives the Poincar\'e-series unfolding with its absolute
integrability hypotheses. From the displayed orbital-sum inputs
the final off-diagonal error can be improved from $b^{-6}$ to
$14b^{-7}$ for $b\ge16$. These are completed local and
functional-analytic deductions. The uniform estimates for the
actual global orbital sums and the arithmetic recursion after Poisson summation
remain to be reconstructed. The local formulas alone do not
settle them.

The actual local orbital kernel has now been reviewed more deeply.
Block Gaussian elimination proves finiteness of its orbital sum,
uniformly on compact parameter sets. A small perturbation preserves
the image lattice and its inverse-coordinate functionals, proving
local constancy uniformly against every comparison configuration.
A finite-module character argument proves the stronger support
condition $t^{-1}-t'^{-1}\in R$: isomorphisms failing a successive
inverse-coordinate relation cancel under a column shear with
nontrivial character. This proof also works for the nonmaximal
split order.

The sum over roots is then converted exactly to a Haar integral.
All mixed critical terms cancel, and the normalizing factors in
the truncated orbital identity cancel exactly. The previously
outstanding germ formula is now proved as well. Reversing rows
turns the two-block orbit into an upper Hessenberg matrix.
Its determinant fixes the first diagonal entry, while the change
to first-row coordinates has Jacobian equal to the modulus of
the trailing determinant. The first-order determinant expansion,
together with local constancy, gives the exact oscillatory integral.
Stationary phase then gives the germ with every power of $|t|$,
$|z|$, and $|m+1|$ retained. A finite block-matrix construction
removes any fixed initial blocks and gives uniform support and
local constancy independently of the comparison module.

The actual single-block Parseval identity is proved by a finite
cokernel calculation. For a functional $\alpha:C\to K/R$, the
module $\{(x,k):[k]=-\alpha(x)\}$ identifies the bordered
cokernel and preserves its marked vector and inverse-coordinate
functional. Character orthogonality gives exactly the product
of the two original cokernel orders, with no omitted factor.
The unrestricted diagonal is consequently bounded by the previous
majorant's total mass. Finite Fourier averaging and a limit in
the fixed finite quotient $L_c/R$ complete the local majorant
induction. Thus neither the germ nor that local induction is
an outstanding hypothesis in this version. The uniform global
orbital estimates remain a separate task.

\begin{example}[Two normalization and integrality requirements]
Over $\Q_2$, take the initial one-by-one matrix $(2)$ on both
sides. Its marked cokernel is $\Z/2\Z$, with generator $1$
and functional $x\mapsto[x/2]$. The single-block kernel has
values $3$ and $1$ on the additive classes $0$ and $1/2$,
respectively: the three pairs $(r,b)\in\{0,1\}^2$ with
$rb=0$ contribute to the first class, and the remaining pair
to the second, with phase one. The total is $4$, whereas the
initial kernel is $1$. Omitting the product of the two
cokernel orders would give a false Parseval identity.

Over $\Q_3$, the initial scalar $3$ and last scalar $-1/3$
give full orbital matrices
\[
 \begin{pmatrix}3&3y\\3x&3xy-1/3\end{pmatrix}.
\]
For an integral matrix, $x=u/3$, $y=v/3$ modulo $\Z_3$,
where $uv=1$ in $\mathbb F_3$. The determinant is $-1$,
so its cokernel is zero. Against the zero marked module the
two phases sum to $\zeta_3^2+\zeta_3^4=-1$.
The last scalar is nonintegral. Restricting the reduced
last-block sum to integral matrices would incorrectly give
zero. The actual source keeps all last-block orbital cosets.
\end{example}

\begin{example}[The functional requires the stronger perturbation scale]
For $R=\Z_p$, the matrices $(p^2)$ and $(p^2+p^3)$ have
the same image ideal. Their inverse entries differ by
$-1/(p(1+p))\notin R$. Their marked generators therefore
have different inverse-coordinate values, and no isomorphism
preserving the generator preserves the functional. Preserving
only the cokernel lattice would be insufficient. The actual
perturbation proof uses the square of the determinant in
its smallness condition and preserves both specifications.
\end{example}

\begin{example}[Omitting the support condition in a mixed term]
Take $K=\Q_3$, $R=\Z_3$, $m=1$, and a character $\psi$
trivial on $\Z_3$ with $\psi(1/3)\ne1$. Put
\[
 \Omega=(3+9\Z_3)\cup(-3+9\Z_3),\qquad
 g=1_{3+9\Z_3},\qquad t=3^N\quad(N\ge4).
\]
Consider the mixed root sum with phase $\psi(4z/t)$,
where $z^2=-tc$ and $c$ ranges through representatives of
the additive cosets in $\{-z^2/t:z\in\Omega\}$.
The set is compact and $\Z_3$-stable for these $t$, as follows
by taking the near-one square root after adding an element
of $\Z_3$ to $c$. Hence this is a finite sum of coset
contributions.

Choose the representative $c=-9/t$ in its coset. Its two roots
are $3,-3$, and only $3$ contributes. Change only that
representative to $c+1$. The contributing root is now
\[
 z'=3\sqrt{1-3^{N-2}}\in3+9\Z_3,\qquad
 \frac{4(z'-3)}t=-\frac23+O(3^{N-3}).
\]
The expansion follows from
$\sqrt{1+u}=1+u/2+u^2b(u)$, with $b(u)\in\Z_3$
for $u\in3\Z_3$. Thus that one contribution is multiplied
by $\psi(-2/3)\ne1$, and all other contributions are
unchanged. The two total sums are different, so they
cannot both vanish. This works for arbitrarily small $t$.

The omitted support condition would require $2/z\in R$
where $g(z)\ne0$, and it fails at $z=3$. The actual kernel
has precisely this support condition. The example therefore
refutes a strengthened vanishing assertion, not the submitted
identity.
\end{example}

\subsection{The strict Selmer comparison and its remaining input}

The next steps of the source's class-counting argument have been
reconstructed with their precise objects. Let $K/L$ be the relevant
cyclic cyclotomic extension of degree $d=p^{\nu-c}$, with generator
$\gamma$ acting on a primitive $p^\nu$-th root by exponent $1+p^c$.
Put $a=(1+p^c)^E$, with $E\ge1$, and let $I$ be the fractional
divisor group. Write $K^\times$ additively in this paragraph,
let $\partial$ be its divisor map, and set
\[
 g=\gamma^{-1},\quad \delta=g-a,\quad
 V=\partial^{-1}(\delta I),\quad
 W=V/(\delta K^\times+\mu(K)).
\]
The manuscript proves the exact sequence from the unit quotient
$(\mathcal O_K^\times/\mu(K))/\delta$ to $W$ and then to the
kernel of $\delta$ on the ideal class group. It proves
\[
 \#W=\frac{(a^d-1)^{r_2(L)}}{a-1}
          \#\bigl(\operatorname{Cl}(K)/\delta\operatorname{Cl}(K)\bigr).
\]
The unit factor is the determinant on the sum-zero hyperplane of
the permutation representation on complex places. A complete
lattice-volume proof of the unit theorem is included. The factor
$a-1$ in the denominator is essential: for $K=L=\Q(i)$ the
unit quotient is zero and has order one; the full permutation
determinant would incorrectly give $a-1$.

The proof also gives an injection of $W[p]$ into invariant
Kummer classes modulo roots of unity. Cyclic Hilbert theorem 90
and a classification by the root-of-unity exponent and the norm
value bound those classes by
$p^{2n+\lfloor\log_2h_L\rfloor+2}$, where $n=[L:\Q]$.
All of this is proved in ordinary mathematics. The polynomial
\[
             N_a=\sum_{j=0}^{d-1}a^j g^{d-j}
       \quad\text{satisfies}\quad
             N_a\delta=(1-a^d)g.
\]
It is the weighted product in the source, not the ordinary norm.
This identity gives the exact kernel exponent used in the finite
cardinality comparison. Distinct classes modulo powers and roots
of unity imply distinct classes modulo powers, which is the
direction needed to apply the source's finite-set assertion.
No stronger injectivity hypothesis has been substituted.

For the strict Selmer comparison, restriction to the infinite
cyclotomic kernel is injective for the nonzero negative twist:
the scalar $(1+p^c)^{-E}-1$ acts surjectively on $\Q_p/\Z_p$.
For any finite set of classes, sufficiently high cyclotomic
levels make their cocycles unramified homomorphisms. The local
coboundaries at $p$ are bounded uniformly over every conjugate
decomposition group. Hilbert class field reciprocity, explicitly
retained as a classical theorem, then gives distinct ideal-class
characters with
$\varphi(\gamma^{-1}P)=(1+p^c)^E\varphi(P)$.
The inverse and the negative twist are consistent.
For twist zero the restriction argument fails:
$x\mapsto p^{-m}x\bmod\Z_p$ gives nonzero continuous
characters of $\Z_p$ into $\Q_p/\Z_p$.
The source requires $E\ge1$.

The separate scalar in the homology construction admits the
exact valuation
\[
 v_p\!\left(\prod_{\substack{0\le j\le2E\\j\ne E}}
       ((1+p^c)^E-(1+p^c)^j)\right)=2Ec+2v_p(E!).
\]
This follows because the distances $|E-j|$ take each value
$1,\ldots,E$ twice, and
$v_p((1+p^c)^m-1)=c+v_p(m)$.
At $p=2$ it requires $c\ge2$; if $c=1$, then
$v_2(3^2-1)=3\ne2$.
The actual source imposes $c\ge2$.
The improved valuation proves a sharper scalar cancellation;
it does not alone prove the finite-field homology construction.

The construction-and-detection comparison is now proved by the
explicit finite-field transfer functionals, integral arithmetic
cycles, and homology count described below. Thus the reconstructed
homology exponent $13$ gives strict Selmer exponent $86$ and,
using the stated classical global and local duality theorems,
coarse $H^2$ exponent $88$. These are now proved consequences
of those classical inputs. This is the original route;
the independent alternative which follows gives a stronger bound.

\subsection{An independent coarse bound from algebraic \texorpdfstring{$K$}{K}-theory}

Let $L$ be totally imaginary, $n=[L:\Q]$, $D=|\operatorname{disc}L|$,
$h=\#\operatorname{Cl}(L)$, $E\ge1$, and let $p$ be a prime
with $\mu_p\subset L$. Let $S$ be the primes above $p$ and put
$M=H^2(G_{L,S},\Z_p(E+1))$. The companion manuscript proves
\[
 \log\#M\le \log\#H_{2E}(\operatorname{GL}_{4E+6}(\mathcal O_L),\Z)_{\mathrm{tors}}
       +(E+3)(n\log p+\log h).
\]
Here the general theorems of Quillen, Borel, homological stability,
Quillen--Lichtenbaum and the \'etale--Galois comparison are used
as explicit classical inputs, rather than reproved.
They imply that $K_{2E}(\mathcal O_L)$ is finite, identify its
$p$-primary subgroup with $M$, and identify the required stable
homology with that of the displayed finite-rank general linear group.
The comparison at $2$ applies because $L$ has no real
places.\cite[IV, 1.14.1, 1.17--1.18, 6.9; VI, 8.2--8.3]{weibel-kbook}
The passage to continuous Galois cohomology uses the invertibility
of $p$ on $\operatorname{Spec}\mathcal O_L[1/p]$ and the finiteness
of the finite-coefficient cohomology groups, so that the inverse-limit
correction term vanishes.\cite[II, 2.9 and 2.13]{milne-adt}

For clarity, the quantitative deductions are short. The Kummer
sequence and the Brauer local-invariant sequence give
\[
 \dim_{\mathbb F_p}(M/pM)\le |S|-1+\log_p h\le n+\log_p h.
\]
Indeed, $M/pM$ injects into the second cohomology with
$\mu_p$ coefficients. Its Picard-group contribution has at
most $h$ elements, while the Brauer contribution is the kernel
of the sum map $\mathbb F_p^S\to\mathbb F_p$.
Mathew's theorem for the ordinary Hurewicz map of a connected
infinite loop space kills its kernel in degree $2E$ by
$p^{E+3}$, including $p=2$.\cite[Theorem 6.3]{mathew-hurewicz}
A finite abelian $p$-group of $p$-rank $r$ has at most
$p^{(E+3)r}$ elements killed by $p^{E+3}$. The image of the
finite $K$-group lies in homological torsion. Multiplying the
kernel and image bounds proves the displayed inequality.

The already proved homology exponent $13$ and class-number bound
then give
\[
 \log\#H^2(G_{L,S},\Z_p(E+1))
       \le\exp\bigl((2En)^{14}\log(2D)\bigr).
\]
The strict Selmer dual is a subgroup of this $H^2$, so the same
bound holds for its order. All hypotheses and the elementary
absorption of the additional term are written out in the
companion manuscript. This closes the coarse input for the
common-interval argument. The distinct finite-field homology
construction is established separately below. The other prime
and character components are not settled by this argument.

The original route has also been reviewed further. Its power-detection
lemma is now proved using Chebotarev: a prescribed automorphism of a
finite Galois extension over a number field lifts to the normal closure
over $\Q$; choosing its actual Frobenius at a prime forces the residue
field of the base to have degree one. Kummer theory then detects a
nontrivial global $q$-th-power class at such a prime, provided
$\mu_q$ belongs to the field. The source retains that condition.
Without it the result is false: $16$ is not an eighth power in
$\Q$, since $v_2(16)=4$, but it is an eighth power modulo every
odd prime. At least one of $2,-2,-1$ is a square in any odd prime
field, because their product is a square. If $t^2=\pm2$, then
$t^8=16$; if $i^2=-1$, then $(1+i)^8=16$.

The source's explicit cyclic bar chain and its relation are now
proved as well. The periodic resolution computes
$H_m(C_q\times C_q,\Z/q)= (\Z/q)^{m+1}$.
For $q$ a prime power, $m+2$ vectors in this module have a relation
with a unit coefficient: lift their matrix to the integers and
choose a primitive integer vector in its rational kernel.
For a general composite modulus this conclusion can fail.
For example, $3a+2b=0$ modulo $6$ forces $a$ even and $b$ divisible
by $3$, so neither coefficient is a unit. The actual construction
uses a prime-power modulus. The finite-field cohomology functional
and the arithmetic classes that it evaluates have now also been
constructed in the companion manuscript.

Two other inputs to the arithmetic cycles now have elementary
proofs in the manuscript. For fractional ideals $I_i,J_i$ and
positive rank, an isomorphism with determinant $\beta$ exists
exactly when $\prod J_i=(\beta)\prod I_i$. A Chinese-remainder
choice gives each required determinant-one two-dimensional
operation. In rank zero the determinant is $1$; the actual source
retains the necessary condition $\beta=1$. The integral
interpolation of the characters $(x,y)\mapsto x+ty$ is proved
in the integral group ring by a polynomial finite-difference
identity. No cancellation or division by a zero divisor is used,
and the endpoint where the relevant group-ring element vanishes
is included. These two proofs are inputs to the finite-field
transfer and arithmetic cycle assembly, now proved there as well.

\subsection{Concrete coefficient maps and the global ideal sum}

The manuscript now constructs the coefficient-extraction map for
each transported polynomial lattice. It proves its exact local
coefficient ideal by Smith normal form: if the elementary-divisor
exponents are $c_1\le\cdots\le c_n$ and the partition is
$\nu_1\ge\cdots\ge\nu_n$, its exponent is
$a\,v(\det B)+\sum_i\nu_ic_i$. The determinant twist $a$
may be negative. At infinity its largest singular value, in
the specified Fock and coefficient metrics, is
$\prod_i s_i^{\nu_i}$ for decreasing singular values $s_i$;
the smallest pairs them with the reversed partition.
The proof obtains the row-weight inequalities directly from
restriction to lower triangular matrices. Gaussian integration
and Gram--Schmidt give the needed factorial normalization bound.
The determinant contribution then cancels exactly against the
degree of the underlying Hermitian lattice.

These are estimates for the actual coefficient map. The final
period maps are different maps and still need their own local
valuation bounds. Likewise, evaluation at class representatives
identifies a full product of coefficient spaces only when the
stabilizers are trivial, as required by the relevant level
condition in the source. With a sign action of $\{\pm1\}$
on a one-dimensional coefficient space, the invariant space
would instead be zero. This is another counterexample to
dropping a hypothesis, not to the qualified source statement.

The corresponding decorated integral lattices now have a direct
construction and the needed degree and short-vector estimates.
Start with the unimodular lattice $\mathcal O_E^n$ and coordinate
residue maps. For any unimodular lattice $N$, the common kernel
$L$ of nonzero residue functionals at distinct primes $P$ has
$[N:L]=\prod_P NP$ and $[L^\#:L]=(\prod_P NP)^2$.
When $[E:\Q]=2f$ and the $f$ selected primes have norm $T$,
its normalized determinant degree is $-\tfrac12\log T$.
An explicit Gaussian proof of Euclidean transference, followed
by a codifferent norm estimate, gives $n$ independent vectors
$y_j\in L^\#$ with
\[
 \frac{\operatorname{Tr}_{E/\Q}h(y_j,y_j)}{[E:\Q]}
                      \le n^2\operatorname{rd}(E)^2.
\]
The proof requires integrality of $L$, retains the distinction
between $n$ and $2f$, and does not assert that these vectors form
an integral basis.

The actual automorphic stabilizer argument uses unipotent
reduction at two distinct rational primes. A finite-order
matrix with unipotent reduction at $p$ has $p$-power order:
its prime-to-$p$ part reduces to the identity and is killed
by the elementary congruence-kernel argument. Two different
primes force order one. Positive definiteness supplies finiteness
of the integral isometry group. This holds at ramified primes
and at $2$. One prime alone is insufficient: multiplication
by $\zeta_p$ on $\Q(\zeta_p)$ reduces to the identity at
$(1-\zeta_p)$ but is nontrivial. This two-prime argument is
distinct from the reduction-modulo-$3$ bound used above for
finite stabilizers of the arithmetic complex.

Finally, the passage to every integral period test now has an
explicit proof. If $M$ is an integral form over $\mathcal O_{K_0}$
and an idempotent $e$ has a nonzero integral multiple preserving
$M$, then for every number-field extension $K/K_0$,
\[
 \mathcal O_K M\cap eV=\mathcal O_K(M\cap eV).
\]
Local Smith normal form proves the equality, including at
ramified primes. The finite-group projector denominator serves
only to make an integral endomorphism; it causes no new
denominator in the resulting valuation extension. Rationality
is essential: with $M=\Z^2$ and $e(x,y)=(x,\sqrt2x)$,
$M\cap e\C^2=0$, while after extension to $\Q(\sqrt2)$
the left side contains $(1,\sqrt2)$.
The actual initial period valuation bound is still needed;
this algebraic argument does not prove it.

The first global theta sum has a shorter proof with stronger
decay. For a fractional ideal $\mathfrak a$ in a totally complex
field of degree $2f$, positive widths $x_v$, and
$\Theta^f=\sqrt{N\mathfrak a}/\prod_vx_v\ge1$, the manuscript proves
for $k>2$
\[
 \sum_{0\ne\alpha\in\mathfrak a}
       \prod_v\min\left(1,\frac{x_v}{|\sigma_v\alpha|}\right)^k
 \le\left(\frac{9k}{k-2}\right)^f
       \Theta^{-kf}
       \sum_{j=0}^{f-1}\frac{((k-2)f\log\Theta)^j}{j!}.
\]
The proof first packs disjoint polydiscs using the lower norm
bound for differences of ideal elements. It then integrates
the resulting box count exactly. It applies to the actual
trace-dual and lift lattices with their different powers and
the trace pairing divided by the level integer.
For the source's $k=5$ and width factor $2$, it gives the
uniform bound $(432\Theta^{-4})^f$ under the displayed
archimedean width assumptions. The manuscript also proves the
exact covolumes of the period and lift lattices when the CM
extension is unramified above the level, and the Poisson formula
under continuous product decay. The Fourier frequency is
$2\overline{\sigma(x)}/(gc)$ for the lattice scaled by $c$;
the factor $2$ comes from the complex trace. These normalizations
include residue characteristic $2$.

The manuscript now also proves the archimedean polynomial and
width estimates, conditional on precisely stated mass, dimension,
and formal-degree identities. An exact binary Vandermonde
calculation controls the perturbed spectral tuples. On the
interval where two tuples first differ it gives
\[
  0<\beta_v(Y)\le 80\,t_v\sqrt{k}
                         \sqrt{\frac{\tau}{\Delta}},
  \qquad \frac52\le\tau\le\frac{\Delta}{2},
\]
improving the source's constant $160$. Here $\Delta$ is the
spacing parameter, $k$ the length of the smaller spectral tuple,
and $t_v$ its normalized geometric-mean scale; all are defined
explicitly in the companion manuscript. The mass and dimension
identities, when supplied, give an exact ratio whose remaining
logarithmic error is at most $8f2^{-k}$. This controls
$\prod_v t_v$ and proves the required ordinary and lift width
bounds, including the terminal square root and its factor $2$.
The resulting uniform power estimate for $\Theta$ is proved
with all growing parameters retained. Thus the width calculation
is no longer an unresolved algebraic step. The actual
representation-theoretic coefficient identities are also now
proved. The mass ratio is supplied below from the classical global
mass theorem. The dimension and formal-degree formulas,
the noncompact integrated radial identities for the generated
functions, and the remaining orbital contributions are still
separate tasks.

The logarithmic factor at full decay cannot always be removed.
Take $E=\Q(\sqrt2,i)$, $\mathfrak a=\mathcal O_E$, and
$x_1=x_2=\Theta^{-1}$.
The unit $u=1+\sqrt2$ has absolute values $u$ and $u^{-1}$
at the two complex places. Every integer
$|m|\le\log\Theta/\log u$ contributes exactly
$\Theta^{-2k}$ through the distinct element $u^m$.
The sum is therefore at least
\[
 \left(2\left\lfloor\frac{\log\Theta}{\log(1+\sqrt2)}
               \right\rfloor+1\right)\Theta^{-2k}.
\]
No bound $C^2\Theta^{-2k}$ with a fixed $C$ is possible even
in this one field. The source uses a weaker decay exponent
which avoids this obstruction; its assertion is not contradicted.

\subsection{Evidence and the boundary of this continuation}

The source traces accompanying this version identify the exact
files and statements read, including the concrete links to the
final divisor and theta comparison. The finite checks use exact
rational arithmetic for the prime error and reduction modulo
cyclotomic polynomials for Gaussian and stationary-phase sums.
They include negative controls that violate an explicit
hypothesis. General assertions are supported by the written
proofs, not by extrapolation from these finite computations.

The public MIT verification page was checked again on
21 November 2026. Its displayed last update remains
20 November, 19:45 UTC, with the same independent \emph{nanoda}
acceptance and with the additional comparator/kernel runs still
listed as in progress. No new completed kernel verdict is
claimed. The previous complete local export inventory and
canonical statement comparison remain the formal evidence
described in Section~\ref{sec:complete-export-verification};
they were not rerun without a new reason.

\subsection{Completion of the original coarse Selmer argument}

The finite-field comparison that was previously left open has now
been reconstructed completely. For a finite field $k$ and a prime
power $q\mid\#k^\times$, an explicit cyclic cocycle with coefficients
in $A=\Z/q$ is transferred over the projective lines of a vector
space. Inclusion--exclusion over coordinate subspaces makes the
result additive under direct sums. Its value on a scalar cyclic
cycle with one inserted commuting matrix is
\[
             x^E v(\det g),\qquad v:k^\times\longrightarrow A.
\]
The proof checks the determinant character directly, including
the small fields. No assertion about an exceptional abelianization
of a general linear group is needed.

Integral arithmetic cycles are then constructed using the explicit
cyclic relation, determinant-prescribed ideal-lattice isomorphisms,
and integral character interpolation. The cycles and their common
finite exceptional set of primes are chosen before the residue
functionals. This order of choices is essential for the final
simultaneous power-detection argument.
The rational homology input is the graded structure theorem for
connected cocommutative Hopf algebras, together with Borel's ranks
and integral homological stability. These are stated as classical
inputs, with the precise degree range retained.\cite{milnor-moore,weibel-kbook}
The primitive contribution in odd degree has rank at most
$r_2(L)$. The universal coefficient sequence also retains both
adjacent torsion groups; discarding either would be unjustified.

Here is the resulting quantitative statement. Let $F/L$ be the
cyclic cyclotomic layer of degree $\ell^{\nu-c}$ with
$\gamma\zeta=\zeta^{1+\ell^c}$, where $L$ is totally imaginary,
$\zeta$ has order $q=\ell^\nu$, $c\ge1$, and $c\ge2$ if
$\ell=2$. Put $\kappa=1+\ell^c$, $E\ge1$, and require
$2E+1<\ell^c$. For elements satisfying the source divisor relation,
define
\[
             \mathcal N(\alpha)
                 =\prod_{j=0}^{[F:L]-1}(\gamma^j\alpha)^{\kappa^{Ej}}.
\]
If a finite family $T$ is separated by these norms modulo $q$-th
powers after exponentiation by $\ell^b$, and
\[
                  b\ge 2E(c+1)+2v_\ell(E!),
\]
the proved bound is
\[
 \#T\le q^{r_2(L)}
       \#H_{2E+1}(\operatorname{GL}_{4E+6}(\mathcal O_L),\Z)_{\rm tors}
       \#H_{2E}(\operatorname{GL}_{4E+6}(\mathcal O_L),\Z)_{\rm tors}.
\]
In particular, the additional source hypothesis
$2E+1<r_2(L)$ is unnecessary for this proof. This closes the original
coarse Selmer route. It does not settle the other prime and character
components of the full theorem.

\subsection{Exact integral branching and the actual arithmetic classes}

The compact branching law has an elementary proof in the companion
manuscript. Normal ordering of matrix units gives the highest-weight
classification, and two Gaussian eliminations compute the joint
highest vectors. Their expression as products of initial minors
forces exactly the interlacing inequalities and multiplicity one.
The integral spanning statement is proved separately by a two-column
calculation and subtraction of monic leading terms; no rational
spanning argument is silently promoted to integral spanning.

For interlacing partitions $\lambda=(\lambda_1,\ldots,\lambda_n,0)$
and $\mu=(\mu_1,\ldots,\mu_n)$, set
\[
 a_i=\mu_i-\lambda_{i+1},\quad b_i=\lambda_i-\mu_i,
 \quad D_0=\sum_i a_i,\quad
 R_i=[\lambda_{i+1},\min(\lambda_i,\lambda_{i+1}+D_0)]\cap\Z.
\]
The exact scalar used by the source is
\[
 c=\prod_{i=1}^n\prod_{v\in R_i\setminus\{\mu_i\}}
                      \prod_{j=1}^n\bigl((v-i)-(\mu_j-j)\bigr).
\]
It is nonzero under interlacing. An integral lift of coefficient
extraction is projected by ordered row Capelli operators. Their
eigenvalues show that the resulting map is exactly $cJ$, where $J$
is normalized by its highest-vector coefficient. This proves the
required integral map. Fock adjointness gives the reverse map with
the actual norm ratio and multi-index factorial, and perfect integral
duality proves its integrality. The source's determinant shifts and
conjugate uniformizer are retained in the block covariance identity.

The necessity of a scalar can be seen without computation software.
Take $\lambda=(2,1,0)$ and $\mu=(2,0)$, write
$D_{ij}=x_i y_j-x_j y_i$, and let the highest vector be
$\xi=x_1D_{13}$. The lowering operator gives
\[
       E_{21}\xi=2x_2D_{13}-x_3D_{12}.
\]
Since $E_{21}(x_1^2)=2x_1x_2$, equivariance forces
\[
       J(x_1x_2)=x_2D_{13}-\tfrac12x_3D_{12}.
\]
Its coefficient of $x_1x_3y_2$ is $-1/2$, so it is not integral.
For these partitions the displayed scalar is $c=4$; the submitted
map includes it. This is a counterexample to deleting that scalar,
not to the actual theorem.

The remaining equality of the two coefficient lattices at a large
prime now has a direct proof as well. The differential Capelli
identity, proved by exterior-algebra normal ordering, gives
\[
 C_\nu=[H_\nu,H_\nu]
       =\prod_{(i,j)\in\nu}
                  (\nu_i-j+\nu'_j-i+1).
\]
The factors are the usual hook lengths. Integral duality implies
$C_\nu\nabla_\nu(\Z)\subseteq\Delta_\nu(\Z)\subseteq\nabla_\nu(\Z)$.
For the canonical partition $\nu_N=0$, every hook is at most
$\nu_1+N-2$. Thus $\nu_1+N-1\le p$ makes $C_\nu$ a $p$-unit
and proves equality after localization at $p$. The source uses
the stronger inequality $\nu_1+N<p$. This proof avoids the
modular simplicity assertion altogether for this application.

The adelic class construction is also now explicit. Local Hermitian
bases and approximation identify the classes with decorated
lattices; the additional cosets at the penultimate step are computed
by the Chinese remainder theorem. The proofs retain vector
stabilizers, inverse level characters, and inverse automorphism
weights. Weighted class summation gives the initial period bound
with denominator the least common multiple of the stabilizer orders.
The exact branching covariance then proves the lowering map on
these integral classes for every auxiliary character. Descent from
complex evaluations to the rational integral lattice is proved.
The sharper theta-column divisibility remains a separate arithmetic
requirement; it does not follow just from this initial integrality.

\subsection{Primitive unitary densities and the mass normalization}

The primitive-frame density is now proved for every Gram corank.
Over a residue field of size $q$, for rank $m$ frames in dimension
$n$ with Gram corank $c$, it is
\[
 \prod_{j=n-m-c+1}^n(1-\eta^j q^{-j}),\qquad
 \eta=-1\ \hbox{inert},\quad \eta=1\ \hbox{split},
\]
when $m+c\le n$, and is zero otherwise. The normalization divides
the frame count by $q^{2nm-m^2}$. In the inert case the proof
chooses an orthonormal frame followed by an isotropic frame in its
orthogonal complement. In the split case it counts injective
matrices whose bottom block is injective on the kernel of the
prescribed Gram matrix. These arguments include all degenerate
cases. For example, at $q=2,n=2,m=1$, the zero-Gram counts are
$9$ and $3$, while both nonzero-Gram counts are $6$.

The corresponding integral lifting calculation is also exact.
At an unramified dyadic prime the diagonal Gram corrections are
solved by the surjective quadratic trace map. No division by $2$
is used. Thus the local factor requires $E/F$ to be unramified;
it does not require $F/\Q$ to be unramified at $2$.

Let $f=[F:\Q]$, $d_F=|\operatorname{disc}F|$, and
$\rho=d_F^{1/f}$. Applying the classical unitary Siegel mass
formula with the proved local factors gives
\[
 \frac{M_n}{M_{n-1}}
   =\left((n-1)!\left(\frac{\rho}{2\pi}\right)^n\right)^f
                  d_F^{-1/2}L_F(n,\chi^n).
\]
The theorem and its normalization were checked against
Kirschmer's general statement and the independent displayed
formula of Terakado--Yu.\cite{kirschmer-herm,terakado-yu-mass}
The source's involution condition gives precisely a finite-unramified
CM extension, hence $|d_E|=d_F^2$. Its factors indexed by primes of
$E$ group into the displayed $F$-Euler factors, including the
positive square roots at split primes. Finally
$|\log L_F(n,\chi^n)|\le f\log\zeta(n)$ follows termwise.
This verifies the exact mass input. The general global mass
formula is an explicitly cited classical theorem, not claimed
to have been proved anew in this audit.

\subsection{Archimedean coefficient identities and pointwise positivity}

The polynomial coefficient identities, their leading terms, and the
iterated coefficient extraction are proved in full. Separated finite
interpolation sets are constructed inside the actual spectral
interval, including the short end interval. The resulting estimates
improve the corresponding numerical constants without changing
the allowed range of parameters.

Ordered Capelli minors give the actual terminal radial operator.
Their adjoint and covariance identities identify its action on the
compact highest-weight types. The compact branching law now
constructs the finite idempotents that had previously been an input.
They are orthogonal at every point, so their contraction bound has
constant $1$ and survives integration against any positive measure.
The prescribed shifts and interlacing inequalities make every
terminal density coefficient positive.

The noncompact integration has also been proved under a precise
weighted Fourier hypothesis. If $E(t)=(1+|t|^2)^2D(t)$ is integrable
and its Fourier transform is $O((1+|s|)^{-m-3})$, Fourier inversion
extends $D$ to a $C^m$ function on the whole plane. All derivatives
through order $m$ decay as $O((1+|t|)^{-4})$, and every positive-order
derivative has integral zero. The source's one-coordinate operators
give an explicit sum over subsets of an admissible word; all nonempty
subsets are exactly such derivative terms. This proves the required
integral identity with its shifts and measure retained.
The application must still establish the weighted Fourier bound for
the actual Jacquet-generated functions and identify all of the frame
recursions. Pointwise positivity alone does not supply those estimates.

\subsection{Exact formal degree and source-space dimensions}

The subsequent reconstruction identifies the formal degree exactly.
For $G=U(p,q)$, $n=p+q$, put $s_n=\prod_{j=0}^{n-1}j!$.
If $\lambda$ is a regular Harish--Chandra parameter, then for
the source's compact-dual normalized Haar measure,
\[
                 d(\pi_\lambda)
                   =\frac{\prod_{i<j}|\lambda_i-\lambda_j|}{s_n}.
\]
The proof makes the normalization explicit. On the matrix ball,
the quotient measure is
$\det(1-Z^*Z)^{-n}\,dZ$ with complex-entry Lebesgue measure.
The source measure is this quotient measure, with probability
measure on $U(p)\times U(q)$, multiplied by
$s_n/(\pi^{pq}s_ps_q)$. Successive column integrations of the
matrix-ball beta integral compute the scalar holomorphic
formal degree and fix the parameter-independent constant.
The proof also includes the compact signatures and the trivial
rank-zero group.

The representation in the source is specified through a minimal
compact type, so a formula for an unspecified discrete series
would not suffice. The reconstruction checks the exact shift
$\mu_0=\lambda+\rho_{\rm n}-\rho_{\rm c}$.
Blattner's formula proves that this type occurs once and has
strictly smaller $\|\mu+2\rho_{\rm c}\|^2$ than every other type.
Conversely, Parthasarathy's characterization identifies any
irreducible unitary representation satisfying the source's
minimum condition. The Harish--Chandra, Blattner and
Parthasarathy theorems are explicit classical inputs, checked
against primary references; they are not presented as new
proofs of the whole representation theory.
\cite{harish-chandra-ds-ii,duflo-heckman-vergne,baldoni-vergne,parthasarathy-minimal}

The exact source-space dimension is also proved. Evaluation on
decorated double cosets gives one copy of the dual coefficient
space for each class. Trivial stabilizers follow from the actual
two-prime level. Counting flags and nonzero functionals supplies
the finite-level factors, while the compact branching law proves
the Weyl dimension formula. Fixed spaces and the simultaneous
kernels of the auxiliary norm operators are computed by the
free finite-group action. These results verify the dimension
input; comparison of the ambient multiplicity with its expected
mass/formal-degree leading term remains a separate estimate.

\subsection{Ordinary moments, torsion phases and the integral iteration}

The ordinary moment argument now has a full functional-analytic
proof. A polynomial functional extends to a bounded measure on
$\mathbb Z_p^r$ exactly when its values on the multivariable
binomial polynomials are uniformly bounded. The proof constructs
the Mahler expansion and identifies its norm. Formal substitution
of $\log(1+T_i)$ gives the precise coefficient criterion.
Integrality descends from a tensor power by taking the coefficient
with the same multi-index in independent variable blocks.
This would not follow by replacing the tensor power by the
ordinary pointwise power of a function.

There is a concrete test of the moment hypothesis. The functional
$P\mapsto P'(0)$ has bounded power moments, but
\[
 \left.\frac{d}{dX}\binom{X}{p^a}\right|_{X=0}
                       =\frac{(-1)^{p^a-1}}{p^a}
\]
has unbounded $p$-adic absolute value as $a\to\infty$.
Thus bounded ordinary derivatives alone are insufficient.
The submitted construction uses the stronger binomial-moment
condition, so this is not a counterexample to its assertion.

For the actual large split primes, a shorter classical route
supplies the one-dimensional CM theta laws.
The cited integral theory of reduced theta functions on a good
ordinary CM elliptic curve supplies the geometric input.
\cite{bannai-kobayashi-2007}
The reconstruction identifies the even theta function and its
zero divisor, fixes one discriminant-normalized unit period,
and proves integrality for all required prime-to-$p$ translates.
A saturated expansion map on global sections proves preservation
of the norm of arbitrary finite linear combinations. This retains
cancellations among their coefficients; estimates for the
individual summands would be insufficient.

The selected $p$-power translation is checked with its exact phase.
If $y=s\nu/p^a$ and the chosen $p$-adic valuation of
$\bar\nu$ is at least $2a$, then $|y|^2\in\mathbb Z$.
An elliptic quotient built from a CM element and its conjugate
identifies translation in the formal multiplicative coordinate
with $e^{\pi i|y|^2}U_y$ on the complex theta function.
No translation by a nonzero torsion point is replaced by its
formal logarithm, which would lose that point.
The normalization is simultaneous for all orders of torsion.

The finite-dimensional integral iteration is proved with the
actual determinant factors and local norms. It retains the
identification of the higher-dimensional normalized theta
function, its support, and the exact incidence commutation
as explicit inputs. In particular its quadratic loss in the
number of iterations is derived from that commutation, not
inferred from integrality of each individual map.
The one-dimensional CM laws alone do not establish those
remaining higher-dimensional identities.

\subsection{Regular global differential operators and exact radial actions}

The inverse terminal operator has a globally regular
coefficient on the required positive-Gram region. Writing
$D$ for the final minor and $d$ for the preceding minor, its
only scalar denominator is a power of
$\mathcal G=|D|^2-TT'|d|^2$. A cofactor vector has squared
norm $\mathcal G/T'$ and supplies the compact Capelli
contraction. Finite polynomial interpolation gives the inverse
on every occurring type, including repeated eigenvalues.
Its derivatives have the required polynomial bounds away from
$\mathcal G=0$.

An explicit rank-one example in the companion manuscript has
$d=0$ and $\mathcal G=1$. Thus formulas dividing by $d$
would fail inside the allowed region. The source's cofactor
construction remains regular there, so this check supports
the actual construction.

All four polynomial differential lifts are now constructed.
The two isotropic Capelli minors give the Toda recurrence by
last-column expansion and the two unipotent character laws.
The other two operators follow from a bordered-minor identity;
the full boundary contraction has exactly the sign of the
terminal $(2,2)$ entry. Both ordinary and conjugate parameter
shifts agree with the source. Their full-group coefficients
are exterior determinants. The compact stabilizer can have
nontrivial determinant, and the proof explicitly cancels that
character in the covariance calculation.

These operators and their ordinary Haar adjoints preserve the
specified smooth rapidly decreasing class. This fact does not
by itself identify adjoints for the Jacquet pairing.
The terminal quantum determinant and all its commutation relations
are now proved in the companion manuscript. The quantum
minor calculation preserves the shifts by one, the order of
every operator, and differentiation of variable coefficients.
The full bordered determinant has two extra noncommutative terms;
these are retained until the actual character laws annihilate them.
A one-pair Bruhat-cell extension and an elementary polynomial
independence argument recover every terminal commutation relation,
including the exceptional parameters by polynomial continuation.
Both shifted adjugate identities follow, and an inverse is asserted
only away from the roots of the scalar determinant.
The actual integral pairing formula, full frame recursion, and
uniform square-integrability and Fourier decay for the generated
functions remain to be established.
The additional scalar-rigidity proof for Hermitian flags likewise
retains the actual conclusion: the successive scalars and norm
are fixed, while ordered frame vectors can differ by lower
flag terms. It makes no false assertion of frame equality.

\subsection{The theta normalization and lowest compact-type decay}

The higher-dimensional analytic theta transformation now has
a conventional proof. Normalized finite averages between
commensurable self-dual lattices are unitary. Their triple
composition is a scalar whose eighth power is one.
An elementary quaternion calculation proves the necessary finite
quadratic Gauss-sum formula, including even-order groups.
The nondegenerate Cayley quadratic form and an absolutely
convergent Gaussian trace then give the exact relation
$c^8=\det(g_1g_0^{-1})^4$.
The invariant sign semicharacter is a checked hypothesis.
The nonprincipal ideal used as the elliptic base has the
required norm semicharacter after the specified half-period
translation. The determinant valuation of the separate scalar
$\kappa$ makes $\kappa c$ a unit at the selected split prime.
This closes the analytic normalization; arithmetic local
standardization and the actual column-support and incidence
identities are separate.

The lowest compact-type compression is identified with Sun's
classical integral formula, with probability compact Haar
measure and the exact dimension and character factors.
\cite{sun-eisenstein,vogan-zuckerman}
A finite exponential-moment estimate gives coefficient decay
with constant $h/(2\exp(1))$, stronger than the needed $h/8$.
The restriction to vectors in the lowest compact type is
essential: choosing $u=\sigma(a)v$ for arbitrary unit vectors
would give a coefficient of absolute value one.
The rational-element estimates also pass review.
Integral traces and the large split congruence prime force
nearby rational isometries to be unipotent; they need not
be the identity. A nontrivial integral unipotent tends to
the identity after conjugation into a cusp.
The covolume separation, matrix packing and complete distant
orbital sum are proved. The kernel/trace construction and
the near-unipotent, rare-cusp and deep-cusp estimates remain
substantive inputs.

\subsection{A stronger opposite-parity Selmer estimate}

There is a shorter classical proof of the entire odd-prime
opposite-parity estimate. For an abelian CM field $E$ of
degree $n$ and conductor $N$, write $b_p(E,e)$ for the
$\mathbb Z_p$-length of the strict Selmer eigenspace of sign
$(-1)^{e+1}$. The companion manuscript proves, for $t\ge2$,
\[
 \sum_{p\ne2}\log p\sum_{1\le e<t}b_p(E,e)
       \le n(t^2+t-2)\log(tN)
       \le\frac98nt^2\log(tN).
\]
The proof uses Greither--Popescu's established special-value
theorem over $E/E^+$, with Ferrero--Washington supplying
its vanishing-$\mu$ hypothesis.
\cite{greither-popescu,ct-fw}
This Galois group has order two, so its integral idempotent
exists for every odd $p$, including those dividing $n$.
Equivariant Poitou--Tate duality and injective localization
bound the exact strict Selmer eigenspace by that cohomology
group. A rational product of imprimitive Dirichlet values,
an explicit Bernoulli denominator, and the bound for higher
roots of unity give the displayed uniform estimate.
The valuation sum is bounded using a nonzero integer, so
possible denominators at two are retained.

This proves the source's opposite-parity estimate with a
stronger constant and no extra degree-dependent term.
It is an independent conventional route using the named
classical theorems. It does not verify every step of the
source's characterwise proof or control the complementary
parity, where the corresponding special value vanishes.

\subsection{Exact arithmetic equivariance and the weighted unipotent integral}

The theta reconstruction now treats the nonprincipal elliptic
ideal with one ordinary period and proves local standardization
of the fourfold lattice. Its ordinary integrality theorem applies
to split primes at least five prime to the ideal norm, which
includes all sufficiently large selected primes. The source's
more general interface also allows the prime three; the present
classical route does not establish that additional case.

The exact equivariance of the arithmetic family also passes
review. Unramifiedness at every finite place supplies a common
hyperbolic decomposition for two self-dual lattices, including
at dyadic places. A finite quadratic sum then has positive sign,
which proves the exact commutation of the lattice moves.
The unitary covariance retains the determinant character.
An eighth-power identity alone would leave a scalar ambiguity;
the arithmetic argument removes it.

For the Jacquet estimates, the basic coordinate integral has
an explicit gamma-product value and a necessary and sufficient
open convergence region. Both weighted averages and the full
unipotent induction have been proved with their Jacobians.
The negative norm of the terminal vector makes the needed
character coefficient nonzero. The examples of divergence
at the boundary of the exponent region do not satisfy the
source's actual hypotheses. The full-group passage has also passed
review: its Haar density and compact rank-zero integral give the
required separated coefficient bound and continuous majorant.
The geometric integration formula and full regularized pairing now
also pass review. Fourier averaging supplies an integrable majorant
before the global unfolding. A fifth derivative gives continuity under
dilation and upgrades equality almost everywhere in the characters to
equality everywhere. The adjoint reverses the derivative word and
differentiates its coefficient; both operations are necessary. The
actual operator commutation and regularity conditions for the four
Fourier lifts are now matched below; the full initial matrix and its
successive restriction identity remain separate.

The automorphic kernel identity has been proved pointwise. Applying
the orthogonal projection to approximate identities first constructs
its evaluation vectors in $L^2$; the diagonal formula then follows
from Parseval. Smoothness and the polynomial lattice bound allow
the classical finiteness theorem to be applied with its actual
hypotheses\cite{getz-hahn-2022}. Thus this step does not assume an
integrable diagonal or infer a diagonal identity from equality almost
everywhere on a product.

The rare cusp-volume estimate has passed an independent check of
its exact constants. A sublattice of least size less than one is
isotropic. Counting primitive isotropic sublattices of each rank,
with the stated classical Tamagawa and residue inputs, gives
the stronger bound $\exp(-nf\log n/3)$ for the probability that
the least size is less than $1/2$\cite{weil-siegel,louboutin-residue}.
The nearby-unipotent integral has also been proved. The canonical
isotropic flag has strictly increasing least sizes; an integral
automorphism sufficiently close to the identity acts trivially on
its graded pieces. Packing its independent matrix blocks and unfolding
over all primitive isotropic flags gives the required
$\exp(-nf\log n/4)$ bound. The deep-cusp integral now also has a
complete ordinary proof. The trace-dual calculation for the radical
retains the different factors of two in its Hom and skew-Hermitian
parts, including at dyadic primes. Euclidean transference gives a
small fundamental domain\cite{banaszczyk-transference,bost-kunnemann-extensions}.
The actual projection vectors have zero constant terms, by local
discrete-series cuspidality and the absolutely convergent kernel
unfolding\cite{khare-larsen-savin}. A direct product-chart Sobolev
inequality and $8qn^2f$ iterations of the zero-average identity give
deep-cusp mass at most $\exp(-n^2f)$ times the projection rank.
The volume estimate uses only the Hecke functional equation and a
strip maximum principle\cite{zaman-hecke-estimates}. Together these
arguments establish the full ambient multiplicity comparison under
the stated numerical and representation-theoretic hypotheses.

The higher-dimensional ordinary moment laws and the two actual
theta-column supports have now passed independent review. A finite
orthogonal-frame decomposition transfers the exact elliptic torsion
law to the lattice theta function. The finite Fourier transform of
$\ell'X=\ell$ has an annihilator of order $p^{fa}$. In the conjugate
ordinary type, Fourier inversion restricts a norm-one measure. In the
selected type, the annihilator is a companion lattice quotient; its
index square root and determinant character have cancelling valuations.
The surviving factor is precisely $p^{-fa}$. The actual indicator's
two volume factors cancel before taking any valuation. The adjacent
Hecke incidence identity has also passed independent review. For each
matrix the two neighbor sets are simultaneously empty or singletons.
Finite Fourier projection to a common theta lattice gives the exact
square root of its volume ratio; the determinant character then gives
the scalar $p^{(b-a)/2}$ for a move from rank $a$ to rank $b$.
The selected primes have residue norm $p$. Common representatives
satisfy both original indicator definitions, by simultaneous Chinese
remaindering at the selected and conjugate primes. The global class
sum retains inverse-stabilizer weights. An explicit $S_3$ example
shows that omitting those weights would give unequal sums $3$ and $1$.
The actual weighted sums both equal $1/2$.

The coefficient-field issue is also resolved for each fixed finite
construction. Finiteness of the norm-one ideal class quotient,
self-dual lattice classes, residue quotients and derivative orders
gives finitely many algebraic generators. Rational unitary covariance
introduces no additional generators. All prime and level data are
fixed before this field is enlarged. This does not assert one number
field for unbounded independent levels or derivative orders.

The actual polynomial kernel has also been checked directly against
its definition. The determinant multiplication and matrix substitution
give the universal polynomial used in the valuation argument.
Each coefficient already belongs to the target representation, so
the arbitrary linear projection onto that representation has no
effect. The derivative convention uses ordinary mixed derivatives;
no factorial division is present. Finite summation gives precisely
the stated moment table and one determinant-character ratio.
Together with the preceding verified inputs, this proves the
selected-prime column estimate on its prescribed modified initial
lattice. It does not establish that this lattice has a nonzero theta
image or prove the different estimates at the auxiliary level primes.

For the complementary Selmer parity, the exceptional trivial-character
case is now proved directly. Equality of the global and $p$-local
cyclotomic invariant subgroups gives a unique residual lift satisfying
the strict local condition. It injects into characters of the class
group of the residual character field. That field is $\mathbb Q$ or
$\mathbb Q(\sqrt p)$ in the exceptional case; in the latter case
$h(\mathbb Q(\sqrt p))<p$. Hilbert class-field reciprocity therefore
gives vanishing of the entire strict Selmer group, including ramified
coefficient fields\cite{milne-cft}.

Two further complementary-parity cases now have direct proofs.
For the fixed-conductor case, the odd-order condition removes the
components at two and the two varying auxiliary primes. Restriction
to a fixed cyclotomic field then has a kernel bounded by
$n(a+v_p(e))$ in normalized length. This uses no division by the
field degree and remains valid for arbitrarily ramified coefficient
fields. Summing over exponents retains the common-interval bound
additively, with a fixed coefficient. For the bounded-order part at
exponent one, the exact small-conductor threshold bounds the degree
of the character-killing field before any auxiliary parameter grows.
The coarse Selmer estimate then gives a fixed polynomial bound.
An explicit uniform bound on the number of these characters permits
their sum over the fixed finite prime set. Neither argument controls
the remaining automorphic cases.

The actual recursive radial frame has now passed a complete ordinary
algebraic review. The initial connection is constructed from the
leading coefficients of the terminal matrix. Its commutators,
both recursive transformations, the weighted congruences and all
spectral exclusions are proved with noncommuting multiplication
order retained. In particular the minus transformation uses the
opposite Schur factorization and never assumes an inverse at an
unneeded adjacent parameter. The shifted spectral margins make
every scalar density positive. Integration by parts proves the
resulting positive form for compactly supported vectors. The
$309$ exact controls include noncommuting matrices and a genuine
two-variable differential-operator calculation; they corroborate
the general proof and are not used as a substitute for it. The
domain of the scalar determinant and the domain of the upper
Gauss coordinate were made explicit during review. These were
clarifications to the reconstruction, not counterexamples to the
submitted theorem. The further analytic steps are described next.

\subsection{Actual Fourier decay and the degenerate orbital terms}

The actual four group operators now yield a complete Fourier-decay
argument. Their Haar adjoints reduce to the radial adjoints because
the nilpotent groups are unimodular, their characters are unitary,
and the radial normalization is an isometry for the exact Haar density.
A shifted two-by-two differential identity telescopes along the
outer coordinates. The remaining terminal Gauss terms cancel in
the specified multiplication order. Analytic adjointness for the
regularized pairing then permits iteration. Its initial bound uses
unit-modulus multiplication and Cauchy--Schwarz, so the bound is
uniform in the Fourier parameter. Polynomial first-minor insertions
give decay of every derivative as well.

For compact group inputs, a direct one-pair Jacquet integral gives
a smooth compact family on the smaller group, with the exact factor
$|t|^{2r+k}$. The full Gram polynomial becomes $|t|^2$ times the
smaller Gram polynomial. This supplies a common positive cutoff
and compact support on every compact parameter set away from zero.
The resulting marginal density is continuous at every nonzero
parameter and therefore agrees there with its Schwartz Fourier
representative. This repairs a possible invalid passage from an
almost-everywhere Fourier identity to a pointwise conclusion.
It is a verified argument, not a counterexample to the source.
Generated slices retain this local regularity and have compact
terminal support. The full initial matrix and its complete restriction
chain are still required for the remaining residue-form application.

For degenerate flags, a direct geometric normal form now gives a
common positive Hermitian space of dimension $k$. Its compact
unitary group acts on the two column representations. The arithmetic
unipotent stabilizer is a cocompact lattice, but the compact group
need not preserve that lattice. The proof first disintegrates over
the real unipotent quotient, where compact conjugation preserves
Haar measure. Irrelevant terms vanish by character averaging;
relevant degenerate terms intertwine the two compact types.
Distinct parameters give distinct highest weights, directly from
the pointed cone of positive roots. Schur's lemma proves vanishing
of every absolutely convergent degenerate off-diagonal term.
The $2399$ exact controls corroborate the normal form and include
a failure example when relevance is deleted. That example does
not contradict the stated vanishing theorem.

The Euclidean volume estimate in the diagonal argument also has
a shorter proof. Successive orthogonal projections give an exact
product formula for negative moments of a projected Gram determinant.
It improves the two submitted orbital volume constants by factorial
factors. Direct integration in the last perpendicular component
gives the optimal threshold exponent $2(d-j+1)$ for $j$ vectors
projected onto a $d$-dimensional complex space; an explicit positive
lower bound proves that a larger uniform exponent is impossible.
The original final dyadic estimate has also been strengthened:
its constant $16$ becomes $4\cdot2^{-j}$, and its decay exponent
gains one power of the size parameter. Deleting the relation
$D\le(SS')^r$ makes even the original estimate false, as an exact
finite example with $1025$ points shows. The source retains that
relation. The flag-coordinate reductions and weighted lattice-pair
count are separate inputs to the full diagonal amplitude bound.

\subsection{The whole-field comparison and rational adjoint vanishing}

The whole-field Selmer comparison now has an integral proof with
its error quantified. For a positive negative-Tate exponent, the
unramified local cohomology at a prime away from $p$ is zero:
Frobenius minus one is a nonzero scalar and is surjective on the
divisible coefficient module. Shapiro's lemma thus preserves the
required zero decomposition conditions. The integral character
transform has both compositions equal to multiplication by the
field degree. It is an isomorphism when that degree is a $p$-adic
unit; otherwise its kernel is killed by the corresponding power
of $p$. A nonzero vector in the Fourier kernel for a cyclic group
of order $p$ shows why this error must be retained.

For the growing CM fields, only the fixed odd primes dividing
the original character order contribute this error. Finite-coefficient
restriction, the ray class exact sequence, and the principal-unit
structure bound it by a fixed multiple of
$N_h^2(t-1)\log N_h$\cite{milne-cft}. Normalized coefficient lengths
cancel the ramification index exactly. Consequently each character
bound, and each additive common-interval bound, keeps coefficient
one in the whole-field comparison.

The actual adjoint-vanishing implication follows from Thorne's
Theorem~6.1\cite{thorne-adjoint-2022}. Its cyclotomic-kernel
irreducibility hypothesis is exactly the one supplied by the
automorphic point. The Frobenius normalization, finite coefficient
field, arbitrary polarization matrix and local Bloch--Kato conditions
have been compared explicitly. The adjoint action at complex
conjugation has the required minus sign\cite{newton-thorne-adjoint-2023}.
An equivariant projection retains the primitive in the symplectic
Lie algebra when that summand is needed. The cited theorem's full
published proof is an external classical input to this reconstruction.
Rational vanishing does not bound the length of an integral torsion
module: $\mathbb Z_p/p^a\mathbb Z_p$ has arbitrarily large length
and becomes zero over $\mathbb Q_p$. The automorphy construction
and quantitative congruence estimates therefore remain essential.

The unresolved tasks now have more precise boundaries. They are
the remaining cohomology groups of the complementary sign; the concrete integral
theta families and their
uniform global estimates; and the final assembly into the fixed
smoothed prime combination. These are unfinished parts of our
human reconstruction. No fatal flaw in the submitted proof has
been established, and no completed conventional proof of GRH is
claimed in this version.


\section{Findings and the remaining mathematical work}
No counterexample satisfying all hypotheses of a consequential assertion
actually used in the submitted proof has been established. No fatal
mathematical defect or additional foundational axiom has been found.
This conclusion is based on the specific checks described in this report;
it is not a claim that every possible counterexample has been excluded
by a human reading of the source.

The complete local axiom inventory and structural comparison, together
with the published independent kernel acceptance, are substantially
stronger formal evidence than the selected compilations available in
the first version of this report. They resolve the earlier absence of
evidence about the complete exported proof. Their exact provenance,
scope, and remaining implementation distinctions are in
Section~\ref{sec:complete-export-verification}.

The continued mathematical checks establish additional positive results:
the Bernoulli-number unit criterion with its exact inequalities; the
restricted cyclotomic Selmer vanishing with its full local conditions;
the regular-semisimple element assertion in the indicated four-dimensional
cosets; and several useful analytic and determinant estimates from the
related literature. An explicit nonzero even-character Selmer class,
a Heisenberg-group example, and the arbitrary-coefficient mean-square
counterexample test plausible overgeneralizations. None contradicts
the qualified assertions actually used.

For the requested conventional proof, three connected tasks remain:
\begin{enumerate}
\item Complete the other prime and character components,
      including the relevant automorphy and deformation arguments.
      The cyclotomic generator and common-interval arguments are now
      reconstructed, and their coarse cohomological input follows
      from the separate $K$-theory argument.
\item Construct the actual integral theta families and prove their
      norm, near-orthogonality, and divisibility estimates uniformly
      along the growing chains, with all measures and lattices fixed.
      The local harmonic-analysis formulas, divisor comparison, and
      finite-group dual exponent have been reconstructed, but they
      do not by themselves estimate the global orbital sums.
\item Combine those constructions to bound the specific fixed positive
      combination of smoothed prime sums by $x^{1/2+\epsilon}$ on
      arbitrarily large multiplicative intervals. The companion
      manuscript proves the subsequent analytic implication completely
      and supplies the selected-character bounds under the explicit
      uniform QRH hypothesis.
\end{enumerate}
These are unfinished parts of our human reconstruction, not alleged
missing proofs in the submitted archive. The possible alternative
through improved M\"obius cancellation has its own precisely stated
shorter mean-square hypothesis. Neither approach is presented as a
completed ordinary proof of GRH.

\appendix
\section{Reproducibility files}
The accompanying \path{grh_audit_reproducibility.zip} contains both
standalone TeX documents, the source-identification records, exact
arithmetic checks, export-inventory and comparison programs, positive
and negative controls, and selected primary verification records.
The guide \path{EXPORTED_PROOF_AUDIT.md} specifies the public inputs
and commands for the full export comparison. The compressed export,
its reproducible indices, the original source archive, and the
toolchain binaries are not duplicated in this bundle.

The source downloader verifies each archive part and the combined
SHA-256 digest before extraction. Extraction rejects links and paths
outside the destination. The import graph distinguishes module
reachability from theorem-proof dependency. The finite arithmetic
programs distinguish their exact checked identities from the
cohomological and analytic proofs given in the text.

The full proof inventory reads every record of the decompressed
stream and checks its digest. The definition comparison follows
constant references from the five target types, with the exclusions
specified in Section~\ref{sec:complete-export-verification}. Its
controls explicitly demonstrate that it detects changed definitions
but does not validate root proof bodies. A separate kernel is needed
for that last task. The saved external logs document the published
kernel result and are labelled as external evidence.

Earlier build records are retained for provenance. Deliberate resource
termination and an unsupported external-checker feature are not
recorded as mathematical rejections. No local complete kernel replay
is asserted. Full SHA-256 values and file inventories in the bundle
allow these results to be tied to the versions actually examined.

\begin{thebibliography}{99}
\small
\raggedright
\bibitem{announcement}
The announcement at \url{https://alpo.ge/remark.pdf}, four pages,
and its fourteen linked archive parts, downloaded 20 November 2026.
The exact hashes are recorded in Section~1 and in the accompanying files.

\bibitem{mathlib}
Mathlib, \path{Mathlib/NumberTheory/LSeries/RiemannZeta.lean},
definition \code{RiemannHypothesis}, at revision
\path{0df444a360eaa60ab8c11dca51a86af692955474}.
\url{https://github.com/leanprover-community/mathlib4/blob/0df444a360eaa60ab8c11dca51a86af692955474/Mathlib/NumberTheory/LSeries/RiemannZeta.lean}.

\bibitem{gan-qiu-takeda}
W.~T. Gan, Y. Qiu and S. Takeda,
\emph{The regularized Siegel--Weil formula (the second term identity)
and the Rallis inner product formula},
arXiv:1207.4709, version 3 (21 January 2014),
especially Sections 1.1--1.3 and Theorem 11.7.
\url{https://arxiv.org/pdf/1207.4709}.

\bibitem{gan}
W.~T. Gan, \emph{Automorphic forms and the theta correspondence},
Arizona Winter School 2022 notes,
Section 1.5 (the definite unitary example) and Section 3.9 (Witt towers).
\url{https://swc-math.github.io/aws/2022/2022GanNotes.pdf}.

\bibitem{kato}
K. Kato, \emph{$p$-adic Hodge theory and values of zeta functions of
modular forms}, Ast\'erisque \textbf{295} (2004), 117--290.
The noncentral/central distinction was checked in the introductory
statement on printed page 119 of the publisher's sample:
\url{https://smf.emath.fr/sites/default/files/2026-03/AST_2004__295__117_2__sample.pdf}.
This check is not a new verification of the complete proof of Kato's theorem.

\bibitem{lean-release}
Lean 4.33.1 release notes, official Lean language reference,
especially the backported kernel fixes:
\url{https://lean-lang.org/doc/reference/latest/releases/v4.33.1/}.
Accessed 20 November 2026.
\bibitem{S-endpoint}
Submitted archive. Endpoint definitions and the final declarations.
\begin{itemize}[leftmargin=1.4em,itemsep=0pt,topsep=1pt]
\item \path{Statements.lean}
\item \path{Main.lean}
\end{itemize}
\bibitem{S-character}
Submitted archive. Primitive reduction and reflection.
\begin{itemize}[leftmargin=1.4em,itemsep=0pt,topsep=1pt]
\item \path{Development/AutomorphicArithmetic/CharacterBoxes/ClaimStrip.lean}
\item \path{Development/AutomorphicArithmetic/CharacterBoxes/Imprimitive.lean}
\end{itemize}
\bibitem{S-assembly}
Submitted archive. Assembly of the arithmetic, geometric and analytic inputs.
\begin{itemize}[leftmargin=1.4em,itemsep=0pt,topsep=1pt]
\item \path{Development/AutomorphicArithmetic/CharacterBoxes/MainConclusion.lean}
\item \path{Development/AutomorphicArithmetic/ThetaCorrespondence/ThetaPairLifts/ThirdPlaceSeries/EveryCharacterOfLiteratureOfVWin.lean}
\end{itemize}
\bibitem{S-size}
Submitted archive. The scalar size inequality.
\begin{itemize}[leftmargin=1.4em,itemsep=0pt,topsep=1pt]
\item \path{Formalism/SizeWindowBrackets/GrowthBoundGivesRiemannHypothesis.lean}
\item \path{Formalism/SizeWindowBrackets/ChainSlopeWindows/LipschitzWeightAndMinimum.lean}
\end{itemize}
\bibitem{S-weak}
Submitted archive. The explicitly refuted weaker error estimate.
\begin{itemize}[leftmargin=1.4em,itemsep=0pt,topsep=1pt]
\item \path{Formalism/SizeWindowBrackets/ChainSlopeWindows/SubThreeMulLogTwoLe.lean}
\end{itemize}
\bibitem{S-landau}
Submitted archive. Pole data and the analytic criterion.
\begin{itemize}[leftmargin=1.4em,itemsep=0pt,topsep=1pt]
\item \path{Development/AutomorphicArithmetic/DedekindLFunctions/LandauZeroFreeCriterion.lean}
\item \path{Development/AutomorphicArithmetic/DedekindLFunctions/LandauTheorems/landauWindowIneq.lean}
\item \path{Development/AutomorphicArithmetic/DedekindLFunctions/LandauTheorems/landauTails.lean}
\item \path{Development/AutomorphicArithmetic/DedekindLFunctions/LandauTheorems/lemmaL.lean}
\end{itemize}
\bibitem{S-cohomology}
Submitted archive. Twisted Galois cohomology and local conditions.
\begin{itemize}[leftmargin=1.4em,itemsep=0pt,topsep=1pt]
\item \path{Formalism/GaloisCohomology/Statement.lean}
\end{itemize}
\bibitem{S-vwin}
Submitted archive. The precise cohomological interval assertion.
\begin{itemize}[leftmargin=1.4em,itemsep=0pt,topsep=1pt]
\item \path{Development/AutomorphicArithmetic/CharacterBoxes/ExistsPageN.lean}
\end{itemize}
\bibitem{S-counts}
Submitted archive. Character action, parity, and restricted vanishing.
\begin{itemize}[leftmargin=1.4em,itemsep=0pt,topsep=1pt]
\item \path{Development/AutomorphicArithmetic/AutomorphicVarieties/CharacterLevelBounds/ShaFamMono.lean}
\item \path{Development/AutomorphicArithmetic/AutomorphicVarieties/HeightAndVanishing/Counts.lean}
\item \path{Development/AutomorphicArithmetic/AutomorphicVarieties/LevelCohomology/PropZ.lean}
\end{itemize}
\bibitem{S-kato}
Submitted archive. The noncentral Kato input and its rational Selmer application.
\begin{itemize}[leftmargin=1.4em,itemsep=0pt,topsep=1pt]
\item \path{Development/AutomorphicArithmetic/AutomorphicVarieties/LiteratureResults/KatoEulerSystems/KatoLitOOf.lean}
\item \path{Development/AutomorphicArithmetic/AutomorphicVarieties/ModularityArgument/KatoQuadraticCharacters/SphericalDecomposition/KatoTwoW.lean}
\item \path{Development/AutomorphicArithmetic/AutomorphicVarieties/ModularityArgument/ThresholdArguments/DenseThresholdCases/KatoTwoWOfPublishedH.lean}
\end{itemize}
\bibitem{S-fw}
Submitted archive. The abelian Ferrero--Washington assertion.
\begin{itemize}[leftmargin=1.4em,itemsep=0pt,topsep=1pt]
\item \path{Development/AutomorphicArithmetic/AutomorphicVarieties/PublishedTheorems/FerreroWashington.lean}
\end{itemize}
\bibitem{S-lemmaK}
Submitted archive. The coarse absolute cohomology bound.
\begin{itemize}[leftmargin=1.4em,itemsep=0pt,topsep=1pt]
\item \path{Development/AutomorphicArithmetic/AutomorphicVarieties/GeneratorCountTuples/LemmaK.lean}
\end{itemize}
\bibitem{S-automorphy}
Submitted archive. Substantial rank-three and rank-four reductions.
\begin{itemize}[leftmargin=1.4em,itemsep=0pt,topsep=1pt]
\item \path{Development/AutomorphicArithmetic/AutomorphicVarieties/ModularityArgument/ThresholdArguments/DenseThresholdCases/WindowBoundOfLiteratureLevel.lean}
\item \path{Development/AutomorphicArithmetic/AutomorphicVarieties/ModularityArgument/ThresholdArguments/DenseThresholdCases/SymmetricCubeTheoremOfEndomorphisms.lean}
\item \path{Development/AutomorphicArithmetic/AutomorphicVarieties/ModularityArgument/ThresholdArguments/DenseThresholdCases/ModularityRankFourOfPromodularity.lean}
\item \path{Development/AutomorphicArithmetic/AutomorphicVarieties/ModularityArgument/ThresholdArguments/DenseThresholdCases/PmTwo.lean}
\end{itemize}
\bibitem{S-irreducibility}
Submitted archive. Residual separation in the irreducibility argument.
\begin{itemize}[leftmargin=1.4em,itemsep=0pt,topsep=1pt]
\item \path{Development/AutomorphicArithmetic/AutomorphicVarieties/UnitaryThreeDeformations/IrrOpen.lean}
\end{itemize}
\bibitem{S-idempotent}
Submitted archive. The denominator of a central idempotent.
\begin{itemize}[leftmargin=1.4em,itemsep=0pt,topsep=1pt]
\item \path{Formalism/PadicSubsetSums/CentralIdempotents/CentralIdempotentIntegral.lean}
\end{itemize}
\bibitem{S-cut}
Submitted archive. Covolume normalization and orthogonal restriction.
\begin{itemize}[leftmargin=1.4em,itemsep=0pt,topsep=1pt]
\item \path{Formalism/HarderNarasimhan/StatementPl.lean}
\item \path{Formalism/MeanHeightDescent/AverageDegreeBounds/GoodBadCut.lean}
\item \path{Development/AbelianVarietyTheory/Proofs/DivisorProofs/MajAtAddRight.lean}
\end{itemize}
\bibitem{S-rational}
Submitted archive. Rational subspaces and integral intersections.
\begin{itemize}[leftmargin=1.4em,itemsep=0pt,topsep=1pt]
\item \path{Formalism/SizeWindowBrackets/CutFailsForRSubmodule.lean}
\item \path{Development/AbelianVarietyTheory/Proofs/TwistedOperators/LiftsRationalOfRatAt.lean}
\item \path{Development/AutomorphicArithmetic/ThetaCorrespondence/ThetaPairLifts/ThirdPlaceSeries/RatAtOfRows.lean}
\end{itemize}
\bibitem{S-gmin}
Submitted archive. Small sections and the theta-spanning input.
\begin{itemize}[leftmargin=1.4em,itemsep=0pt,topsep=1pt]
\item \path{Development/AutomorphicArithmetic/ThetaCorrespondence/ThetaPairLifts/HodgeBundleWeights/ExistsGoodCoverAtMinimalLevel.lean}
\item \path{Development/AbelianVarietyTheory/Proofs/DivisorProofs/NormPhiHatIndPhi.lean}
\item \path{Development/AbelianVarietyTheory/Proofs/DivisorProofs/ExistsBoundedIntegerGoodCoverThetaSeriesIntegral.lean}
\item \path{Development/AbelianVarietyTheory/Proofs/SiegelThetaTransformation/HermitianThetaSeriesSpanBoundedForms.lean}
\end{itemize}
\bibitem{S-theta-count}
Submitted archive. The lattice count, overlap, level factor and inversion.
\begin{itemize}[leftmargin=1.4em,itemsep=0pt,topsep=1pt]
\item \path{Development/AbelianVarietyTheory/Proofs/SiegelThetaTransformation/RelIndexDualSubmoduleInfOfSelfDual.lean}
\item \path{Development/AbelianVarietyTheory/Proofs/SiegelThetaTransformation/ThetaTestFunctionNormAndOverlapBound.lean}
\item \path{Development/AbelianVarietyTheory/Proofs/SiegelThetaTransformation/HermitianThetaLevelDataFormsSystem.lean}
\item \path{Development/AbelianVarietyTheory/Proofs/SiegelThetaTransformation/LatticeOperatorCoefficients.lean}
\item \path{Development/AbelianVarietyTheory/Proofs/SiegelThetaTransformation/IsOpenFrames.lean}
\item \path{Development/AbelianVarietyTheory/Definitions/SpanningFamilies/IsSystem.lean}
\item \path{Development/AbelianVarietyTheory/Proofs/SiegelThetaTransformation/ExistsClassicalLevelDataCoverSelfDual.lean}
\item \path{Formalism/LocalToGlobal/ProjectiveLattices/HalfOneAddSqrtNegIsIntegral.lean}
\end{itemize}
\bibitem{S-behrend}
Submitted archive. Cyclotomic orders, short intervals and random subsets.
\begin{itemize}[leftmargin=1.4em,itemsep=0pt,topsep=1pt]
\item \path{Formalism/BehrendSets/Statement.lean}
\item \path{Formalism/BehrendSets/BehrendVariance/Select.lean}
\end{itemize}
\bibitem{S-tlift}
Submitted archive. The exact theta-lifting assertion and its supplied proof.
\begin{itemize}[leftmargin=1.4em,itemsep=0pt,topsep=1pt]
\item \path{Development/AutomorphicArithmetic/Foundations/GaloisRepresentations/ThetaParseval/HalfIndexLattices/TLiftK.lean}
\item \path{Development/AutomorphicArithmetic/Foundations/GaloisRepresentations/ThetaParseval/ThetaKernels/TLiftK.lean}
\end{itemize}
\bibitem{S-tempered}
Submitted archive. Temperedness and the base-change hypotheses.
\begin{itemize}[leftmargin=1.4em,itemsep=0pt,topsep=1pt]
\item \path{Development/AbelianVarietyTheory/Proofs/SatakeParameters/LocalComponentTempered.lean}
\item \path{Development/AutomorphicArithmetic/Foundations/BaseChangeIdentities/Narrow.lean}
\item \path{Development/AutomorphicArithmetic/Foundations/BaseChangeIdentities/DefiniteNHolds.lean}
\end{itemize}
\bibitem{S-models}
Submitted archive. The traced uses of abstract local models.
\begin{itemize}[leftmargin=1.4em,itemsep=0pt,topsep=1pt]
\item \path{Formalism/SharpRestriction/Divisibility/LocalNonemptiness/SubsingletonRankIndicatorIff.lean}
\item \path{Development/AbelianVarietyTheory/Definitions/AutomorphicMembership/LocStOfLoc.lean}
\item \path{Development/AbelianVarietyTheory/Proofs/UnitaryExistence/LocalDataAtConcretePointsExist.lean}
\item \path{Formalism/LocalToGlobal/ProjectiveLattices/PhiSetPullSub.lean}
\item \path{Development/AutomorphicArithmetic/Foundations/UnitarySphericalFixing/SatakeParameters/SwOn.lean}
\end{itemize}
\bibitem{S-meta}
Submitted archive. The inspected term-assignment and declaration-copying code.
\begin{itemize}[leftmargin=1.4em,itemsep=0pt,topsep=1pt]
\item \path{Formalism/LocalToGlobal/IndefiniteReal/HeckeNonvanishing/Imports.lean}
\item \path{Formalism/LocalToGlobal/NumberFieldFamilies/OwnStruct.lean}
\item \path{Formalism/HermitianCurvature/HermitianCurvatureProofs/CutCptCopy.lean}
\end{itemize}

\bibitem{S-bernoulli-continued}
Archive sources under
\path{Development/AutomorphicArithmetic/AutomorphicVarieties/LevelCohomology/}:
\path{LemmaX.lean},
\path{UnramifiedCocycleIsCoboundaryOfBernoulliUnit.lean},
\path{LxCore.lean}, and \path{LxPartial.lean}.
The statement audited is \code{UpperTheory.V.X.LemmaX}, with its
\code{XData} hypotheses.

\bibitem{S-propz-continued}
Archive sources under
\path{Development/AutomorphicArithmetic/AutomorphicVarieties/}:
\path{HeightAndVanishing/Counts.lean};
\path{LevelCohomology/PropZ.lean},
\path{LevelCohomology/TorsClass.lean},
\path{LevelCohomology/ZqField.lean}, and
\path{LevelCohomology/HQuad.lean}.

\bibitem{S-even-selmer}
Archive sources under
\path{Development/AutomorphicArithmetic/AutomorphicVarieties/}:
\path{CharacterLevelBounds/ShaFamMono.lean},
\path{GaloisLibrary/SelmerTransport/ConjStableShaFam.lean}, and
\path{SymplecticGSp4/DecompSigmaOrd.lean}.

\bibitem{S-detector-continued}
Archive sources under
\path{Development/AutomorphicArithmetic/AutomorphicVarieties/PatchingModules/DeterminantPadicFree/}:
\path{FrameWith.lean}, \path{DetAlg.lean}, and \path{RelPair.lean}.

\bibitem{S-regfour}
Archive sources under
\path{Development/AutomorphicArithmetic/AutomorphicVarieties/}:
\path{AlgebraTools/MonomialRepresentations/RegCosetFour.lean},
\path{AlgebraTools/MonomialRepresentations/AdjugateCornerEqMinor.lean},
\path{AlgebraTools/MonomialRepresentations/DiscSwap.lean},
\path{AlgebraTools/MonomialRepresentations/RelWeight.lean},
\path{AlgebraTools/MonomialRepresentations/TorusTransfer.lean},
\path{AlgebraTools/MonomialRepresentations/IdCompMat.lean},
\path{AlgebraTools/MonomialRepresentations/MonomialRed.lean},
\path{AlgebraTools/MonomialRepresentations/PermChar.lean},
\path{WeakShimuraTheory/CompactIntegralFields/H4.lean}, and
\path{WeakShimuraTheory/StabilizerIndexBounds/NonComm.lean}.

\bibitem{milne-ant}
J. S. Milne, \emph{Algebraic Number Theory}, version 3.08,
19 July 2020, Theorem 4.3, printed p.~70 (Minkowski's bound).
\url{https://www.jmilne.org/math/CourseNotes/ANT.pdf}.

\bibitem{milne-cft}
J. S. Milne, \emph{Class Field Theory}, version 4.03,
6 August 2020, Theorem 0.4 and Chapter V, Example 3.9
(the ideal class group and the maximal everywhere unramified
abelian extension).
\url{https://www.jmilne.org/math/CourseNotes/CFT.pdf}.

\bibitem{milne-ag}
J. S. Milne, \emph{Algebraic Groups}, Cambridge University Press,
2022, Proposition 1.34 and Corollary 1.35
(the identity component and irreducibility).
\url{https://www.jmilne.org/math/Books/iAG2022.pdf}.


% Insert these items inside the host document's bibliography.
\bibitem{qrh-september}
OpenAI, \emph{The Quasi-Riemann Hypothesis: A Zero-Free Half-Plane
$\operatorname{Re}s>7/8$}, 30 September 2026, 199 pp.
Version at repository commit
\path{4eaee87548250b2682affda7d05d947c05554bae}.
Theorem 1.1; Proposition 2.1; Sections 17--20.
\url{https://github.com/openai/math/blob/4eaee87548250b2682affda7d05d947c05554bae/preprints/The-Quasi-Riemann-Hypothesis-September-30-2026/paper.pdf}.

\bibitem{qrh-october}
OpenAI, \emph{The Quasi-Riemann Hypothesis}, 5 October 2026,
48 pp. Version at the same repository commit.
Theorem 1.1, Proposition 3.1, Section 3, Propositions 5.1, 5.2,
5.4, and Appendix A.2.
\url{https://github.com/openai/math/blob/4eaee87548250b2682affda7d05d947c05554bae/preprints/The-Quasi-Riemann-Hypothesis-October-5-2026/paper2.pdf}.

\bibitem{qrh-catalogue}
H. Collet and contributors, \emph{QRH Catalogue}, repository
commit \path{5b7217781269248f39fbcec2432394d3d1610344},
consulted 20 November 2026. In particular the README's evidence
qualifications and the mathematical entries G276 and G280.
\url{https://github.com/HaileyStorm/QRH-Catalogue/tree/5b7217781269248f39fbcec2432394d3d1610344}.
This reference identifies a catalogue of source claims and
conditional reductions; it is not a certificate for every
listed theorem.

\bibitem{alpo-prime-supply}
L. Alp\"oge, submitted source archive, the ordinary progression
and quadratic-field prime conditions in
\path{Formalism/SharpRestriction/InertJacobiLevels/PrimeSupply.lean};
the varying quadratic ray-class conditions in
\path{Formalism/SharpRestriction/GlobalSchurBounds/ShortIntervalPrimes.lean}.
Archive SHA-256:
\path{9631b4f020199e580ba86178cb56c6e9b09f690c2369a80dc4334aa93f49c7fc}.

\bibitem{alpo-final-scalar}
L. Alp\"oge, same submitted archive,
\path{Formalism/SizeWindowBrackets/GrowthBoundGivesRiemannHypothesis.lean},
in particular the height hypotheses, the conditions on the
three auxiliary characters, and the final size inequalities.

\bibitem{sz-openai}
OpenAI, \emph{Uniform exclusion of Landau--Siegel zeros},
October 1, 2026, 9 pp.
\url{https://github.com/openai/math/blob/main/preprints/Uniform-exclusion-of-Landau-Siegel-zeros-October-1-2026/paper.pdf}.

\bibitem{sz-guo}
Victor Zhenyu Guo, \emph{An explicit elementary proof of no Landau--Siegel zeros},
arXiv:2610.30476v1, October 21, 2026, 17 pp.
\url{https://arxiv.org/html/2610.30476v1}.

\bibitem{sz-tao}
Terence Tao, \emph{A sketch of the proof of non-existence of Siegel zeroes},
What's New, October 19, 2026.
\url{https://terrytao.wordpress.com/2026/10/19/a-sketch-of-the-proof-of-non-existence-of-siegel-zeroes/}.

\bibitem{sz-pescador}
GPT-6 Pro and Claude Opus 5.5, posted and maintained by Jacob Pescador,
\emph{Landau--Siegel zeros and a Stepanov-type nonvanishing lemma},
version 3, October 22, 2026.
\url{https://jacobpescador.github.io/ai-math/papers/landau-siegel-stepanov/}.

\bibitem{sz-vandijk}
Roelof van Dijk,
\emph{A Frobenius determinant proof of the non-existence of Landau--Siegel zeros},
hexagon:2610.00409v2, 6 pp.
\url{https://hexagonmath.org/pdf/2610.00409v2}.

\bibitem{sz-kerr}
Bryce Kerr, \emph{A simpler proof that there are no Landau--Siegel zeros},
arXiv:2610.21099v1, October 15, 2026, 6 pp.
\url{https://arxiv.org/abs/2610.21099}.

\bibitem{ct-fw}
B. Ferrero and L. C. Washington,
The Iwasawa invariant $\mu_p$ vanishes for abelian number fields,
\emph{Annals of Mathematics} (2) \textbf{109} (1979), 377--395.
\url{https://doi.org/10.2307/1971116}.
Primary bibliographic record:
\url{https://annals.math.princeton.edu/1979/190-2/p05}.

\bibitem{ct-sharifi}
R. Sharifi, \emph{Iwasawa Theory}, author's lecture notes,
Chapter 3 (cyclotomic extensions and limits of class groups) and
Theorem 6.2.1 (Ferrero--Washington).
\url{https://www.math.ucla.edu/~sharifi/iwasawa.pdf}.
Author's reading edition, checked 21 November 2026:
\url{https://www.math.ucla.edu/~sharifi/notes/iwasawa-ch03.html},
\url{https://www.math.ucla.edu/~sharifi/notes/iwasawa-ch06.html}.

\bibitem{milne-adt}
J. S. Milne, \emph{Arithmetic Duality Theorems}, second edition,
BookSurge, 2006. In particular, Chapter II, Proposition 2.9
and Theorem 2.13 give the \'etale--Galois comparison and finite
coefficient cohomology used here.
\url{https://www.jmilne.org/math/Books/ADTnot.pdf}.

\bibitem{weibel-kbook}
C. A. Weibel, \emph{The $K$-book: An Introduction to Algebraic $K$-theory},
Graduate Studies in Mathematics \textbf{145}, American Mathematical
Society, 2013. Author's chapters III, IV and VI, especially III,
Example 1.3.6; IV,
Corollary 1.14.1, Theorems 1.17--1.18 and 6.9, and VI,
equation (8.1.1), Theorem 8.2 and Corollary 8.3.
\url{https://sites.math.rutgers.edu/~weibel/Kbook/Kbook.III.pdf};
\url{https://sites.math.rutgers.edu/~weibel/Kbook/Kbook.IV.pdf};
\url{https://sites.math.rutgers.edu/~weibel/Kbook/Kbook.VI.pdf}.

\bibitem{milnor-moore}
J. W. Milnor and J. C. Moore, On the structure of Hopf algebras,
\emph{Annals of Mathematics} (2) \textbf{81} (1965), 211--264,
especially Corollary 4.18, Theorem 5.18 and the appendix,
pp.\ 262--264.
\url{https://doi.org/10.2307/1970615}.
Author's article in the University of Rochester collection:
\url{https://www.sas.rochester.edu/mth/sites/doug-ravenel/otherpapers/Milnor-Moore.pdf}.

\bibitem{mathew-hurewicz}
A. Mathew, Torsion exponents in stable homotopy and the Hurewicz
homomorphism, \emph{Algebraic \& Geometric Topology}
\textbf{16} (2016), 1025--1041, especially Theorem 6.3.
\url{https://doi.org/10.2140/agt.2016.16.1025}.
Author's article:
\url{https://msp.org/agt/2016/16-2/agt-v16-n2-p12-s.pdf}.

\bibitem{alpo-uniform-source}
L. Alp\"oge, submitted source archive identified in the audit,
especially the directories
\path{Development/AutomorphicArithmetic/AutomorphicVarieties/LayeredIsomorphisms},
\path{Development/AutomorphicArithmetic/AutomorphicVarieties/AdditiveValuationNorms},
\path{Development/AutomorphicArithmetic/AutomorphicVarieties/QuadraticExtensions},
\path{Formalism/LocalToGlobal/Divisibility}, and the theta and
selected-character files listed in the accompanying source manifests.
The source traces specify the exact scope read; these locators do not
assert that every file in the directories was reviewed.

\bibitem{kirschmer-herm}
M. Kirschmer, \emph{Definite quadratic and hermitian forms with
small class number}, habilitation thesis, RWTH Aachen,
\S4.2, especially Theorem 4.2.3 and the definition of the local density,
pp.\ 50--51.
\url{https://www.math.uni-bielefeld.de/~mkirschm/papers/herm.pdf}.
The general mass formula there is attributed to
W. T. Gan and J.-K. Yu, \emph{Group schemes and local densities},
Duke Math.\ J.\ \textbf{105} (2000), 497--524, \S10.

\bibitem{terakado-yu-mass}
Y. Terakado and C.-F. Yu, \emph{Mass formulas and the basic locus
of unitary Shimura varieties}, arXiv:2210.04054v1 (2022),
\S2.3, especially equation (2.6).
\url{https://arxiv.org/pdf/2210.04054}.

\bibitem{harish-chandra-ds-ii}
Harish--Chandra, Discrete series for semisimple Lie groups. II:
Explicit determination of the characters,
\emph{Acta Mathematica} \textbf{116} (1966), 1--111,
especially Theorem 16 and \S41, pp.\ 96--100.
\url{https://doi.org/10.1007/BF02392813}.
Primary article in the Yau Mathematical Sciences Center collection:
\url{https://archive.ymsc.tsinghua.edu.cn/pacm_download/117/6013-11511_2006_Article_BF02392813.pdf}.

\bibitem{duflo-heckman-vergne}
M. Duflo, G. Heckman and M. Vergne,
Projection d'orbites, formule de Kirillov et formule de Blattner,
\emph{M\'emoires de la Soci\'et\'e Math\'ematique de France},
new series, no.\ 15 (1984), 65--128, especially III.1,
theorem (2), p.\ 111.
\url{https://www.numdam.org/article/MSMF_1984_2_15__65_0.pdf}.
This article gives another proof of the formula established by
H. Hecht and W. Schmid, A proof of Blattner's conjecture,
\emph{Inventiones Mathematicae} \textbf{31} (1975), 129--154,
\url{https://doi.org/10.1007/BF01404112}.

\bibitem{baldoni-vergne}
V. Baldoni and M. Vergne,
Discrete series representations and $K$ multiplicities for
$U(p,q)$. User's guide, arXiv:1008.5360v1 (2010),
especially \S2.2, pp.\ 12--14.
\url{https://webusers.imj-prg.fr/~michele.vergne/publications/discreteseries.pdf}.

\bibitem{parthasarathy-minimal}
R. Parthasarathy, Classification of discrete series by minimal $K$-type,
\emph{Representation Theory} \textbf{19} (2015), 167--185,
especially Theorem 1.1, pp.\ 167--168.
\url{https://doi.org/10.1090/ert/467}.
Author-posted full text, used to verify the theorem and its hypotheses:
\url{https://www.researchgate.net/publication/282775275_Classification_of_discrete_series_by_minimal_-type}.

\bibitem{bannai-kobayashi-2007}
K. Bannai and S. Kobayashi,
\emph{Algebraic theta functions and $p$-adic interpolation of
Eisenstein--Kronecker numbers},
arXiv:math/0610163v4, 11 December 2007;
published in Duke Mathematical Journal \textbf{153} (2010), 229--295.
The theorem numbers used here refer to the arXiv version:
Propositions 2.14--2.15, Corollary 2.16(i),
the reduced Mumford-translation formula in Lemma 2.7 and
Proposition 2.8, and \S3.1.
\url{https://arxiv.org/abs/math/0610163v4}.
\url{https://doi.org/10.1215/00127094-2010-024}.
\bibitem{stacks-07S7}
The Stacks Project Authors,
\emph{The Stacks Project}, Lemma 66.14.2, Tag 07S7,
free finite locally free quotients of projective schemes.
\url{https://stacks.math.columbia.edu/tag/07S7}.

\bibitem{sun-eisenstein}
Binyong Sun, \emph{Eisenstein integrals for theta stable parabolic
subalgebras}, arXiv:math/0403440v2, 28 November 2004.
Theorem 6.4 and \S6.4 give the compact-type integral;
\S\S1.1--1.4 fix the group, measure and conjugation conventions.
\url{https://arxiv.org/abs/math/0403440v2}.

\bibitem{vogan-zuckerman}
David A. Vogan, Jr. and Gregg J. Zuckerman,
\emph{Unitary representations with non-zero cohomology},
Compositio Mathematica \textbf{53} (1984), 51--90.
See (5.1), Theorem 5.3 and Proposition 6.1.
\url{https://www.numdam.org/item/CM_1984__53_1_51_0/}.

\bibitem{greither-popescu}
Cornelius Greither and Cristian D. Popescu,
\emph{An equivariant main conjecture in Iwasawa theory and applications},
arXiv:1103.3069v1, 16 March 2011.
Theorem 6.11 and Lemma 6.9, with the conventions in \S6.2.
\url{https://arxiv.org/abs/1103.3069v1}.


\bibitem{getz-hahn-2022}
J. R. Getz and H. Hahn,
\emph{An Introduction to Automorphic Representations},
author manuscript dated 22 April 2022, Theorems 6.2.3
and 6.3.2, with their fixed-level and compact-type conventions.
The theorem numbers here refer to this manuscript.
\url{https://sites.duke.edu/jgetz/files/2022/04/Graduate_Text.pdf}.

\bibitem{weil-siegel}
Andr\'e Weil,
\emph{La formule de Siegel et les groupes classiques},
Acta Mathematica \textbf{113} (1965), 1--87.
The unitary case is \S27, type $(I_2)$; see \S54 for
independence of the normalized quotient volume and
\S57 for the special linear Tamagawa number.
\url{https://archive.ymsc.tsinghua.edu.cn/pacm_download/117/5988-11511_2006_Article_BF02391774.pdf}.

\bibitem{louboutin-residue}
St\'ephane Louboutin,
\emph{Explicit upper bounds for residues of Dedekind zeta functions
and values of $L$-functions at $s=1$, and explicit lower bounds
for relative class numbers of CM-fields},
Canadian Journal of Mathematics \textbf{53} (2001), 1194--1222.
Theorem~1(1).
\url{https://doi.org/10.4153/CJM-2001-045-5}.

\bibitem{banaszczyk-transference}
W. Banaszczyk, New bounds in some transference theorems
in the geometry of numbers,
\emph{Mathematische Annalen} \textbf{296} (1993), 625--635,
especially Theorem~2.2.
\url{https://doi.org/10.1007/BF01445125}.

\bibitem{bost-kunnemann-extensions}
J.-B. Bost and K. K\"unnemann,
Hermitian vector bundles and extension groups on arithmetic
schemes. I. Geometry of numbers,
\emph{Advances in Mathematics} \textbf{223} (2010), 987--1106.
Author preprint, 11 January 2007, especially
Theorem~3.7.2, (3.54), printed p.~46.
\url{https://epub.uni-regensburg.de/13828/1/MP77.pdf}.

\bibitem{khare-larsen-savin}
C. Khare, M. Larsen and G. Savin,
Functoriality and the inverse Galois problem,
\emph{Compositio Mathematica} \textbf{144} (2008), 541--564.
Author version arXiv:math/0610860v3 (29 May 2015),
especially \S4.1, pp.~14--15.
\url{https://arxiv.org/abs/math/0610860}.

\bibitem{zaman-hecke-estimates}
A. Zaman, Explicit estimates for the zeros of Hecke
\(L\)-functions,
\emph{Journal of Number Theory} \textbf{162} (2016), 312--375.
Author version arXiv:1502.05679v4 (29 March 2016),
especially \S2.1 and Lemma~2.1.
\url{https://arxiv.org/abs/1502.05679}.

\bibitem{thorne-adjoint-2022}
J. A. Thorne, \emph{On the vanishing of adjoint Bloch--Kato
Selmer groups of irreducible automorphic Galois representations},
arXiv:2207.04925v1, 11 July 2022, Theorem~6.1.
\url{https://arxiv.org/pdf/2207.04925}.

\bibitem{newton-thorne-adjoint-2023}
J. Newton and J. A. Thorne, \emph{Adjoint Selmer groups of
automorphic Galois representations of unitary type},
J. Eur. Math. Soc. \textbf{25} (2023), 1919--1967.
doi:10.4171/JEMS/1228. Conventions on pp.~1920 and~1923.
\url{https://ems.press/content/serial-article-files/32858}.

\end{thebibliography}
\end{document}

